\pagebreak

\chapter{Inbuilt Procedures}

    \index{procedure!inbuilt}
    The following procedures are built into the compiler. User
    called modules or procedures must not use identifiers that conflict
    with these.
    
    The following abbreviations are used to denote modes -

    \begin{tabular}{| l | l |}
    \hline
    I  & immediate                      \\ \hline
    V  & value                          \\ \hline
    SV & selectvalue                    \\ \hline
    S  & static                         \\ \hline
    Q  & queue                          \\ \hline
    CM & combinatorial output memory    \\ \hline
    RM & registered output memory       \\ \hline
    PR & priority                       \\ \hline
    IN & input                          \\ \hline
    O  & output                         \\ \hline
    C  & clock                          \\ \hline
    \end{tabular}
    
    The following abbreviations are used to denote types -

    \begin{tabular}{| l | l |}
    \hline
    bits  & bits                                                    \\ \hline
    uint  & unsigned integer                                        \\ \hline
    int   & integer or unsigned integer                             \\ \hline
    fixed & fixed point                                             \\ \hline
    float & floating point                                          \\ \hline
    log   & logical (boolean)                                       \\ \hline
    str   & string                                                  \\ \hline
    enum  & enumerated type                                         \\ \hline
    type  & type                                                    \\ \hline
    file  & file                                                    \\ \hline
    map   & map                                                     \\ \hline
    list  & list                                                    \\ \hline
    any   & any type                                                \\ \hline
    prim  & a primitive type                                        \\ \hline
    comp  & a compound type (arrays, structs or any combination)    \\ \hline
    []    & array of type                                           \\ \hline
    \end{tabular}

    Procedures with 'yes' to 'Parameter=argument allowed' can have arguments
    matched to parameters by position in the argument list or by
    parameter=argument (keymatch) or a mixture as long as positional arguments
    do not follow keymatch arguments. The keymatch arguments may be in
    any order. Otherwise these keymatch arguments are
    not allowed. This may be because - \\
    \blist
    \item   the procedure has a single input or output argument in which case
    no positional or keymatch parameter association is needed
    \item   the procedure has an indefinite number of arguments and hence there
    are no 'named' parameters with which to associate arguments
    \blend
    The keymatching may be restricted to just input or just output.

    Procedures with 'yes' to in-line target execution will produce in-line
    code in the FPGA. This code is controlled by the execution stream in which
    it is embedded and by the availability of any queues in the arguments.
    Other procedures do not produce in-line code in the FPGA. However, they may 
    produce target components which are accessed by in-line target code
    elsewhere in the source code.
        


    \section{addmodeattribute - add attribute to the mode attribute map}\label{sec:ipaddmodeatt}
        
        {\bf Call:} \\
        addmodeattribute(mode, key, isarray, type);
        
        {\bf Output parameters:} \\
        none
        
        {\bf Input parameters:} \\
        \begin{tabular}{| l | l | l | l |}
        \hline
        param  & mode & type & \\
        \hline \hline
        mode    & I & str & mode           \\ \hline
        key     & I & str & attribute      \\ \hline
        isarray & I & log & is an array    \\ \hline
        type    & I & str & primitive type \\ \hline
        \end{tabular}
        
        {\bf Parameter=argument allowed:} \\
        Yes
        
        {\bf In-line target execution:} \\
        No
        
        {\bf Description:} \\
        Add an attribute to the attribute map for a variable mode. 3PL contains
        an internal map initialised with entries for allowed attributes for each
        mode. These attributes are internal to 3PL and affect the netlist
        generation. Attributes from which constraints are generated are not
        handled internally by 3PL but by library code appropriate to the FPGA
        platform. To allow centralised vetting of attributes in calls to inbuilt
        procedures input(), output() and clock() this procedure is used to add
        these constraint attributes to the map.
        
        {\bf mode} is a string identifying the mode and is "input", "output"
        or "clock"
        
        {\bf key} is the attribute name.
        
        {\bf isarray} is true if the attrbute value is an array (one member
        for each bit). If not an array the value will be duplicated for
        each bit.
        
        {\bf type} is the attribute value primitive type and can be "uint",
        "float", "log" or "str" as appropriate.



        
    \section{append - append a list entry}\label{sec:ipappend}
        
        {\bf Call:} \\
        append(list, index, value);
        
        {\bf Output parameters:} \\
        none
        
        {\bf Input parameters:} \\
        \begin{tabular}{| l | l | l | l |}
        \hline
        param  & mode & type & \\
        \hline \hline
        list   & I     & list & a list   \\ \hline
        index  & I     & int  & an index   \\ \hline
        value  & I V C & any  & the value   \\ \hline
        \end{tabular}
        
        {\bf Parameter=argument allowed:} \\
        Yes
        
        {\bf In-line target execution:} \\
        No
        
        {\bf Description:} \\
        The item {\bf value} is appended at the position indicated by {\bf
        index} in list {\bf list}. The index may range from 0 to {\bf
        entries(list)} inclusive. If {\bf index} is outside this range a fatal
        error occurs. An index of 0 inserts at the beginning of the list and
        an index equal to the list size appends to the end of the list. If
        the index argument is omitted the item will be appended to the end
        of the list. If the last argument is omitted null will be appended.
        
        The appended item may be immediate, clock, cmemory or rmemory mode.


    \section{appendall - append entries from another list}\label{sec:ipappendall}
        
        {\bf Call:} \\
        appendall(l, l2);
        
        {\bf Output parameters:} \\
        none
        
        {\bf Input parameters:} \\
        \begin{tabular}{| l | l | l | l |}
        \hline
        param  & mode & type & \\
        \hline \hline
        l      & I     & list & a list   \\ \hline
        l2     & I     & list & another list   \\ \hline
        \end{tabular}
        
        {\bf Parameter=argument allowed:} \\
        No
        
        {\bf In-line target execution:} \\
        No
        
        {\bf Description:} \\
        The contents of list {\bf l2} are appended at the position indicated
        by {\bf index} in list {\bf l}. The index may range from 0 to {\bf
        entries(l)} inclusive. If {\bf index} is outside this range a fatal
        error occurs. An index of 0 inserts at the beginning of the list and
        an index equal to the list size appends to the end of the list. If
        the index argument is omitted the items will be appended to the end
        of the list.
        


    \section{argadjust - adjust arguments for arithmetic operations}\label{sec:ipargadjust}
        
        {\bf Call:} \\
        argadjust(out1, out2, rt \verb'<-' in1, in2, op, matchleft);
        
        {\bf Output parameters:} \\
        \begin{tabular}{| l | l | l | l |}
        \hline
        param  & mode     & type & \\
        \hline \hline
        out1   & V        & -    & adjusted first operand   \\ \hline
        out2   & V        & -    & adjusted second operand  \\ \hline
        rt     & I        & type & result type              \\ \hline
        \end{tabular}
        
        {\bf Input parameters:} \\
        \begin{tabular}{| l | l | l | l |}
        \hline
        param      & mode & type & \\
        \hline \hline
        in1        & V    & -    & first operand   \\ \hline
        in2        & V    & -    & second operand  \\ \hline
        op         & I    & str  & operation       \\ \hline
        matchleft  & I    & log  & use left offset \\ \hline
        \end{tabular}
        
        {\bf Parameter=argument allowed:} \\
        Yes
        
        {\bf In-line target execution:} \\
        No
        
        {\bf Description:} \\
        This procedure takes two value mode arguments and an operator and
        supplies two adjusted values and a result type. The input operands
        (the 1st and 2nd arguments) must be type "uint:", "int:", "ufixed:"
        or "fixed:". The 3rd input argument, the operator type, must be
        one of the strings "+", "*" or "/".
        
        This procedure accesses code within the 3PL interpreter used for
        adjusting target variables in arithmetic expressions.
        
        The "+" operator is used for addition and "-" for subtraction but the
        action is identical. The operand with the smaller offset will be
        adjusted to match the other, i.e. the binary points will be aligned
        by left shifting the operand with fewer fractional bits. The type
        assigned to the 3rd output argument will have the same binary point
        offset as the now-adjusted operands and the integer part will be
        one bit wider than the operand with the largest integer part.
        
        If the 4th input argument is supplied and is true, the result offset
        will be the same as the first input operand and the second input operand
        will be adjusted by left or right shifting as required. This feature
        is included for the case where the first operand is the result of
        an intermediate hardware operation and is not accessible for adjustment
        (for example in an XC5V DSP block where a multiplication can be followed
        by an addition or subtraction but the intermediate product cannot be
        adjusted).
        
        The "*" operator is used for multiplication. The operands are not
        adjusted and are returned as the output values. The output result type
        fraction width is the sum of the operand fraction widths and the
        integer width is the sum of the operand integer widths.
        
        The "/" operator is used for division. The operands are not
        adjusted and are returned as the output values. The output result type
        fraction width is the sum of the numerator fraction width and the
        denominator integer width. The result integer width is the sum of the
        numerator integer width and the denominator fraction width.



    \section{assert - assert a value mode variable for a clock cycle}\label{sec:ipassert}
        
        
        {\bf Call:} \\
        assert(v \verb'<-' o);
        
        {\bf Output parameters:} \\
        \begin{tabular}{| l | l | l | l |}
        \hline
        param & mode & type & \\
        \hline \hline
        v & V & log & value carrying the pulse \\ \hline
        \end{tabular}
        
        {\bf Input parameters:} \\
        \begin{tabular}{| l | l | l | l |}
        \hline
        param & mode & type & \\
        \hline \hline
        or & I & log & if true, OR with previous value \\ \hline
        \end{tabular}
        
        {\bf Parameter=argument allowed:} \\
        No
        
        {\bf In-line target execution:} \\
        yes
        
        {\bf Description:} \\
        This procedure asserts a logical value variable true for one clock cycle
        when this statement executes. The value is true for the clock cycle
        leading up to the clock edge on which execution occurs.

        If the input argument is provided and is true the output value variable
        will be assigned the OR of its previous value and the execution pulse.
        If it has no previous assigned value it is assigned the execution pulse.
        Using this optional input argument a signal pulsed from more than one
        place can be generated. {\bf Note: when ORed in this way, use of the
        value of the output variable must follow the last call to assert() since
        each call re-assigns the value mode variable!}

        If the ORing option is not selected and multiple calls to
        {\bf assert()} are made with the same output argument, each call will
        re-assign the output value variable, which may be what is wanted or
        may not!
        
        The following code illustrates some points to consider -
        \begin{lstlisting}[gobble=8, language=threepl]
           value({l, v1, v2}, "log");
           static(v, "log");

           when (l) seq {
               ....
               .
               .
               par {
                   assert(v1 <-);
                   v2 <- true;
               }
               v2 <- false;
               .
           }
        \end{lstlisting}
        
        'v1' and 'v2' will both be be true for one clock cycle, but 'v1' will
        be asserted one clock cycle earlier than 'v2'. Use of {\bf assert()}
        has the potential disadvantage that since it is actually the execution
        signal it suffers any delays inherent in the logic path of any
        conditional and hence might introduce timing constraint difficulties.
        
        Note that the following usage is allowed -
        \begin{lstlisting}[gobble=8, language=threepl]
           value(v, "log");

           when (v) 
               ....
           .
           .
           assert(v <-);
           .
        \end{lstlisting}
        Here the call to {\bf assert()} occurs only once. 3PL notices when
        assigning a value to {\bf v} within the call to {\bf assert()} that
        {\bf v} has been referenced previously at which time it had not
        been assigned a value, so it takes care of the post-assignment
        (see discussion of post-assignment in section \ref{sec:im_valass}).
        The following {\bf would not work in the same way} -
        \begin{lstlisting}[gobble=8, language=threepl]
           value(v, "log");

           when (v)
               ....
           .
           .
           assert(v <-);
           .
           assert(v <- true);   // PROBABLY WRONG!
           .
        \end{lstlisting}
        The second call to {\bf assert()} would indeed OR the two execution
        signals together and assign the result to variable {\bf v}, but the
        {\bf when ()} would use the value of {\bf v} after the first call
        to {\bf assert()}. Since with parallel code the order of statements
        is not significant, it would clearly be better to use a policy of
        referencing the variable following all calls to {\bf assert()}.

        see also inbuilt function {\bf used()} \autorefph{sec:ifused}.


    \section{attributeprocedure - register a procedure to handle attributes}\label{sec:ipattributeproc}
        
        {\bf Call:} \\
        attributeprocedure(p);
        
        {\bf Output parameters:} \\
        none
        
        {\bf Input parameters:} \\
        \begin{tabular}{| l | l | l | l |}
        \hline
        param & mode & type  & \\
        \hline \hline
        p     & - &  -  & procedure identifier \\ \hline
        \end{tabular}
        
        {\bf In-line target execution:} \\
        No
        
        {\bf Description:} \\
        This procedure registers the argument procedure as a user-defined procedure
        to handle input, output and clock mode attributes. The registered procedure
        will be called after every input, output and clock mode variable is
        created. The called procedure will be passed a single argument which is a
        pointer to the newly-created variable.

        This facility is intended to allow attributes to be attached to I/O pin
        specifications and have them transferred to relevant input, output or
        clock mode variables when those variables are created. It can also provide
        any necessary re-formatting of attributes, either those brought in from pin
        associations or those passed as arguments to inbuilt procedures input(),
        output() or clock().
        
        There is a Xilinx library procedure for this purpose;
        {\bf processPinAttributes()} \autorefph{sec:lpprocPinAtt},
        in xilinx/common.3pl.
        
        attributeprocedure() should be called prior to calls to inbuilt procedures
        input(), output() or clock().
        



    \section{attributes - assign attributes to a variable}\label{sec:ipattributes}
        
        {\bf Call:} \\
        attributes(v, attr);
        
        {\bf Output parameters:} \\
        none
        
        {\bf Input parameters:} \\
        \begin{tabular}{| l | l | l | l |}
        \hline
        param & mode & type  & \\
        \hline \hline
        v     & any &  -  & variable identifier \\ \hline
        attr  & I   & map & attributes map      \\ \hline
        \end{tabular}
        
        {\bf Parameter=argument allowed:} \\
        No
        
        {\bf In-line target execution:} \\
        No
        
        {\bf Description:} \\
        This procedure applies attributes to one or more variables. The first
        argument may be a compound value containing a number of identifiers for
        created variables. Identifiers may have fields or subscripts appended
        but these will be ignored, the attribute(s) being applied to the
        entire variable.
        
        The optional second argument can be
        \blist
        \item   a map of attributes
        \item   a one-dimensional compound constant containing key=value entries
        \item   a key=value argument
        \blend
        
        In the last case this may be followed by more key=value arguments. If this
        is not supplied the attributes can be added later using inbuilt procedure
        {\bf attributes()} \autorefph{sec:ipattributes}. For a description of
        applicable attributes see \autorefph{sec:attributes}.
        
        Attributes are applied to a variable in a map. The map
        keys are the attribute identifiers and the associated values
        are the attribute values. The entry values are mixed types,
        the type depending on the individual attribute.

        A compound value can be used for the attributes argument. It
        must be a one-dimensional list and every entry must be
        key=value format.

        The attributes() procedure does not enter a new attributes map in
        the variable, it adds the entries in the argument map to that in the
        variable. Thus the attributes() procedure can be called multiple times
        to add one or more attributes to the variable. If the same attribute is
        entered again the new value replaces the old value. This can be used to
        enable and disable 3PL attributes that are used during immediate
        execution.

        Netlist attributes can be deleted by assigning the value {\bf null}.
        3PL attributes cannot be deleted; 3PL attributes either
        relate to parameters that are mandatory such as clock domains, or are
        boolean in which case they can be disabled rather than deleted, so
        there is no reason to delete them.
        
        3PL attributes {\bf domain}, {\bf readdomain} and {\bf writedomain} may be
        assigned null which is interpreted to mean that the current clock domain is
        to be used.
        
        {\bf Note that for input, output and clock mode variables many of the
        attributes are ignored by 3PL itself and must be processed by 3PL source
        code. These modes are handled by inbuilt procedures input(), output() and
        clock() which call a library procedure named by
        attributeprocedure()\autorefph{sec:ipattributeproc} to process the
        attributes. These attributes will also be added to the variable but 3PL
        itself will ignore them.}
        
        Attributes are described in the sections pertaining to the relevant
        variable creation procedures.



    \section{change - change a list entry}\label{sec:ipchange}
        
        {\bf Call:} \\
        change(list, index, value);
        
        {\bf Output parameters:} \\
        none
        
        {\bf Input parameters:} \\
        \begin{tabular}{| l | l | l | l |}
        \hline
        param  & mode & type & \\
        \hline \hline
        list   & I     & list & a list   \\ \hline
        index  & I     & int  & an index   \\ \hline
        value  & I V C & any  & the value   \\ \hline
        \end{tabular}
        
        {\bf Parameter=argument allowed:} \\
        Yes
        
        {\bf In-line target execution:} \\
        No
        
        {\bf Description:} \\
        The item in list {\bf list} at position {\bf index} is replaced by item
        {\bf value}. The index may range from 0 to {\bf entries(l) - 1}
        inclusive. If {\bf index} is outside this range a fatal error
        occurs. The new value does not have to be the same type as the
        previous value. If the last argument is omitted the list item will
        be changed to null.
        
        The new item may be immediate, clock, cmemory or rmemory mode.


    \section{class - create a variable of type "class"}\label{sec:ipclass}

        {\bf Call:} \\
        class(c, init);

        {\bf Output parameters:} \\
        None.

        {\bf Input parameters:} \\            
        {\bf init} is an immediate value of type "class".
        
        {\bf Parameter=argument allowed:} \\
        No

        {\bf In-line target execution:} \\
        no

        {\bf description:} \\
        This procedure creates one or more variables of type "class".
        
        The first argument is the variable identifier, or it can be a
        comma-separated list of identifiers enclosed in braces.
        
        The second argument is optional and is an initial value for the variable.
        It is only allowed where a single variable is being created. The "class"
        value may be from another variable of the same type but more usually it
        would be the returned value from a call to a class, for example the
        following creates a variable {\bf cl} of type "class" and assigns to it
        an instance of the class {\bf c} returned from a call to {\bf c()}.

        \begin{lstlisting}[gobble=8, language=threepl]
            class(cl, c());
            .
            .
            
            class c () {
                .
                .
                .
            }
        \end{lstlisting}



    \section{clock - create one or more clock variables}\label{sec:ipclock}
        
        {\bf Call:} \\
        clock(ident, ....);
        
        {\bf Output parameters:} \\
        none
        
        {\bf Input parameters:} \\
        \begin{tabular}{| l | l | l | l |}
        \hline
        param  & mode & type  & \\
        \hline \hline
        ident or \verb'{'id1, id2... \verb'}' & - &  -  & variable identifier(s)   \\ \hline
        .... & - & - & optional attributes \\ \hline
        \end{tabular}
        
        {\bf Parameter=argument allowed:} \\
        No
        
        {\bf In-line target execution:} \\
        No
        
        {\bf Description:} \\
        Create a clock variable.

        Create one or more clock mode variables.
        
        {\bf ident} is the variable identifier. If more than one
        variable is to be created the list of variables must be in braces
        {\em \{id1, id2, \dots\}}. The variable identifier can also be constructed
        from a string by using the indirect operator on a string variable or
        parenthesised expression, see subsection \autorefph{subsec:strvarcreat}.
        
        When a {\bf loc} attribute is provided only a single clock variable
        can be created by this procedure.
        
        When {\bf clock()} is used to create a module, procedure or function
        parameter it cannot be passed attributes. The created parameter variable
        will adopt the attributes of the associated argument.
        
        Any following arguments specify attributes. Each argument may specify an
        attribute by keymatch format, i.e. attribute=value where 'attribute' is
        an attribute identifier and 'value' is an appropriate mode and type.
        Alternatively an argument can be a map containing attributes where the
        key (string) is the attribute identifier and the associated value is
        the attribute value.

        Attributes may also be associated with the input pin or pin pair. This is
        described in chapter \autorefph{chap:cpa} {\bf Constraints, Properties and
        Attributes} where other aspects of processing attributes and constraints
        for input, output and clock mode variables are also described.

        Attributes can be added by using inbuilt procedure {\bf attributes()}
        \autorefph{sec:ipattributes}.             
        
        Attributes are required for two main purposes; control of netlist
        generation and constraint file entries. Attributes giving input pin
        locations and input buffer types directly control netlist generation.
        I/O standards and timing group names exercise control of place-and-route
        via a constraint file.
              
        The applicable attributes are shown in the following table.

                These attributes are not case-sensitive but will be converted to
        lower case.
         
        The attribute values will be applied to all pins.
        
        \btabularx
        \hline
        \multicolumn{4}{|l|}{\bf clock mode} \\
        \hline
        frequency    & 3PL     & float         & frequency, Hz \\
        dutycycle    & 3PL     & float         & duty cycle    \\
        differential & netlist & log           & differential  \\
        ibufg        & netlist & log           & buffer type   \\
        loc          & netlist & str or [2]str & locations     \\
        \hline
        \btabularxend
        
        Attributes {\bf frequency} and {\bf dutycycle} specify these parameters
        for clock nets and any associated clock domains. {\bf Note that at
        present dutycycle is not used. It has been included for future use.}

        The following attribute is read-only and is created when inbuilt
        procedure {\bf clock()} is called and can be retrieved from the clock
        mode variable if required by the function call attribute("ids").
        
        \btabularx
        \hline
        \multicolumn{3}{|l|}{\bf clock mode} \\
        \hline
        ids          & 3PL & []str & element IDs       \\
        \hline
        \btabularxend
        
        The following attributes only apply to a clock variable if it is a clock
        input from a pad. If the {\bf loc} attribute is supplied to a clock
        variable an input pad and buffer will be created, otherwise the clock
        variable source is internal to the FPGA.

        {\bf loc} specifies one or two pad locations (FPGA pins) for the clock
        input - two locations if it is differential (+ve, -ve order). For a
        differential clock input this attribute will be an array of two strings.

        {\bf differential} indicates that the clock input is differential - two
        locations must be supplied and the buffer will be an IBUFDS or IBUFGDS
        rather than the default IBUF or IBUFG. This attribute is used internally
        by 3PL for code generation.

        {\bf ibufg} if true specifies that the input buffer must be an IBUFG or
        IBUFGDS instead of the default IBUF or an IBUFDS. This attribute is used
        internally by 3PL for code generation. It is not clear from Xilinx
        documentation that an IBUFG needs to be explicitly specified, or if
        instead this can be inferred from the package pin(s). other software. 

        The following attributes are required for constraint generation and are
        handled at completion of immediate execution by library procedure {\bf
        generateconstraints()} \autorefph{sec:lpgenerateconstraints}. The
        attribute values are single values which will be applied to the pin or
        pins.
        
        \btabularx
        \hline
        \multicolumn{4}{|l|}{\bf clock mode constraint attributes} \\
        \hline
        tnm          & constraint & str & timing group id                       \\
        diff\_term   & constraint & log & differential  $100\Omega$ termination \\
        iostandard   & constraint & str & voltage standard                      \\
        \hline
        \btabularxend

        {\bf tnm} is the name of a timing group in which all state devices
        clocked by this clock are to be included for timing constraint purposes.
        All clocks which drive logic clock inputs (as distinct from inputs to
        buffers of clock managers) will have a period constraint generated by
        generate\_constraints() in xilinx/common.3pl. By default a timing group
        name will be constructed from the clock identifier translated to upper
        case, but if a "tnm" attribute is provided then this string will be used
        instead.        

        {\bf diff\_term} indicates that a $100\Omega$ termination resistance should
        be applied to a differential input.

        {\bf iostandard} is a string specifying one of the allowed I/O standards
        (see Xilinx documentation).

        The following attribute is read-only and is added when target statement
        variables are assigned or evaluated and the clock domain is established
        or verified. It can be retrieved from the clock mode variable if
        required by the function call attribute("valid").
        
        \btabularx
        \hline
        \multicolumn{3}{|l|}{\bf clock mode} \\
        \hline
        valid   & 3PL & log & is a valid logic clock    \\
        \hline
        \btabularxend

        Any attributes that are not listed above will result in a fatal error.

                
        Where the call to {\bf clock()} is to create a parameter to a module,
        procedure or function, only the first argument may be given and it must be
        a single identifier. The clock frequency, period and duty cycle are
        copied from the matching argument.
        
        Where a clock variable is an external input the {\bf loc} attribute
        must be supplied. This would be an array of two location strings for
        a differential clock input or a singlr location string otherwise.

        As far as the compiler is concerned any clock variable may be used to
        control logic (by calling inbuilt procedure {\bf useclock()}). However,
        many clock variables are used as routing to connect clock managers,
        input pads, buffers etc together and these interconnect clock variables
        must not be passed to {\bf useclock()}.
       
        Note: for each pin a netlist port "PORT\_pin" is created where 'pin' is
        the pin identifying string. This port name is also the netlist identifier
        of the netlist from the output of the pad to the input of the IBUFG.
        This convention can be used when generating associated constraints.



    \section{clockfromexpr - drive a clock from an expression}\label{sec:ipclockfromexpr}

        {\bf Call:} \\
        clockfromexpr(clk \verb'<-' expr);

        {\bf Output parameters:} \\
        {\bf clk} is a clock variable.

        {\bf Input parameters:} \\            
        {\bf expr} is an expression of type log which is to drive the clock.
        
        {\bf Parameter=argument allowed:} \\
        No

        {\bf In-line target execution:} \\
        no

        {\bf description:} \\
        This procedure connects a clock signal to an expression. That
        clock signal would normally then be buffered using a global clock
        buffer created by library procedure {\bf clockbufg}.
        The argument may be input, selectvalue, value, static or priority mode
        but must be type log and must contain no queue references.



    \section{close - close a file}\label{sec:ipclose}
        
        {\bf Call:} \\
        close(file);
        
        {\bf Output parameters:} \\
        none
        
        {\bf Input parameters:} \\
        \begin{tabular}{| l | l | l | l |}
        \hline
        param  & mode & type  & \\
        \hline \hline
        file & I & file & file variable   \\ \hline
        \end{tabular}
        
        {\bf Parameter=argument allowed:} \\
        No
        
        {\bf In-line target execution:} \\
        No
        
        {\bf Description:} \\
        This procedure closes an open file.



%    \section{connect - connect two signals}\label{sec:ipconnect}
%        
%        {\bf Call:} \\
%        connect(vsink \verb'<-' vsrc);
%        
%        {\bf Output parameters:} \\
%        \begin{tabular}{| l | l | l | l |}
%        \hline
%        param  & mode & type  & \\
%        \hline \hline
%        vsink & V & log & value variable   \\ \hline
%        \end{tabular}
%        
%        {\bf Input parameters:} \\
%        \begin{tabular}{| l | l | l | l |}
%        \hline
%        param  & mode & type  & \\
%        \hline \hline
%        vsrc & V & log & value variable   \\ \hline
%        \end{tabular}
%        
%        {\bf Parameter=argument allowed:} \\
%        No
%        
%        {\bf In-line target execution:} \\
%        No
%        
%        {\bf Description:} \\
%        This procedure connects two signals. The two arguments must be value
%        mode.
%        
%        {\bf This is intended solely for use where basic library elements are being
%        used so that, for example, location constraints can be applied.
%        \index{location constraint}
%        \index{constraint!location}
%        It is not intended for use in normal high-level 3PL code.} This procedure
%        is needed where the outputs of basic library elements are to be connected
%        on 3-state nets, e.g. BUFT, BUFE, PULLUP, PULLDOWN, KEEPER.


    \section{cmemory - create a combinatorial output memory}\label{sec:ipcmemory}
        \index{variable!memory!combinatorial output}
        
        {\bf Call:} \\
        cmemory(id, ports, atype, dtype, init, ...);
        
        {\bf Output parameters:} \\
        none
        
        {\bf Input parameters:} \\
        \begin{tabular}{| l | l | l | l |}
        \hline
        param  & mode       & type  & \\
        \hline \hline
        id     & -  &  -    & memory identifier         \\ \hline
        ports  & I  & int   & number of ports           \\ \hline
        atype  & I  & int   & address type              \\ \hline
        dtype  & I  & int   & data type                 \\ \hline
        init   & I  & []int & optional initialisation (variable or compound
                              constant)                 \\ \hline
        ...    & I  &  -    & optional attributes       \\ \hline
        \end{tabular}
        
        {\bf Parameter=argument allowed:} \\
        Yes
        
        {\bf In-line target execution:} \\
        No
        
        {\bf Description:} \\
        Procedure {\bf cmemory()} creates a combinatorial output RAM.        
        {\bf cmemory()} may create a module, procedure or function
        parameter in which case only the first argument is needed (the memory
        identifier).

        In the case of Xilinx FPGAs combinatorial output memory is
        distributed CLB RAM. This can have one or two access ports. One
        logic slice can hold 1 bit x 32 words single port or 1 bit x 16
        words dual port.

        Port 0 can be read or written (or both). Port 1 is read-only.
        
        {\em id} is an identifier for the memory. The variable identifier can
        also be constructed from a string by using the indirect operator on a
        string variable or parenthesised expression, see subsection
        \autorefph{subsec:strvarcreat}.

        {\em ports} is the number of ports. This argument may be by position or
        by ports=x where x is a constant or immediate variable. For Xilinx FPGAs
        theis must be 1 or 2.

        {\em atype} is the address type. This can be any target type but
        would normally be type "uint:". The type has the constraint that the total number of
        bits must not exceed 9. This argument may be by position or by
        atype=x where x is a string specifying a type or else a type variable.       

        {\em dtype} is the data type. This can be any target type. This argument
        may be by position or by dtype=x where x is a string specifying a type
        or else a type variable.
        
        {\em init} is an optional initialisation array. The initialisation data
        is a one-dimensional variable or array which is no larger than the depth
        of the memory. It may be smaller in which case the memory is zero
        padded. The types and structure of the initialisation variable or
        structured constant after the first dimension must match the memory data
        type. Immediate constants are padded to fit target elements but must not
        exceed the target width. If the initialisation argument is omitted the
        memory will be initialised to all zero. This argument may be by position
        or by init=x where x is an immediate array variable. For compound types
        the immediate type must match the target data type but omitting the data
        widths. The target type can be converted to a matching immediate type
        by the builtin function {\em targtoimtype()}.
        
        Where the {\em init} argument is supplied in the call to cmemory() the
        argument may be compound constant, though this is not usually practical as
        {\em init} is usually a large array. Otherwise the argument will be
        an immediate variable of appropriate type. The initialising variable must be a
        matching type to the target data type of the memory. The initial value
        variable can be created from the target data type by calling builtin function
        {\bf targtoimtype()} \autorefph{sec:iftargtoimtype} -

        \begin{lstlisting}[gobble=8, language=threepl]
            struct (st,
                f1, "uint:30",
                f2, "uint:30"
            );
            var(init, "[8]" + targtoimtype(st));
            init[0].f1 = ....;
            init[0].f2 = ....;
            .
            .
            .
            var(init0, targtoimtype(st));
            init0.f1 = ..;
            init0.f2 = ..;
            cmemory(m, 1, atype="uint:3", dtype=st, init=init);
        \end{lstlisting}
            
        {\em init} can be assigned to a memory
        variable, along with attributes, later in execution by calling
        builtin procedure {\bf attributes()}\autorefph{sec:ipattributes}, e.g.
        \begin{lstlisting}[gobble=8, language=threepl]
            cmemory(m, 1, atype="uint:3", dtype=st);
            attributes(m, init=init);
        \end{lstlisting}
        but in this case the argument must be an immediate variable, not a compound
        constant. The call to attributes() can be any time after calling
        cmemory() and (obviously) before termination of execution.

        There may be trailing arguments following the optional {\em init} which supply
        memory attrbutes, described in a table below. {\em init} is included amongst
        the attributes for convenience. These arguments must all be of the form
        attribute=value. These may be included in the call to cmemory() or else passed
        later using builtin procedure attributes().

                \btabularx
        \hline
        \multicolumn{4}{|l|}{\bf cmemory mode attributes} \\
        \hline
        init       & 3PL       &                      & initial value                    \\
        domain     & 3PL       & clock                & force clock domain for all ports \\     
        \hline
        \btabularxend
        
        {\bf These attributes are case-insensitive.}
        
        {\bf init} sets an initial value for the memory. Any previous initial
        value will be overwritten. The initial value must be a compound type and
        the nested type must be an immediate type equivalent to the memory data
        type, i.e. same type but without a width field. The initial value is not
        used until completion of immediate execution when variables are
        processed into target code and hence it can be changed at any time
        during immediate execution.

        {\bf domain} forces the clock domain (both input and output) for all
        ports to the specified clock domain. If assigned {\bf null} the current
        clock domain is selected. If the variable already has a clock domain as
        a result of having been assigned or evaluated then that clock domain is
        compared with the intended domain. If the assigned clock domain agrees
        with that already in place then there is no further action, otherwise a
        fatal error occurs.


        Architectures supported -

        \begin{tabular}{| l | l |} \hline
        Spartan-II, Spartan IIE         & YES   \\ \hline
        Spartan-3                       & YES   \\ \hline
        Virtex, Virtex-E                & YES   \\ \hline
        Virtex-II, Virtex-II Pro        & YES   \\ \hline
        XC9500, XC9500XV, XC9500XL      & NO    \\ \hline
        Coolrunner XPLA3, Coolrunner-II & NO    \\ \hline
        \end{tabular}

        The range of memory depths that have been implemented is shown in the
        following table. The depth required is determined from the number of bits
        in the address type. Memories with depths smaller than the minimum listed
        are still implemented, but the physical memory will be the minimum size
        shown and the smaller depth will be achieved by tying higher address bits
        low. The column marked ROM is for memories that are initialised and are
        read but not written -

        \begin{tabular}{| l | l | l | l |} \hline
                                & ROM       & 1 port   &   2 port   \\ \hline \hline
        Spartan-II, Spartan IIE & 16 -  256 & 16 -  256 & 16 -  128 \\ \hline
        Spartan-3               & 16 - 2048 & 16 -  512 & 16 -  128 \\ \hline
        Virtex                  & 16 -  256 & 16 -  256 & 16 -  256 \\ \hline
        Virtex-E                & 16 -  256 & 16 -  256 & 16 -  256 \\ \hline
        Virtex-II               & 16 - 2048 & 16 - 1024 & 16 -  512 \\ \hline
        Virtex-II Pro           & 16 - 2048 & 16 - 1024 & 16 -  512 \\ \hline
        Virtex 4                & 16 - 2048 & 16 -  512 & 16 -  128 \\ \hline
        Virtex 5                & 16 - 4096 & 16 - 4096 & 16 - 2048 \\ \hline                                             
        \end{tabular}

        The data may be any width.
        
        Applicable attributes are shown in the following table. The
        attributes can be passed as arguments to this procedure or may be
        applied to the variable using the inbuilt procedure
        {\bf attributes()} \autorefph{sec:ipattributes}.
        
        When {\bf cmemory()} is used to create a module, procedure or function
        parameter it cannot be passed attributes. The created parameter variable
        will adopt the attributes of the associated argument.



    \section{cmemorywrite - write to a combinatorial output memory}\label{sec:ipcmemorywrite}
        
        {\bf Call:} \\
        cmemorywrite(id, port, addr, data);
        
        {\bf Output parameters:} \\
        none
        
        {\bf Input parameters:} \\
        \begin{tabular}{| l | l | l | l |}
        \hline
        param  & mode     & type  & \\
        \hline \hline
        id   & M             &  -  & memory being written   \\ \hline
        port & I             & int & port                   \\ \hline
        addr & I V SV S Q PR & any & memory address         \\ \hline
        data & I V SV S Q PR & any & data to be written     \\ \hline
        \end{tabular}
        
        {\bf Parameter=argument allowed:} \\
        Yes
        
        {\bf In-line target execution:} \\
        yes
        
        {\bf Description:} \\
        Write to a port of a combinatorial output memory.

        {\em id} is the memory mode variable identifier.

        {\em port} is the memory port.

        {\em addr} is the address.

        {\em data} is the data to be written.
        
        The data and address may be any type, including compound types, but the
        types must match the types given when the combinatorial output memory was
        created (see {\bf cmemory()} \autorefph{sec:ipcmemory}) but primitive
        component widths may differ in which case padding or truncation will occur.

        Xilinx FPGAs do not allow writes to port 1. Within this limitation
        multiple reads and writes on a port are allowed and the address and
        data inputs will be multiplexed, however there is no access arbitration
        and it is up to the application code to ensure exclusive access.
        Coincident accesses will result in data corruption.

        A call to {\bf cmemorywrite()} (\autorefph{sec:ipcmemorywrite}) or {\bf
        cmemoryread()} (\autorefph{sec:ifcmemoryread}) sets the clock domain for
        the memory port to the current clock domain. All further calls to either of
        these inbuilt procedures for the same port of the same combinatorial output
        memory must be in the same clock domain (or a fatal error results).


    \section{designname - change the design name}\label{sec:ipdesignname}
        
        {\bf Call:} \\
        designname(name);
        
        {\bf Output parameters:} \\
        none
        
        {\bf Input parameters:} \\
        \begin{tabular}{| l | l | l | l |}
        \hline
        param  & mode & type & \\
        \hline \hline
        name   & I    & str  & design name string   \\ \hline
        \end{tabular}
        
        {\bf Parameter=argument allowed:} \\
        No
        
        {\bf In-line target execution:} \\
        no
        
        {\bf Description:} \\
        Change the design name, overriding the default derived from
        the file name. This can be overridden by the command line option
        {\bf -N}. It is also overridden if inbuilt procedure {\bf export()}
        is called.



    \section{directive - enter a directive}\label{sec:ipdirective}
        
        {\bf Call:} \\
        directive(key, value);
        
        {\bf Output parameters:} \\
        none
        
        {\bf Parameter=argument allowed:} \\
        No
        
        {\bf Input parameters:} \\
        \begin{tabular}{| l | l | l | l |}
        \hline
        param  & mode     & type  & \\
        \hline \hline
        key   & I & str & directive key                 \\ \hline
        value & I & any & directive value (optional)    \\ \hline
        \end{tabular}
        
        {\bf In-line target execution:} \\
        No
        
        {\bf Description:} \\
        There is a single global directive map. The key is a string.  The value
        may one of the 3PL immediate primitive types {\bf int}, {\bf log}, {\bf
        float} or {\bf str}. If the second argument is not provided the entry
        will be type {\bf log} and value {\bf true}.
        
        If the directive is already present its value will be changed but
        the type must be the same or a fatal error will occur.
               
        see also inbuilt procedures \\
        {\bf presetdirective()} \autorefph{sec:ippresetdirective} \\
        {\bf removedirective} \autorefph{sec:ipremovedirective}. \\
        and inbuilt functions \\
        {\bf directive()} \autorefph{sec:ifdirective}.\\
        {\bf direxists()} \autorefph{sec:ifdirexists}.\\
        {\bf directives()} \autorefph{sec:ifdirectives}.



    \section{divrem - combined quotient and remainder}\label{sec:ipdivrem}

        {\bf Prototype:} \\
        \vsp
        \begin{lstlisting}[gobble=8, language=threepl]
           procedure
           divrem (quot, rem <- num, denom) {}
        \end{lstlisting}
        
        {\bf Output parameters:} \\
        \begin{tabular}{| l | l | l | l |}
        \hline
        param  & mode     & type  & \\
        \hline \hline
        quot   & V & int uint & quotient    \\ \hline
        rem    & V & int uint & remainder   \\ \hline
        \end{tabular}
        
        {\bf Input parameters:} \\
        \begin{tabular}{| l | l | l | l |}
        \hline
        param  & mode     & type  & \\
        \hline \hline
        num    & I & int uint & numerator   \\ \hline
        denom  & I & int uint & denominator \\ \hline
        \end{tabular}
        
        {\bf Parameter=argument allowed:} \\
        Yes

        {\bf In-line target execution:} \\
        no

        {\bf description:} \\
        This procedure takes numerator and denominator values and
        computes the quotient and remainder. The results appear on
        output value variables since this procedure is combinatorial.
        The input and output types must be all{\em int} or all {\em uint}
        (widths may be different). The inputs may be immediate.

        {\nf Note: this procedure uses cascaded adders and subtractors
        and may have a very long data path delay!} However, the code is
        optimised for some cases where one operand is constant -  for
        division of 0 (return 0:0), division by 1 (return n:0) and
        division by power of 2 (quotient by shifting, remainder by
        masking).




    \section{element - add an element to the EDIF netlist}\label{sec:ipelement}
        
        {\bf Call:} \\
        element(name, id);
        
        {\bf Output parameters:} \\
        none
        
        {\bf Input parameters:} \\
        \begin{tabular}{| l | l | l | l |}
        \hline
        param  & mode     & type  & \\
        \hline \hline
        name & I & str & element name                     \\ \hline
        id   & I & str & optional unique block identifier \\ \hline
        \end{tabular}
        
        {\bf Parameter=argument allowed:} \\
        No
        
        {\bf In-line target execution:} \\
        No
        
        {\bf Description:} \\
        This procedure adds an element to the EDIF netlist. This is intended for
        special elements, perhaps particular to a specific FPGA family. For
        example global clock buffers are not considered generic enough to be
        included as fundamental logic elements so they are included using this
        procedure. A family-specific example would be the high-speed serial
        interfaces. Note that this is a low level process that bypasses 3PL
        execution streams. It is rather similar to embedding assembly code within
        a high level language.
        
        The first argument is the Xilinx library element name, e.g.
        {\bf element("BUFG")}.
        
        The second argument is an optional unique block identifier. If this
        argument is omitted then an identifier of the form "e123" will be alloted.
        {\bf Explicit identifiers must avoid the forms "e123" and "b123" to avoid
        conflicts with 3PL internally generated element identifiers.}
        
        Regardless of explicit use of the block identifer or not, the element
        constructed by this procedure remains as the current element until
        another call to this procedure. Connections to the element are made by
        the inbuilt procedures {\bf elementinput()}, {\bf elementclockinput()},
        {\bf elementoutput()} or {\bf elementtsoutput()} (see links below). The
        construction of an element must be completed by calling inbuilt
        procedure {\bf elementend()} \autorefph{sec:ipelementend}. Properties
        for the element are added by calling inbuilt procedure {\bf
        elementproperty()} \autorefph{sec:ipelementproperty}. Associated entries
        to the NCF file for the element can be made using inbuilt
        procedure {\bf elementaddncf()} \autorefph{sec:lpelementaddncf}.
        
        Note that this procedure can also be used to add standard elements to
        the netlist, e.g. {\bf element(, "FDRE")}.
        
        Elements defined by using {\bf element()} are added to cell definitions in
        the EDIF file. If another element of the same cell name is created they
        will share the same cell definition. The ports in the cell definition will
        be the set of all input, output or three-state ports specified in all
        instances of the element. Thus an instance of an element may not specify
        all the available ports for the element but the cell definition will
        accumulate all ports specified in the instances of the element. Any
        available ports not specified will be omitted from the cell definition.
        In the case of array ports where the port width
        changes from instance to instance the largest dimension will be used in
        the cell definition.
        
        For further description see \autorefph{chap:asfpgahe}.

        See also - \\
        {\bf elementend()} \autorefph{sec:ipelementend}, \\
        {\bf elementproperty()} \autorefph{sec:ipelementproperty}, \\
        {\bf elementinput()} \autorefph{sec:ipelmntin}, \\
        {\bf elementclockinput()} \autorefph{sec:ipelmntinclk} \\
        {\bf elementoutput()} \autorefph{sec:ipelmntout} \\
        {\bf elementtsoutput()} \autorefph{sec:ipelmntoutts}.




    \section{elementclockinput - connect a clock input signal to an EDIF netlist element}\label{sec:ipelmntinclk}
        
        {\bf Call:} \\
        elementclockinput(id, port\_name, signal);
        
        {\bf Output parameters:} \\
        none
        
        {\bf Input parameters:} \\
        \begin{tabular}{| l | l | l | l |}
        \hline
        param  & mode     & type  & \\
        \hline \hline
        id          & I          & str & optional block identifier \\ \hline
        signal      & V SV S Q I & any & optional signal           \\ \hline
        port\_name  & I          & str & port name                 \\ \hline
        \end{tabular}
        
        {\bf Parameter=argument allowed:} \\
        Yes
        
        {\bf In-line target execution:} \\
        No
        
        {\bf Description:} \\
        This procedure connects a clock input signal to an element (see {\bf
        element()} \autorefph{sec:ipelement}).

        The block identifier {\bf id} specifies the element to which a port is
        to be added. If this argument is omitted then the most recent element
        begun by a call to {\bf element} is assumed.

        Creation of multiple elements may be interleaved, but it is preferable
        to avoid interleaving and to complete an element before starting the
        next.

        {\bf port\_name} is a string identifying the element port to be
        connected.

        {\bf signal} is the signal to be connected to the element port. It must
        be clock mode or else an immediate type "log" whose value is false, in
        which case the clock input will be connected to GND. This argument may
        also be an immediate mode string in which case it refers to a single
        FPGA port (PAD), the string value being the pin designation. For a
        string "AB12" the port name will be "PORT\_AB12". The element port must
        be a single bit.

        If the {\bf signal} argument is missing or null the port will be
        connected to ground.

        This procedure should only be used for clock inputs to elements where
        the input functions as an edge trigger. Thus for inputs to clock buffers
        or the source inputs to clock managers {\bf elementinput()}
        \autorefph{sec:ipelmntin} should be used instead. This distinction was
        once required so that 3PL could keep track of clock destinations for
        timing constraint generation, however this is no longer the case. A
        remnant remains in the {\bf clk\_sinks} and {\bf clk\_CE} read-only
        attributes of clock mode variables.

        See also - \\
        {\bf element()} \autorefph{sec:ipelement}, \\
        {\bf elementend()} \autorefph{sec:ipelementend}, \\
        {\bf elementproperty()} \autorefph{sec:ipelementproperty}, \\
        {\bf elementinput()} \autorefph{sec:ipelmntin}, \\
        {\bf elementoutput()} \autorefph{sec:ipelmntout}, \\
        {\bf elementtsoutput()} \autorefph{sec:ipelmntoutts}.




    \section{elementend - complete construction of an EDIF netlist element}\label{sec:ipelementend}
        
        {\bf Call:} \\
        elementend(id);
        
        {\bf Output parameters:} \\
        none
        
        {\bf Input parameters:} \\
        \begin{tabular}{| l | l | l | l |}
        \hline
        param  & mode     & type  & \\
        \hline \hline
        id         & I    & str & optional block identifier \\ \hline
        \end{tabular}
        
        {\bf Parameter=argument allowed:} \\
        No
        
        {\bf In-line target execution:} \\
        No
        
        {\bf Description:} \\
        This procedure is called to complete entry of a netlist element.
        The optional argument gives the block identifier of the element
        to be completed. If omitted, the most recent uncompleted element
        is assumed.



    \section{elementinput - connect an input signal to an EDIF netlist element}\label{sec:ipelmntin}
        
        {\bf Call:} \\
        elementinput(id, port\_name, signal, arrayformat);
        
        {\bf Output parameters:} \\
        none
        
        {\bf Input parameters:} \\
        \begin{tabular}{| l | l | l | l |}
        \hline
        param  & mode     & type  & \\
        \hline \hline
        id          & I          & str & optional block identifier \\ \hline
        port\_name  & I          & str & port name                 \\ \hline
        signal      & V SV S Q I & any & signal                    \\ \hline
        arrayformat & I          & int & array format indicator    \\ \hline
        \end{tabular}
        
        {\bf Parameter=argument allowed:} \\
        Yes
        
        {\bf In-line target execution:} \\
        No
        
        {\bf Description:} \\
        This procedure connects an input signal to an element (see {\bf element()}
        \autorefph{sec:ipelement}).
        
        The block identifier {\bf id} specifies the element to which a port is to
        be added. If this argument is omitted then the most recent
        element begun by a call to {\bf element} is assumed.
        
        Creation of multiple elements may be interleaved, but it is preferable to
        avoid interleaving and to complete an element before starting the next.
        
        {\bf port\_name} is a string identifying the element port to be
        connected. The port width is determined by the signal width.

        
        {\bf signal} is the signal to be connected to the element port. It may be
        value, selectvalue, or static and may be any type. If value mode the value
        must not contain any queue reads. This argument may also be an immediate mode
        string in which case it refers to a single FPGA port (PAD), the string value
        being the pin designation. For a string "AB12" the port name will be
        "PORT\_AB12". If the signal is more than one bit wide an array format
        argument must be provided.

        If the {\bf signal} argument is missing or null no port connection is
        made.
               
        {\bf arrayformat} indicates the form of ports which are arrays.
        If this argument is 1 it indicates that array ports should be
        expanded to individual pins of the form A0, A1, A2 etc. If the value
        is 2 then the port will be a formal EDIF array. ISE will accept format
        1 for all elements. Some versions of ISE (e.g. 10) give errors when
        format 2 is used for block RAM elements where different port sizes
        are used in the same configuration. Vivado requires format 2 for port
        arrays. Since Vivado cannot be used for earlier FPGAs, block RAMs for
        these (up to XCS6, XCV6) are defined using format 1. Series 7 block RAMs
        are defined using format 2 since early versions of ISE do not support
        these devices. A single bit port can have a format type of 1, but not
        type 2.
        
        In the case of clock mode inputs this procedure should only be used for
        inputs to elements where the input does not function as an edge trigger, e.g.
        for inputs to clock buffers or the source inputs to clock managers. For edge
        trigger inputs {\bf elementclockinput()} \autorefph{sec:ipelmntinclk} should
        be used instead. This distinction was once required so that 3PL could keep
        track of clock destinations for timing constraint generation, however this
        is no longer the case. A remnant remains in the {\bf clk\_sinks} and
        {\bf clk\_CE} read-only attributes of clock mode variables.

        As an example of port input, the code -
        
        \begin{lstlisting}[gobble=8, language=threepl]
           input(is, "log", loc="AB12", tnm="IC");
        \end{lstlisting}
        
        could be effectively replaced by -
        
        \begin{lstlisting}[gobble=8, language=threepl]
           value(is, "log");

           ibuf(is <- "AB12", "IC");

           module
           ibuf (value(out) <- str(pin), str(tnm)) {
               element("IBUF");
               elementinput(, "I", pin);
               elementoutput(out <- , "O");
               elementStrProperty(, "loc", pin);
               elementend();
               IOconstraint(pin, "tnm", "IC");
           }
        \end{lstlisting}
        
        though note that the variable {\bf is} would now be value mode, not
        input mode.

        See also - \\
        {\bf element()} \autorefph{sec:ipelement}, \\
        {\bf elementend()} \autorefph{sec:ipelementend}, \\
        {\bf elementproperty()} \autorefph{sec:ipelementproperty}, \\
        {\bf elementclockinput()} \autorefph{sec:ipelmntinclk} \\
        {\bf elementoutput()} \autorefph{sec:ipelmntout}, \\
        {\bf elementtsoutput()} \autorefph{sec:ipelmntoutts}.





    \section{elementoutput - connect an output signal to an EDIF netlist element}\label{sec:ipelmntout}
        
        {\bf Call:} \\
        elementoutput(signal \verb'<-' id, port\_name, arrayformat);
        
        {\bf Output parameters:} \\
        \begin{tabular}{| l | l | l | l |}
        \hline
        param  & mode     & type  & \\
        \hline \hline
        signal      & V SV S Q I & any & signal                    \\ \hline
        \end{tabular}
        
        {\bf Input parameters:} \\
        \begin{tabular}{| l | l | l | l |}
        \hline
        param  & mode     & type  & \\
        \hline \hline
        id          & I          & str & optional block identifier \\ \hline
        port\_name  & I          & str & port name                 \\ \hline
        arrayformat & I          & int & array format indicator    \\ \hline
        \end{tabular}
        
        {\bf Parameter=argument allowed:} \\
        Yes
        
        {\bf In-line target execution:} \\
        No
        
        {\bf Description:} \\
        This procedure connects an output signal to an element (see {\bf element()}
        \autorefph{sec:ipelement}).
        
        The block identifier {\bf id} specifies the element to which a port is to
        be added. If this argument is omitted then the most recent
        element begun by a call to {\bf element} is assumed.
        
        Creation of multiple elements may be interleaved, but it is preferable to
        avoid interleaving and to complete an element before starting the next.
        
        {\bf signal} is the signal to be connected to the element port. It must be
        value mode, of any type, or clock mode. This argument may also be an
        immediate mode string in which case it refers to a single FPGA port (PAD),
        the string value being the pin designation. For a string "AB12" the port
        name will be "PORT\_AB12" (see input example in {\bf elementinput()}
        \autorefph{sec:ipelmntin}). If the signal is more than one bit wide an
        array format argument must be provided.        
        
        {\bf port\_name} is a string identifying the element port to be
        connected. The port width is determined by the signal width.
        
        {\bf arrayformat} indicates the form of ports which are arrays.
        If this argument is 1 it indicates that array ports should be
        expanded to individual pins of the form A0, A1, A2 etc. If the value
        is 2 then the port will be a formal EDIF array. ISE will accept format
        1 for all elements. Some versions of ISE (e.g. 10) give errors when
        format 2 is used for block RAM elements where different port sizes
        are used in the same configuration. Vivado requires format 2 for port
        arrays. Since Vivado cannot be used for earlier FPGAs, block RAMs for
        these (up to XCS6, XCV6) are defined using format 1. Series 7 block RAMs
        are defined using format 2 since early versions of ISE do not support
        these devices. A single bit port can have a format type of 1, but not
        type 2.

        See also - \\
        {\bf element()} \autorefph{sec:ipelement}, \\
        {\bf elementend()} \autorefph{sec:ipelementend}, \\
        {\bf elementproperty()} \autorefph{sec:ipelementproperty}, \\
        {\bf elementinput()} \autorefph{sec:ipelmntin}, \\
        {\bf elementclockinput()} \autorefph{sec:ipelmntinclk} \\
        {\bf elementtsoutput()} \autorefph{sec:ipelmntoutts}.






    \section{elementproperty - set an element property}\label{sec:ipelementproperty}
        
        {\bf Call:} \\
        elementproperty(id, property, value);
        
        {\bf Output parameters:} \\
        none
        
        {\bf Input parameters:} \\
        \begin{tabular}{| l | l | l | l |}
        \hline
        param  & mode     & type  & \\
        \hline \hline
        id       & I        & str & optional block identifier \\ \hline
        property & I        & str & property name             \\ \hline
        value    & I        & str & property value            \\ \hline
        \end{tabular}
        
        {\bf Parameter=argument allowed:} \\
        Yes
        
        {\bf In-line target execution:} \\
        No
        
        {\bf Description:} \\
        This procedure sets a property value for an element.
        The optional argument {\bf id} gives the block identifier of the element
        to be completed. If omitted, the most recent uncompleted element
        is assumed.

        {\bf elementproperty()} only takes strings as value arguments so a number
        of library procedures are available in xilinx/common.3pl to format other
        types. These are - \\
        {\bf elementLogProperty()} \autorefph{sec:lpelementLogProperty}, \\
        {\bf elementIntProperty()} \autorefph{sec:lpelementIntProperty}, \\
        {\bf elementBinProperty()} \autorefph{sec:lpelementBinProperty}, \\
        {\bf elementHexProperty()} \autorefph{sec:lpelementHexProperty}, \\
        {\bf elementFloatProperty()} \autorefph{sec:lpelementFloatProperty}, \\
        {\bf elementStrProperty()} \autorefph{sec:lpelementStrProperty}.
        
        These procedures also provide limit checks or valid argument lists.

        See also - \\
        {\bf element()} \autorefph{sec:ipelement}, \\
        {\bf elementend()} \autorefph{sec:ipelementend}, \\
        {\bf elementinput()} \autorefph{sec:ipelmntin}, \\
        {\bf elementclockinput()} \autorefph{sec:ipelmntinclk} \\
        {\bf elementoutput()} \autorefph{sec:ipelmntout}.




    \section{elementtsoutput - connect a 3-state output signal to an EDIF netlist element}\label{sec:ipelmntoutts}
        
        {\bf Call:} \\
        elementtsoutput(signal \verb'<-' id, port\_name, arrayformat);
        
        {\bf Output parameters:} \\
        \begin{tabular}{| l | l | l | l |}
        \hline
        param  & mode     & type  & \\
        \hline \hline
        signal      & V SV S Q I & any & signal                    \\ \hline
        \end{tabular}
        
        {\bf Input parameters:} \\
        \begin{tabular}{| l | l | l | l |}
        \hline
        param  & mode     & type  & \\
        \hline \hline
        id          & I          & str & optional block identifier \\ \hline
        port\_name  & I          & str & port name                 \\ \hline
        arrayformat & I          & int & array format indicator    \\ \hline
        \end{tabular}
        
        {\bf Parameter=argument allowed:} \\
        Yes
        
        {\bf In-line target execution:} \\
        No
        
        {\bf Description:} \\
        This procedure connects a 3-state output signal to an element (see {\bf element()}
        \autorefph{sec:ipelement}).
        
        The block identifier {\bf id} specifies the element to which a port is to
        be added. If this argument is omitted then the most recent
        element begun by a call to {\bf element} is assumed.
        
        Creation of multiple elements may be interleaved, but it is preferable to
        avoid interleaving and to complete an element before starting the next.
        
        {\bf signal} is the signal to be connected to the element port. It must be
        selectvalue mode of any type. If the signal is more than one bit wide an
        array format argument must be provided. This argument may also be an
        immediate mode string in which case it refers to a single FPGA port (PAD),
        the string value being the pin designation. For a string "AB12" the port name
        will be "PORT\_AB12".
        
        {\bf port\_name} is a string identifying the element port to be connected. If
        the signal is more than one bit wide an array format argument must be
        provided. The port width is determined by the signal width.

        
        {\bf arrayformat} indicates the form of ports which are arrays.
        If this argument is 1 it indicates that array ports should be
        expanded to individual pins of the form A0, A1, A2 etc. If the value
        is 2 then the port will be a formal EDIF array. ISE will accept format
        1 for all elements. Some versions of ISE (e.g. 10) give errors when
        format 2 is used for block RAM elements where different port sizes
        are used in the same configuration. Vivado requires format 2 for port
        arrays. Since Vivado cannot be used for earlier FPGAs, block RAMs for
        these (up to XCS6, XCV6) are defined using format 1. Series 7 block RAMs
        are defined using format 2 since early versions of ISE do not support
        these devices. A single bit port can have a format type of 1, but not
        type 2.

        See also - \\
        {\bf element()} \autorefph{sec:ipelement}, \\
        {\bf elementend()} \autorefph{sec:ipelementend}, \\
        {\bf elementproperty()} \autorefph{sec:ipelementproperty}, \\
        {\bf elementinput()} \autorefph{sec:ipelmntin}, \\
        {\bf elementclockinput()} \autorefph{sec:ipelmntinclk} \\
        {\bf elementoutput()} \autorefph{sec:ipelmntout}.





    \section{emptyvalue - clear a value-mode variable}\label{sec:ipemptyvalue}
        
        {\bf Call:} \\
        emptyvalue(name);
        
        {\bf Output parameters:} \\
        none
        
        {\bf Input parameters:} \\
        \begin{tabular}{| l | l | l | l |}
        \hline
        param  & mode     & type  & \\
        \hline \hline
        name   & V        & -     & value variable identifier \\ \hline
        \end{tabular}
        
        {\bf Parameter=argument allowed:} \\
        No
        
        {\bf In-line target execution:} \\
        No
        
        {\bf Description:} \\
        This procedure clears the value of a value-mode variable and changes
        its type to "empty". This allows the variable to be subsequently
        assigned a value of any type. This is particularly useful where
        a value variable is to be supplied as an argument to an output
        parameter where the type of that parameter is not known prior to
        the call.


    \section{enum - create an enumerated type}\label{sec:ipenum}
        \index{type!enum}
        
        {\bf Call:} \\
        enum(name, id1, id2,\dots);
        
        {\bf Output parameters:} \\
        none
        
        {\bf Input parameters:} \\
        \begin{tabular}{| l | l | l | l |}
        \hline
        param  & mode     & type  & \\
        \hline \hline
        name &   &     & type identifier        \\ \hline
        id1  & I &     & 1st value identifier   \\ \hline
        id2  & I &     & 2nd value identifier   \\ \hline
         .   & I &     & ...                    \\ \hline
        \end{tabular}
        
        {\bf Parameter=argument allowed:} \\
        No
        
        {\bf In-line target execution:} \\
        No
        
        {\bf Description:} \\
        The {\bf enum()} procedure creates an enumerated type where the ordinal
        values are 0, 1, 2 etc (see also \autorefph{sec:ipenumv}).

        The first argument, {\bf name}, is an identifier giving the enum name
        which will be created as a variable of type "type". This is followed by
        a number of value identifiers which will have associated ordinal values
        0, 1, 2 etc.

        The identifiers must be unique within this enum, but the
        identifiers may be used in other enums or as variable identifiers
        without conflict.
        
        Enumerated type variables are created using the enumerated type,
        e.g.
        \begin{lstlisting}[gobble=8, language=threepl]
           enum(e1, none, typea, typeb);
           var(v1, e1);
           static(s1, e1);
        \end{lstlisting}
        The enumerated type can be used for immediate, value, selectvalue,
        static or queue mode variables.
        
        Enumerated type values are accessed by using the type name and the
        value identifier, e.g.
        \begin{lstlisting}[gobble=8, language=threepl]
           when (s1 == e1.typea)
               s1 <- e1.typeb;
        \end{lstlisting}
        
        There are a number of inbuilt functions for handling enumerated
        types -

        \begin{tabular}{| l | l |}
        \hline
        ord(v)     \autorefph{sec:iford}     & the ordinal value of variable v \\
        \hline
        next(v)    \autorefph{sec:ifnext}    & the next value of variable v \\
        \hline
        prev(v)    \autorefph{sec:ifprev}    & the previous value of variable v \\
        \hline
        hasnext(v) \autorefph{sec:ifhasnext} & true if variable v has a next
                                     value (i.e. it is not the last
                                     value) \\
        \hline
        hasprev(v) \autorefph{sec:ifhasprev} & true if variable v has a previous
                                     value (i.e. it is not the first
                                     value) \\
        \hline
        first(x)   \autorefph{sec:iffirst}   & the first ordered value of
                                     variable or enum type x \\
        \hline
        last(x)    \autorefph{sec:iflast}    & the last ordered value of
                                     variable or enum type x \\
        \hline
        \end{tabular}

        The identifier of an enum value can be obtained by using
        the inbuilt function {\bf tostring()} \autorefph{sec:iftostring}.


    \section{enumv - create an enumerated type}\label{sec:ipenumv}
        \index{type!enum}
        
        {\bf Call:} \\
        enumv(name, id1, v1, id2, v2,\dots);
        
        {\bf Output parameters:} \\
        none
        
        {\bf Input parameters:} \\
        \begin{tabular}{| l | l | l | l |}
        \hline
        param  & mode     & type  & \\
        \hline \hline
        name &   &     & type identifier        \\ \hline
        id1  &   &     & 1st value identifier   \\ \hline
        v1   & I & int & 1st value ordinal      \\ \hline
        id2  &   &     & 2nd value identifier   \\ \hline
        v2   & I & int & 2nd value ordinal      \\ \hline
         .   &   &     & ...                    \\ \hline
        \end{tabular}
        
        {\bf Parameter=argument allowed:} \\
        No
        
        {\bf In-line target execution:} \\
        No
        
        {\bf Description:} \\
        The {\bf enumv()} procedure creates an enumerated type where
        the ordinal values are explicitly given (see also \autorefph{sec:ipenum}).
               
        The first argument, {\bf name}, 
        is an identifier giving the enum name which will be created
        as a variable of type "type". This is followed by a number of
        pairs of identifiers and associated ordinal values. The ordinal
        value arguments must be immediate expressions of type int which
        are not negative.

        The identifiers and ordinals must be unique within this
        enum, but the identifiers may be used in other enums
        or as variable identifiers without conflict.
        
        Enumerated type variables are created using the enumerated type,
        e.g.
        \begin{lstlisting}[gobble=8, language=threepl]
           enum(e1,
               none,   3,
               typea,  7,
               typeb,  5
           );
           var(v1, e1);
               static(s1, e1);
        \end{lstlisting}
        The enumerated type can be used for immediate, value, selectvalue,
        static or queue mode variables.
        
        Enumerated type values are accessed by using the type name and the
        value identifier, e.g.
        \begin{lstlisting}[gobble=8, language=threepl]
           when (s1 == e1.typea)
               s1 <- e1.typeb;
        \end{lstlisting}
        
        There are a number of inbuilt functions for handling enumerated
        types -

        \begin{tabular}{| l | l |}
        \hline
        ord(v)     \autorefph{sec:iford}     & the ordinal value of variable v \\
        \hline
        next(v)    \autorefph{sec:ifnext}    & the next value of variable v \\
        \hline
        prev(v)    \autorefph{sec:ifprev}    & the previous value of variable v \\
        \hline
        hasnext(v) \autorefph{sec:ifhasnext} & true if variable v has a next
                                     value (i.e. it is not the last
                                     value) \\
        \hline
        hasprev(v) \autorefph{sec:ifhasprev} & true if variable v has a previous
                                     value (i.e. it is not the first
                                     value) \\
        \hline
        first(x)   \autorefph{sec:iffirst}   & the first ordered value of
                                     variable or enum type x \\
        \hline
        last(x)    \autorefph{sec:iflast}    & the last ordered value of
                                     variable or enum type x \\
        \hline
        \end{tabular}

        The identifier of an enum value can be obtained by using
        the inbuilt function {\bf tostring()} \autorefph{sec:iftostring}.


    \section{exec - execute a command}\label{sec:ipexec}
        
        {\bf Call:} \\
        exec(status, out, err \verb'<-' command, environ, dir, in);
        
        {\bf Output parameters:} \\
        \begin{tabular}{| l | l | l | l |}
        \hline
        param  & mode     & type  & \\
        \hline \hline
        status & I & int, uint & status returned from command     \\ \hline
        out    & I & str       & standard output from command     \\ \hline
        err    & I & str       & standard error from command     \\ \hline
        \end{tabular}
        
        {\bf Input parameters:} \\
        \begin{tabular}{| l | l | l | l |}
        \hline
        param  & mode     & type  & \\
        \hline \hline
        command & I & []str & command with arguments to be executed        \\ \hline
        environ & I & map   & optional environment variables               \\ \hline
        dir     & I & str   & optional working directory                   \\ \hline
        in      & I & str   & optional string to standard input to command \\ \hline
        \end{tabular}
        
        {\bf Parameter=argument allowed:} \\
        Yes
        
        {\bf In-line target execution:} \\
        No
        
        {\bf Description:} \\
        This procedure executes a command with arguments, given by the first
        input parameter {\bf command}. This is an array of strings, the first
        giving the command path and the remainder supplying arguments.

	Note that if the command path is not absolute, the {\bf PATH} variable
	in the environment on the call to 3PL will be use to find the
	executable file, {\em not} that in the optional {\bf environ} input
	parameter.

        No shell is executed. See {\bf shell()} \autorefph{sec:ipshell} if this
	is desired.
        
        The optional second input parameter, {\bf environ}, is a map of environment
        variables. For each entry the key is the environment variable and the
        associated map value is the variable value. Value types are restricted
        to "int", "uint", "str" and "float", other types resulting in a fatal error.
        If no argument is supplied for this parameter the environment on the
        call to 3PL is passed to the command.
        
        The optional third input parameter, {\bf dir}, is a string specifying the working
        directory. If no argument is supplied for this parameter the current
        directory on the call to 3PL is used.
        
        The optional fourth input parameter, {\bf in}, is a string to be copied to
        standard input of the command process.
        
        If the first output parameter, {\bf status}, is given an argument, the
        status returned from the command is copied to this variable. The argument
        must be immediate mode and type "int" or "uint".
        
        If the second output parameter, {\bf out}, is given an argument, the
        standard output from the command is copied to this variable. The argument
        must be immediate mode and type "str". The resulting string is usually
        one or more lines, each terminated by a newline character.
        
        If the third output parameter, {\bf err}, is given an argument, the
        standard error from the command is copied to this variable. The argument
        must be immediate mode and type "str". The resulting string is usually
        one or more lines, each terminated by a newline character.
        
        If the command cannot be found, or cannot be executed, a fatal error occurs.
        
        This procedure waits until the command exits. If the command is signalled
        then this procedure continues as for a normal exit.

        See also {\bf shell()} \autorefph{sec:ipshell}.


    \section{exit - terminate execution}\label{sec:ipexit}
        
        {\bf Call:} \\
        exit(message);
        
        {\bf Output parameters:} \\
        none
        
        {\bf Input parameters:} \\
        \begin{tabular}{| l | l | l | l |}
        \hline
        param  & mode     & type  & \\
        \hline \hline
        message & I & str & optional message   \\ \hline
        \end{tabular}
        
        {\bf Parameter=argument allowed:} \\
        No
        
        {\bf In-line target execution:} \\
        No
        
        {\bf Description:} \\
        Terminate 3PL immediate execution. If the argument is provided
        a message is printed on standard out.

        See also {\bf texit()} \autorefph{sec:iptexit}, {\bf trace()} \autorefph{sec:iptrace}.


    \section{export - export a logic 'core'}\label{sec:ipexport}
        
        {\bf Call:} \\
        export(OP1=o1, OP2=o2, \dots \verb'<-' IP1=o1, IP2=i2, \dots);
        
        {\bf Output parameters:} \\
        See description
        
        {\bf Input parameters:} \\
        See description
        
        {\bf Parameter=argument allowed:} \\
        No
        
        {\bf In-line target execution:} \\
        No
        
        {\bf Description:} \\
        This procedure signals the compiler that the logic is a 'core'
        (or macro) to be exported so that its netlist can be linked into the
        netlists from other programs (not necessarily from 3PL programs).        
        
        The -e option must be used in the call to 3PL to supply the export core
        name, otherwise a fatal error will occur. This string will also be used
        to name the netlist and constraint output files.
        
        The input arguments are signal inputs to the 'core'. These arguments may
        be passed using the key-matching syntax. The identifier on the LHS of
        the '=' is the port name used when linking the netlists. The RHS of the
        '=' is a 3PL variable to be an input from the port. If the argument is a
        single identifier then it is taken as both the port identifier and the
        identifier of a variable driving the port.
        
        Input port variables may be {\bf value} or {\bf clock} mode. {\bf Value}
        mode variables may be any target type. Input port variables which are
        mode {\bf value} must have their clock domains established by assigning
        attribute {\bf domain} (see \autorefph{sec:attributes}) before the call
        to {\bf export()} which passes them as input arguments.
        
        Output arguments are signal outputs from the 'core'. These arguments
        may also be passed using the key-matching syntax. The identifier
        on the LHS of the '=' is the port name used when linking the netlists.
        The RHS of the '=' is a 3PL variable or expression to drive the port.
        
        Output port variables can be mode {\bf value}, {\bf selectvalue}, {\bf
        static}, {\bf priority} or {\bf clock}. A {\bf selectvalue} mode output
        port variable is used for a 3-state port. Non-clock output port
        variables may be any target type. These variables must not contain queue
        reads. If the argument is a single identifier then it is taken as both
        the port identifier and the identifier of a variable driving the
        port.    
        


    \section{external - declare a signal as an EDIF port}\label{sec:ipexternal}
        
        {\bf Call:} \\
        external(portname, var, direction);
        
        {\bf Output parameters:} \\
        None.
        
        {\bf Input parameters:} \\
        {\bf portname} is the identifier to be used for the EDIF port name ("str")\\
        {\bf var} is a target mode variable providing he signal to the external port.\\
        {\bf direction} is "in" or "out" to select the port direction.
        
        {\bf Parameter=argument allowed:} \\
        Yes
        
        {\bf In-line target execution:} \\
        No
        
        {\bf Description:} \\
        This procedure connects a signal to an external EDIF port in the netlist file.
        
        For an input port the variable mode must be VALUE or CLOCK.
        
        For an output port the variable mode must be VALUE, STATIC or CLOCK.
        For a value mode variable the call to {\bf external()} must occur
        after the variable is assigned.
        


    \section{file - create file variables}\label{sec:ipfile}
        \index{variable!file}
        
        {\bf Call:} \\
        file(id, path);
        
        {\bf Output parameters:} \\
        none
        
        {\bf Input parameters:} \\
        \begin{tabular}{| l | l | l | l |}
        \hline
        param  & mode & type  & \\
        \hline \hline
        ident & - & -   & variable identifier \\ \hline
        path  & I & str & file path string    \\ \hline
        \end{tabular}
        
        {\bf Parameter=argument allowed:} \\
        No
        
        {\bf In-line target execution:} \\
        No
        
        {\bf Description:} \\
        This procedure creates a variable of type {\bf file}. {\bf ident} is
        the variable identifier. The path name may be relative to the current
        directory or may be absolute (a leading \verb'\').
        
        A file must be opened by calling inbuilt procedure
        {\bf open()}(\autorefph{sec:ipopen} before it can be read or written.
        
        A number of files are already open as global variables. {\bf stdin} is the
        standard input, {\bf stdout} and {\bf stderr} are standard input and standard error,
        and {\bf rpt} is the 3PL report file. The report file is {\bf design\_name.rpt} if
        the 3PL was called with the -r option, otherwise it is {\bf stdout}.
        
        Files are read as characters by \\
        {\bf readln()}(\autorefph{sec:ifreadln}) \\
        {\bf scanf()}(\autorefph{sec:ipscanf}) \\
        {\bf fscanf()}(\autorefph{sec:ipfscanf}) \\   
        and are read as bytes by \\
        {\bf readbyte()}(\autorefph{sec:ifreadbyte}).
        
        Files are written as characters by \\
        {\bf writeln()}(\autorefph{sec:ipwriteln}) \\
        {\bf printf()}(\autorefph{sec:ipprintf}) \\
        {\bf fprintf()}(\autorefph{sec:ipfprintf}) \\
        and are written as bytes by \\
        {\bf writebyte()}(\autorefph{sec:ipwritebyte}).
        
        The two modes cannot be mixed; a file read in
        one mode cannot be subsequently read in the opposite mode, and similarly
        for writing.
        
        The open file can be closed by calling inbuilt procedure
        {\bf close()}(\autorefph{sec:ipclose}).

        See also inbuilt functions - \\
        {\bf filecanread()}      (\autorefph{sec:iffilecanread}) \\
        {\bf filecanwrite()}     (\autorefph{sec:iffilecanwrite}) \\
        {\bf filedelete()}       (\autorefph{sec:iffiledelete}) \\
        {\bf fileexists()}       (\autorefph{sec:iffileexists}) \\
        {\bf filegetparent()}    (\autorefph{sec:iffilegetparent}) \\
        {\bf fileisdir()}        (\autorefph{sec:iffileisdir}) \\
        {\bf fileisfile()}       (\autorefph{sec:iffileisfile}) \\
        {\bf filelastmod()}      (\autorefph{sec:iffilelastmod}) \\
        {\bf filelength()}       (\autorefph{sec:iffilelength}) \\
        {\bf filelist()}         (\autorefph{sec:iffilelist}) \\
        {\bf filelistdirs()}     (\autorefph{sec:iffilelistdirs}) \\
        {\bf filelistfiles()}    (\autorefph{sec:iffilelistfiles}) \\
        {\bf filemkdir()}        (\autorefph{sec:iffilemkdir}) \\
        {\bf filerename()}       (\autorefph{sec:iffilerename}) \\
        {\bf filesetreadonly()}  (\autorefph{sec:iffilesetreadonly}) \\
        {\bf hadeof()}           (\autorefph{sec:ifhadeof}) \\
        {\bf readbyte()}         (\autorefph{sec:ifreadbyte}) \\
        {\bf readln()}           (\autorefph{sec:ifreadln}) \\
        
        See also inbuilt procedures - \\
        {\bf fprintf()}          (\autorefph{sec:ipfprintf}) \\
        {\bf fscanf()}           (\autorefph{sec:ipfscanf})) \\ \\
 


    \section{float - create one or more immediate floating point variables}\label{sec:ipfloat}
        \index{variable!immediate!float}
        
        {\bf Call:} \\
        float(ident, init);
        
        {\bf Output parameters:} \\
        none
        
        {\bf Input parameters:} \\
        \begin{tabular}{| l | l | l | l |}
        \hline
        param  & mode & type & \\
        \hline \hline
        ident or \verb'{'id1, id2... \verb'}' & - &  -  & variable identifier(s)   \\ \hline
        init                   & I & float & optional initial value \\ \hline
        \end{tabular}
        
        {\bf Parameter=argument allowed:} \\
        No
        
        {\bf In-line target execution:} \\
        No
        
        {\bf Description:} \\
        Create one or more immediate floating point variables. If more than one
        variable is to be created the list of variables must be in braces
        (\{\dots\}). An integer or floating point expression may be provided as
        the second argument to initialise the variables, otherwise they will be
        initially zero.



    \section{fprintf - print a formatted line to an output file}\label{sec:ipfprintf}

        {\bf Call:} \\
        fprintf(file \verb'<-' fmt, v1, v2,....);
        
        {\bf Output parameters:} \\
        \begin{tabular}{| l | l | l | l |}
        \hline
        param  & mode     & type  & \\
        \hline \hline
        file & I & file & output file   \\ \hline
        \end{tabular}
        
        {\bf Input parameters:} \\
        \begin{tabular}{| l | l | l | l |}
        \hline
        param  & mode     & type  & \\
        \hline \hline
        fmt  & I & str  & format string   \\ \hline
        v1   & I & any  & input variable  \\ \hline
        v2   & I & any  & input variable  \\ \hline
        .    & I & any  & input variable  \\ \hline
        .    & I & any  & input variable  \\ \hline
        \end{tabular}
        
        {\bf Parameter=argument allowed:} \\
        No
        
        {\bf In-line target execution:} \\
        No
        
        {\bf Description:} \\
        This procedure is equivalent to fprintf() in C. The format string is that used
        by the Java Formatter class.        
        
        The format argument may be followed by a single argument which is an array of
        pointers, for example {\bf inargs} in the calling module, procedure or function.
        In that case the pointers are resolved and provide the input variable list to be
        printed.
        
        The output file must have already been opened.
        
        See the note on character and byte reads and writes in the description
        of inbuilt procedure {\bf file()} \autorefph{sec:ipfile}.
        
        See also - \\
        {\bf printf()} \autorefph{sec:ipprintf} \\
        {\bf sprintf()} \autorefph{sec:ipsprintf} \\
        {\bf writeln()} \autorefph{sec:ipwriteln} \\
        {\bf writebyte()} \autorefph{sec:ipwritebyte}.



    \section{fscanf - scan a line from a file into variables}\label{sec:ipfscanf}

        {\bf Call:} \\
        fscanf(v1, v2,.... \verb'<-' file, fmt);
        
        {\bf Output parameters:} \\
        \begin{tabular}{| l | l | l | l |}
        \hline
        param  & mode     & type  & \\
        \hline \hline
        v1 & I & any & output variable  \\ \hline
        v2 & I & any & output variable  \\ \hline
        .  & I & any & output variable  \\ \hline
        .  & I & any & output variable  \\ \hline
        \end{tabular}
        
        {\bf Input parameters:} \\
        \begin{tabular}{| l | l | l | l |}
        \hline
        param  & mode     & type  & \\
        \hline \hline
        file & I & str & an input file  \\ \hline
        fmt  & I & str & format string  \\ \hline
        \end{tabular}
        
        {\bf Parameter=argument allowed:} \\
        No
        
        {\bf In-line target execution:} \\
        No
        
        {\bf Description:} \\
        This procedure scans a line from an input file extracting values from fields
        and assigning these to the output variables.
        
        The input file must have already been opened.
        
        The scanned line must consist of fields (tokens) separated by delimiters of commas
or spaces (any number of either, i.e. the delimiter pattern is [, ]+ ).

The format string is not the same as that used in scanf, sscanf or fscanf in C.
It is a string of single characters, one per conversion. Allowed
conversions are -

\begin{tabular}{| l | l |} 
\hline
'd' & integer                                   \\ \hline
'u' & unsigned integer                          \\ \hline
'x' & hexadecimal integer                       \\ \hline
'o' & octal integer                             \\ \hline
'f' & floating point value                      \\ \hline
's' & string                                    \\ \hline
'.' & convert according to the variable type    \\ \hline
'*' & skip the token                            \\ \hline
' ' & ignored                                   \\ \hline
tab & ignored                                   \\ \hline
\end{tabular}

The format string must not contain any other characters.
Spaces or tabs can be included anywhere if desired. Because the scanner cannot match explicit
fields from the format string and there are no conversion paramaters such as field widths or
precisions there is no need for a \verb"'%'" escape character. The number of actual
conversions (i.e. not including the skip) must be equal to the number of output variables.

The format argument may be null or "" in which case tokens will be converted and
assigned to the variables based on the variable types.

Tokens that are to be converted as hexadecimal values must not have a leading 0x.
Similarly a leading 0 will not be interpreted as indicating an octal value. The
conversion character alone decides which radix will be used.

        
        See the note on character and byte reads and writes in the description
        of inbuilt procedure {\bf file()} \autorefph{sec:ipfile}.
       
        See also - \\
        {\bf scanf()} \autorefph{sec:ipscanf} \\
        {\bf sscanf()} \autorefph{sec:ipsscanf} \\
        {\bf readln()} \autorefph{sec:ifreadln} \\
        {\bf readbyte()} \autorefph{sec:ifreadbyte}





    \section{import - import a 'core'}\label{sec:ipimport}
        
        {\bf Call:} \\
        import(OP1=o1, OP2=o2, \dots \verb'<-' name, bname, IP1=o1, IP2=i2, \dots);
        
        {\bf Output parameters:} \\
        See description
        
        {\bf Input parameters:} \\
        See description
        
        {\bf Parameter=argument allowed:} \\
        No
        
        {\bf In-line target execution:} \\
        No
        
        {\bf Description:} \\
        This procedure signals the compiler that an external 'core'
        (or macro) is to be linked into this program. The core is
        not necessarily one created by a 3PL program.        
        
        The first input argument must be a string expression giving the
        name of the 'core'.
        
        The second argument is immediate mode and type "str" and is taken as an
        optional netlist block name for the imported core element. This
        identifier must begin with an alphabetic character and contain only
        alphabetic characters, digits and underscores. In addition it should not
        consist of 'b' followed by a string of digits since that is the block
        naming scheme used by the compiler to generate unique block names and
        hence the use of such a string in this context may result in a clash. If
        this argument is omitted the block will be labeled with a string of the
        form "bnnnn".
        
        The remaining input arguments are signal inputs to the 'core'. These
        arguments may be passed using the key-matching syntax. The identifier on
        the LHS of the '=' is the port name used when linking the netlists. The
        RHS of the '=' is a 3PL variable or expression to drive the port. If the
        argument is a single identifier then it is taken as both the port
        identifier and the identifier of a variable driving the port.
        
        Input port variables can be mode {\bf input}, {\bf value}, {\bf
        selectvalue}, {\bf static}, {\bf priority} or {\bf clock}. Input port
        variables may be any target type.
        
        Output arguments are signal outputs from the 'core'. These arguments may
        be passed using key-matching syntax. The identifier on the LHS of the
        '=' is the port name used when linking the netlists. The RHS of the '='
        is a 3PL variable to be an output from the port. Value variables must
        not contain queue reads. If the argument is a single identifier then it
        is taken as both the port identifier and the identifier of a variable
        driven by the port.    
        
        Output port variables may be mode {\bf value}, {\bf selectvalue}
        (3-state output) or {\bf clock}. Value and selectvalue mode variables
        may be any target type. Output port variables which are mode {\bf value}
        or mode {\bf selectvalue} will have null clock domains. they may have
        their clock domains established by assigning attribute {\bf domain} (see
        \autorefph{sec:attributes}) after the call to {\bf import()} where they
        are passed as output arguments.
        
        The number of bits in each argument must exactly match the associated
        core port width.

        All input ports must be connected but connections to unwanted output
        ports may be omitted.
        
        
        
        
        
        
        
    \section{ident - extract an identifier from an argument}\label{sec:ipident}
        \index{variable!immediate!int}
        
        {\bf Call:} \\
        ident(v);
        
        {\bf Output parameters:} \\
        none
        
        {\bf Input parameters:} \\
        \begin{tabular}{| l | l | l | l |}
        \hline
        param  & mode & type & \\
        \hline \hline
        v & - &  -  & variable identifier   \\ \hline
        \end{tabular}
        
        {\bf Parameter=argument allowed:} \\
        No
        
        {\bf In-line target execution:} \\
        No
        
        {\bf Description:} \\
        {\bf ident()} can only be used to create an input parameter to a module,
        procedure or function. If it is called in any other context it generates a
        fatal error. {\bf ident()} creates a variable {\bf v} containing the
        identifier passed as a matching argument. The argument matching {\bf
        ident()} must be an identifier with no appended subcripts or fields. The
        identifier may have a scope prepended. It is not relevant if the identifier
        matches a previously created variable.
        
        The created variable {\bf v} is a struct with two fields; {\bf id}, which
        is the argument identifier as a string , and {\bf context} which is a
        string "global", "file",  "calling", "local" or "default"
        describing the scope context of the argument identifier.
        
        This procedure is intended to enable the writing of procedures
        which create variables, e.g. the following procedure creates an
        initialised one-dimensional string array -
        \begin{lstlisting}[gobble=8, language=threepl]
           procedure
           strarray (ident(v), var(init, "[]str")) {
               int(dim, dimension(init));
               str(name);
               switch (v.context) {
               case "global":
                   name = "global."+v.id;
               case "file":
                   name = "file."+v.id;
               case "calling":
                   trace();
                   fatal("strarray() - cannot have 'calling' context");
               case "local":
                   trace();
                   fatal("strarray() - cannot have 'local' context");
               case "default":
                   name = "calling."+v.id; // create in calling scope
               }
               var(name->, "["+dim+"]str", init);
           }
        \end{lstlisting}
        
        This procedure first gets the dimension of the required array from the
        dimension of the second argument, the array initialisation compound
        value. It then creates a string variable and this is assigned the
        required identifier string, {\bf v.id}, with a scope keyword prepended.
        The {\bf switch} makes sure that the scope requested, {\bf v.context},
        is sensible. {\bf "global"} and {\bf "file"} are OK and the default will
        create the variable in the scope of the caller. Scope {\bf "calling"}
        would imply the scope prior to the calling scope which is not
        accessible. Scope {\bf "local"} would be useful in case this procedure
        is called within a nested code block, but unfortunately that scope
        cannot be accessed. Finally a call to {\bf var()} creates the required
        variable by using the indirect operator on local string variable {\bf
        name}. The second argument to {\bf var()} is the type (string array) and
        the third argument supplies the initial values of the array.
        
        If called by 
        \begin{lstlisting}[gobble=8, language=threepl]
           strarray (sar, {"s1", "s2"});
        \end{lstlisting}
        it would create a variable {\bf sar} which is an array of
        "[2]str" initialised with entries "s1" and "s2". It is created
        in the scope of the call. Alternatively
        \begin{lstlisting}[gobble=8, language=threepl]
           strarray (global.sar, {"s1", "s2"});
        \end{lstlisting}
        would create {\bf sar} in global scope and
        \begin{lstlisting}[gobble=8, language=threepl]
           strarray (file.sar, {"s1", "s2"});
        \end{lstlisting}
        would create {\bf sar} in the current file module scope.
        
        Unfortunately as discussed above we cannot use scope {\bf "local"},
        e.g.
        \begin{lstlisting}[gobble=8, language=threepl]
           if (...) {
               ...
               strarray (local.sar, {"s1", "s2"});
           }
        \end{lstlisting}
        because we already need scope "calling" in the procedure to
        create the variable in the calling scope. The call
        \begin{lstlisting}[gobble=8, language=threepl]
           if (...) {
               ...
               strarray (sar, {"s1", "s2"});
           }
        \end{lstlisting}
        will create the variable in the scope of the code block of the
        {\bf if}.\footnote{3PL could be changed to allow an extra step back up
        the scope chain beyond "calling" ("callingcalling"?) but that might be
        regarded as excessive!}


% Not sure why this was ever written!!!!
% By the time this procedure is executed all source code has been read
% and parsed and any included files will have already been read.
%
%    \section{inludedir - add directory for source/module path}\label{sec:ipincludedir}
%        
%        {\bf Call:} \\
%        inludedir(dir);
%        
%        {\bf Output parameters:} \\
%        none
%        
%        {\bf Parameter=argument allowed:} \\
%        No
%        
%        {\bf Input parameters:} \\
%        \begin{tabular}{| l | l | l | l |}
%        \hline
%        param  & mode     & type  & \\
%        \hline \hline
%        dir & I & str & directory string    \\ \hline
%        \end{tabular}
%        
%        {\bf Parameter=argument allowed:} \\
%        No
%        
%        {\bf In-line target execution:} \\
%        No
%        
%        {\bf Description:} \\
%        The string argument is added to the list of directories to be searched for
%        source/module files. If the string has a leading "/" the directory path
%        will be absolute, otherwise it will be relative to the directory containing
%        the initial source file. A trailing "/" on the string is optional.



    \section{input - create an input mode variable}\label{sec:ipinput}
        \index{variable!input}
        
        {\bf Call:} \\
        input(ident, type, ....);
        
        input({ident, ident2}, type, ....);
        
        {\bf Output parameters:} \\
        none
        
        {\bf Input parameters:} \\
        \begin{tabular}{| l | l | l | l |}
        \hline
        param  & mode     & type  & \\
        \hline \hline
        ident   &     & any          & variable identifier  \\ \hline
        ident2  &     & any          & variable identifier  \\ \hline
        type    & I   & str or type  & type or type string  \\ \hline
        ....    &     &              & attributes           \\ \hline
        \end{tabular}
        
        {\bf Parameter=argument allowed:} \\
        No
        
        {\bf In-line target execution:} \\
        No
        
        {\bf Description:} \\
        This procedure creates an input mode variable.

        An input variable represents a connection from off-chip via input pads
        and buffers.

        The created variable has no clock domain. When an input variable is
        referenced in an assignment its clock domain will be established. The
        clock domain may also be established by assigning attribute {\bf
        readdomain} (see \autorefph{sec:attributes}) in which case references to
        the variable must be in that clock domain (unless it is an argument to
        inbuilt function {\bf accept()}).
        
        For bi-directional pads both {\bf input()} and {\bf output()} are
        used with the same pins.
       
        The resulting input mode variable cannot be assigned.
        
        {\bf ident} is the input mode variable identifier. The variable
        identifier can also be constructed from a string by using the indirect
        operator on a string variable or parenthesised expression, see
        subsection \autorefph{subsec:strvarcreat}.
        
        {\bf ident2} is an optional identifier for a differential output. This
        is only applicable for devices which support the IBUFDS\_DIFF\_OUT
        element. In that case {\bf ident} is the positive output and {\bf
        ident2} provides an inverted output. {\bf ident2} is created as a
        value mode variable.
        
        {\bf type} is the variable type which can be specified by a type-mode
        variable or by a string-valued expression. Any type is allowed.
        
        Any following arguments specify attributes. Each argument may specify an
        attribute by keymatch format, i.e. attribute=value where 'attribute' is
        an attribute identifier and 'value' is an appropriate mode and type
        (modes are all immediate except for clock domains which are mode clock).
        Alternatively an argument can be a map containing attributes where the
        key (string) is the attribute identifier and the associated value is
        the attribute value.

        Attributes may also be associated with the input pins. This is described
        in chapter \autorefph{chap:cpa} {\bf Constraints, Properties and
        Attributes} where other aspects of processing attributes and constraints
        for input, output and clock mode variables are also described.

        Attributes can be added by using inbuilt procedure {\bf attributes()}
        \autorefph{sec:ipattributes}.             
        
        Attributes are required for two main purposes; control of netlist
        generation and constraint file entries. Attributes giving input pin
        locations and input buffer types directly control netlist generation.
        I/O standards and timing group names exercise control of place-and-route
        via a constraint file. A clock domain specification is an exception and
        is used to specify the clock domain of the input should that be necessary.
        
        The applicable attributes are shown in the following table.

                These attributes are not case-sensitive but will be converted to
        lower case.

        The following attributes are used by {\bf input()} itself.
        
        \btabularx
        \hline
        \multicolumn{4}{|l|}{\bf input mode attributes} \\
        \hline
        differential & netlist & log          & differential                         \\
        loc          & netlist & str or []str & locations                            \\
        \hline
        \btabularxend
                
        Where an attribute value is an array the dimension must be equal to
        the number of pins and each array member will apply to the corresponding
        output pad or buffer. A single attribute value will be applied to all pins.

        {\bf differential} indicates that the input pads are differential.

        The {\bf loc} attribute is mandatory. The array dimension
        must exactly match the number of bits in the variable type. If the input
        is differential the array dimension must be twice the number of
        bits, the locations being used in pairs for the +ve and the -ve inputs
        (in that order). The location array matches the associated signal
        in order LSB to MSB for the array members and the signal bits.
            
        Pin strings are matched left to right to an input variable
        starting at the the least-significant bit. For a compound type pins are
        matched left to right from inner to outer nested types. For a struct the
        strings are matched left to right in field order.

        FPGA package pins are usually already defined in a platform header file,
        for example -
        \begin{lstlisting}[gobble=12, language=threepl]
           iopins(a_pins, "[12]str",
                        {
                            "A0", "B0", "A1", "B1", "A2", "C2",
                            "A3", "B3", "A4", "B4", "C3", "C4"
                        }
                );
        \end{lstlisting}
        There might be additional attributes included for these pins.

        \begin{lstlisting}[gobble=12, language=threepl]
           input(v, "[2][2]uint:2", loc=a_pins[2..9]);
        \end{lstlisting}

        The resulting input connections would be -\\
        pins "A1", "B1" (LSB, MSB) are connected to v[0][0] \\
        pins "A2", "C2" (LSB, MSB) are connected to v[0][1] \\
        pins "A3", "B3" (LSB, MSB) are connected to v[1][0] \\
        pins "A4", "B4" (LSB, MSB) are connected to v[1][1] \\

        The following attribute is read-only and is created when {\bf
        input()} is called and can be retrieved from the input mode
        variable if required by using inbuilt function {\bf attribute()}.
        
        \btabularx
        \hline
        \multicolumn{4}{|l|}{\bf input mode} \\
        \hline
        ids          & 3PL & []str & element IDs       \\
        \hline
        \btabularxend       
        
        {\bf ids} is a string array that contains the netlist identifiers of the
        input buffer elements. This attribute is read-only and is initialised
        when the input variable is created. This identifier can be used should
        there be a necessity to apply a constraint (other than location which is
        done via the port name) to the input buffer.

        The following attributes are required for constraint generation and are
        handled at completion of immediate execution by library procedure {\bf
        generateconstraints()} \autorefph{sec:lpgenerateconstraints}. The
        attribute values are single values which will be applied to all input
        bits.
        
        \btabularx
        \hline
        \multicolumn{4}{|l|}{\bf input mode} \\
        \hline
        ibuf\_delay\_value & constraint & int or []int & extra delay to IBUF                      \\
        ifd\_delay\_value  & constraint & int or []int & extra delay to IFD                       \\
        tnm                & constraint & str or []str & timing group id                          \\
        pullup             & constraint & log or []log & pullup                                   \\
        pulldown           & constraint & log or []log & pulldown                                 \\
        diff\_term         & constraint & log or []log & differential  $100\Omega$ termination    \\
        iostandard         & constraint & str or []str & voltage standard                         \\
        \hline
        \btabularxend
        
        {\bf ibufdelay} specifies a delay to be inserted in the non-registered
        output path from an input buffer. The values (integers) are 0 to 12
        inclusive (default 0).
        
        {\bf ifddelay} specifies a delay to be inserted in the registered
        path within an input buffer. The values (integers) are 0 to 6
        inclusive (default is a value generated automatically, see device
        databook).
        
        {\bf tnm} is the name of a timing group in which the input pads
        are to be included for timing constraint purposes.
        
        If {\bf pullup} or {\bf pulldown} are true the input will be given
        a weak pullup or pulldown. It is not allowed for both these
        attributes to be true.
        
        {\bf diff\_term} indicates that a $100\Omega$ termination resistance
        be applied to a differential input.
        
        {\bf iostandard} is a string specifying one of the allowed
        I/O standards (see Xilinx documentation).
        
        The following attributes are tolerated but ignored. They may be included
        as attributes for the output side of a bi-directional pin.
        
        \btabularx
        \hline
        \multicolumn{4}{|l|}{\bf input mode} \\
        \hline
        slew              & constraint & str or []str   & timing group id                          \\
        drive             & constraint & uint or []uint & drive current                                  \\
        \hline
        \btabularxend
        
        Any attributes that are not listed above will result in a fatal error.

        
        Note: for each pin a netlist port "PORT\_pin" is created where 'pin' is
        the pin identifying string. This port name is also the netlist identifier
        of the netlist from the output of the pad to the input of the IBUF or
        IBUFG. This convention can be used when generating associated constraints.
        
        In the Xilinx implementation if a static variable is assigned from an input
        variable, and has no other assignments (except for assignment of a constant
        equal to its initial value, which is implemented as a reset), then the
        associated flip-flops will be located in IOBs. This means that the
        propagation delay from pad to flip-flop is the minimum achievable (also
        controlled by attribute {\bf ifddelay}). If an input variable is used in an
        expression rather than being assigned directly to a static variable then a
        timing constraint from the input pads to the clock domain will usually need
        to be applied.
        
        {\bf input()} may be used as a parameter in a module, procedure or
        function definition. It must be passed an argument which is mode INPUT.
        The call to {\bf input()} to create the parameter must have only a
        single argument, the parameter identifier and all references to the
        parameter will apply to the argument variable. That variable will point
        indirectly to the argument INPUT mode variable. Although it is
        conventional to utlilise this for an input parameter, using {\bf input()}
        for an output parameter/argument will nevertheless function correctly.
        

        \subsection{input() as a parameter}
        
            {\bf input()} may be used as a parameter in a module, procedure or
            function definition. It must be passed an argument which is mode INPUT.
            
            The call to {\bf input()} to create the parameter need not have an
            identifier (1st argument), but in that case it cannot be referred
            to in the module, procedure or function body.

            The type argument (2nd argument) is optional. If provided it must match
            the type of the argument exactly. If not provided the parameter type
            will be determined from the argument variable type.
        
            When {\bf input()} is used to create a module, procedure or function
            parameter it cannot be passed attributes. The created parameter
            variable will adopt the attributes of the associated argument.
            
            All operations on the created parameter variable will be applied instead
            to its matching argument variable. If there is no argument to the input
            mode parameter then it will be  created as a local variable.

            Although it is conventional to utilise this for an input parameter,
            using {\bf input()} for an output parameter/argument will nevertheless
            function correctly.

            {\bf Note that the changes to attributes in the argument variable will
            take place even if the parameter input mode variable is not used!}






    \section{int - create one or more immediate integer variables}\label{sec:ipint}
        \index{variable!immediate!int}
        
        {\bf Call:} \\
        int(ident, init);
        
        {\bf Output parameters:} \\
        none
        
        {\bf Input parameters:} \\
        \begin{tabular}{| l | l | l | l |}
        \hline
        param  & mode & type & \\
        \hline \hline
        ident or \verb'{'id1, id2... \verb'}' & - &  -  & variable identifier(s)   \\ \hline
        init                   & I & int & optional initial value   \\ \hline
        \end{tabular}
        
        {\bf Parameter=argument allowed:} \\
        No
        
        {\bf In-line target execution:} \\
        No
        
        {\bf Description:} \\
        Create one or more immediate mode variables of type INT.
        
        {\bf ident} is the variable identifier. If more than one
        variable is to be created the list of variables must be in braces
        {\em id1, id2,\{\dots\}}. The variable identifier can also be constructed
        from a string by using the indirect operator on a string variable or
        parenthesised expression, see subsection \autorefph{subsec:strvarcreat}.
        
        An integer expression may be provided as the second
        argument to initialise the variables, otherwise they will be
        initially zero. Integers are 64 bits so the initialisation value
        is restricted to a range appropriate to that size.



    \section{liststacks - list the current state of the body and scope stacks}\label{sec:ipliststt}
        
        {\bf Call:} \\
        liststacks();
        
        {\bf Output parameters:} \\
        none
        
        {\bf Input parameters:} \\
        none
        
        {\bf Parameter=argument allowed:} \\
        No
        
        {\bf In-line target execution:} \\
        No
        
        {\bf Description:} \\
        List the state of the body and scope stacks on standard output. {\bf This
        procedure is intended only as an aid to 3PL interpreter development and has no
        other use.}
        


    \section{log - create one or more immediate logical variables}\label{sec:iplog}
        \index{variable!immediate!log}
        
        {\bf Call:} \\
        log(ident, init);
        
        {\bf Output parameters:} \\
        none
        
        {\bf Input parameters:} \\
        \begin{tabular}{| l | l | l | l |}
        \hline
        param  & mode & type & \\
        \hline \hline
        ident or \verb'{'id1, id2... \verb'}' & - &  -  & variable identifier(s)   \\ \hline
        init                   & I & log & optional initial value   \\ \hline
        \end{tabular}
        
        {\bf Parameter=argument allowed:} \\
        No
        
        {\bf In-line target execution:} \\
        No
        
        {\bf Description:} \\
        Create one or more immediate mode variables of type LOG.

        {\bf ident} is the variable identifier. If more than one
        variable is to be created the list of variables must be in braces
        {\em id1, id2,\{\dots\}}. The variable identifier can also be constructed
        from a string by using the indirect operator on a string variable or
        parenthesised expression, see subsection \autorefph{subsec:strvarcreat}.
        
        A logical expression may be provided as the second
        argument to initialise the variables, otherwise they will be
        initially false.


    \section{msg - print to standard out and the report file}\label{sec:ipmsg}
        
        {\bf Call:} \\
        msg(s, t);
        
        {\bf Output parameters:} \\
        none
        
        {\bf Input parameters:} \\
        \begin{tabular}{| l | l | l | l |}
        \hline
        param & mode & type & \\
        \hline \hline
        s & I & str & message string    \\ \hline
        t & I & log & print a traceback \\ \hline
        \end{tabular}
        
        {\bf Parameter=argument allowed:} \\
        No
        
        {\bf In-line target execution:} \\
        No
        
        {\bf Description:} \\
        Print a message to standard output and to the optional report file if
        it is active.
        
        {\bf s} is the message string.
        
        {\bf t} (optional) if true requests a stack trace.
        
        See also - \\
        {\bf println()} \autorefph{sec:ipprintln} \\
        {\bf texit()} \autorefph{sec:iptexit} \\
        {\bf rpt()} \autorefph{sec:iprpt}.


    \section{nop - a no-operation (delay)}\label{sec:ipnop}
        
        {\bf Call:} \\
        nop(n);
        
        {\bf Output parameters:} \\
        none
        
        {\bf Input parameters:} \\
        \begin{tabular}{| l | l | l | l |}
        \hline
        param & mode & type & \\
        \hline \hline
        n & I & int & number of cycles  \\ \hline
        \end{tabular}
        
        {\bf Parameter=argument allowed:} \\
        No
        
        {\bf In-line target execution:} \\
        Yes
        
        {\bf Description:} \\
        Do nothing for one or more cycles.
        
        {\bf n} is the number of cycles to delay. If omitted, the default is 1.


    \section{negclock - create a negative clock}\label{sec:ipnegclock}
        
        {\bf Call:} \\
        negclock(ident, c);
        
        {\bf Output parameters:} \\
        none
        
        {\bf Input parameters:} \\
        \begin{tabular}{| l | l | l | l |}
        \hline
        param & mode & type & \\
        \hline \hline
        ident & - & - & identifier of new -ve clock  \\ \hline
        c     & C & - & reference clock variable  \\ \hline
        \end{tabular}
        
        {\bf Parameter=argument allowed:} \\
        No
        
        {\bf In-line target execution:} \\
        Yes
        
        {\bf Description:} \\
        Creates a negative-edge clock where {\bf ident} is the
        identifier of the new clock. The clock is a negative-edge
        version of the the clock {\bf c} (which must not itself
        be a negative clock). The new clock will have the frequency
        and period parameters of the reference clock and will have
        1.0 minus its duty cycle. If the frequency, period or duty cycle
        of either clock is changed the parameters of the associated clock will
        also change as appropriate.
        
        A negative clock is not actually generated. Instead the created
        negative clock variable is linked to the related positive clock variable
        but is flagged as inverted and the signal identifier in the intermediate
        code list has an "\_" appended. When a negative clock signal is
        encountered during generation of the EDIF netlist the positive clock
        signal is used followed by an inverter. Xilinx clocked elements
        have an inbuilt inverter so that is used rather than an explicit
        external inverter element. This avoids the use of a separate
        negative clock network.



    \section{open - open a file}\label{sec:ipopen}
        \index{variable!input}
        
        {\bf Call:} \\
        open(ident, rwa);
        
        {\bf Output parameters:} \\
        none
        
        {\bf Input parameters:} \\
        \begin{tabular}{| l | l | l | l |}
        \hline
        param  & mode     & type  & \\
        \hline \hline
        var &     & any & variable identifier               \\ \hline
        rwa & I   & str & read, write or append specifier   \\ \hline
        \end{tabular}
        
        {\bf Parameter=argument allowed:} \\
        No
        
        {\bf In-line target execution:} \\
        No
        
        {\bf Description:} \\
        This procedure opens a file for reading or writing.
        The first argument is the file immediate variable identifier.
        The second argument is "r" for reading, "w" for writing or
        "a" for appending.



    \section{output - create an output mode variable}\label{sec:ipoutput}
        \index{variable!input}
        
        {\bf Call:} \\
        output(ident, type, src, ....);
        
        {\bf Output parameters:} \\
        none
        
        {\bf Input parameters:} \\
        \begin{tabular}{| l | l | l | l |}
        \hline
        param  & mode     & type  & \\
        \hline \hline
        var     &     & any          & optional variable identifier  \\ \hline
        type    & I   & str or type  & optional type or type string  \\ \hline
        src     & T   & str          & optional source               \\ \hline
        ....    &     &              & attributes                    \\ \hline
        \end{tabular}
        
        {\bf Parameter=argument allowed:} \\
        No
        
        {\bf In-line target execution:} \\
        No
        
        {\bf Description:} \\
        This procedure creates an output mode variable.

        An output mode variable has no initial clock domain. The clock domain is
        determined by the first assignment. Any subsequent assignments, whether to
        the whole variable or a component of a compound type, must have the same
        clock domain. The clock domain may be established by assigning attribute
        {\bf writedomain} (see attributes below). If the value to be
        assigned has a different clock domain it must be converted by one of the
        inbuilt functions {\bf sample()} \autorefph{sec:ifsample} or {\bf accept()}
        \autorefph{sec:ifaccept}.
        
        The resulting output mode variable cannot be evaluated, i.e. it cannot
        appear on the RHS of an assignment statement, or anywhere else a value
        is expected.
        
        For bi-directional pads both {\bf input()} and {\bf output()} are
        used with the same pins (locations).
        
        {\bf ident} is the variable identifier. If parameter {\bf src} is matched
        then this argument is optional since there may be no reason to refer to the
        variable anywhere. The variable identifier can also be constructed from a
        string by using the indirect operator on a string variable or parenthesised
        expression, see subsection \autorefph{subsec:strvarcreat}.

        {\bf type} is the variable type which can be specified by a variable of
        type "type" or by a string-valued expression. Any type is allowed except
        "map" or "list". If parameter {\bf src} is matched then this argument is
        optional since the type can be extracted from the third argument, but if
        the third argument is immediate mode the type argument must be provided
        unless the third argument is type "log" (a target type "log" can be
        unambiguously inferred from immediate type "log" whereas other types
        require a width to be specified for the target type).
        
        {\bf src} is an optional data source for the output. Its use avoids the
        need for assignments to the output mode variable to establish the data
        source (see \ref{sec:outputiass}). This can only be used for
        unconditional output, i.e where 3-state output is not required. The
        argument may be immediate mode. If later assignment to the output
        variable is made this source is overridden. If this argument is not used
        an explicit assignment must be made to this output variable. If this
        argument is provided the type argument (second argument) can be omitted
        if the source is a target mode or if it is immediate mode type "log". If
        this argument is immediate mode and any type except "log" then the type
        argument must also be provided. The argument to parameter {\bf src}
        cannot be queue, cmemory, rmemory or output mode.
        
        {\bf src} may also be a 1-dimensional list of variables and constants
        of type "log" in braces.
        
        Note: if {\bf src} contains value mode variables which have not been
        assigned values yet, the interim values will be set to the current clock
        domain. When ultimately these value variables are assigned values the
        values must have the same clock domain. Other than the above specific
        case no attempt is made by {\bf output()} to establish the output clock
        domain of argument variables to parameter {\bf src}.      
        
        Any following arguments specify attributes. Each argument may specify an
        attribute by keymatch format, i.e. attribute=value where 'attribute' is
        an attribute identifier and 'value' is an appropriate mode and type
        (modes are all immediate except for clock domains which are mode clock).
        Alternatively an argument can be a map containing attributes where the
        key (string) is the attribute identifier and the associated value is
        the attribute value.

        Attributes may also be associated with the output pins. This is described
        in chapter \autorefph{chap:cpa} {\bf Constraints, Properties and
        Attributes} where other aspects of processing attributes and constraints
        for input, output and clock mode variables are also described.

        Attributes can be added by using inbuilt procedure {\bf attributes()}
        \autorefph{sec:ipattributes}.             
        
        Attributes are required for two main purposes; control of netlist
        generation and constraint file entries. Attributes giving input pin
        locations and output buffer types directly control netlist generation.
        I/O standards and timing group names exercise control of place-and-route
        via a constraint file. A clock domain specification is an exception and
        is used to specify the clock domain of the output should that be necessary.
        
        When {\bf output()} is used to create a module, procedure or function
        parameter it cannot be passed attributes. The created parameter variable
        will adopt the attributes of the associated argument.
        
        The applicable attributes are shown in the following table.

                These attributes are not case-sensitive but will be converted to
        lower case.
         
        The following six attributes are used by {\bf output()} itself.
        
        \btabularx
        \hline
        \multicolumn{4}{|l|}{\bf output mode} \\
        \hline
        direct       & 3PL        & log          & enable direct from test               \\
        differential & netlist    & log or []log & differential                          \\
        defaultout   & 3PL        & log or []log & default for unconnected outputs       \\
        noshadow     & 3PL        & log or []log & don't duplicate output flip-flop      \\
        loc          & constraint & str or []str & locations                             \\
        \hline
        \btabularxend
        
        Where an attribute value is an array the dimension must be equal to
        the number of pins and each array member will apply to the corresponding
        output pad or buffer. A single attribute value will be applied to all pins.
        
        {\bf direct} must have been set prior to any assignment statements using the
        output mode variable. The other attributes are not used until completion of
        immediate execution (excluding postprocessing).

        {\bf direct} specifies that if possible the output 3-state enable will
        be driven directly from a conditional signal, not through an execution
        thread and WHEN block. This can significantly simplify the logic and
        reduce the probability of failing to meet timing constraints locally.
        Note that when this attribute is in effect the 3-state enable is
        immediately active following configuration loading regardless of any 3PL
        startup mechanisms - see {\bf
        start()}(\autorefph{sec:ipstart}).
        
        {\bf direct} cannot be passed via iopins() as it cannot be expanded to
        an array, being a single attribute applying to the whole variable and not
        to individual outputs.
        
        {\bf differential} indicates that the output pads are differential.
        Two locations must be supplied for each data bit (see {\bf loc} above).
        
        {\bf defaultout} specifies that unconnected outputs variables, or parts of
        output variables, should be tied high or low. If this attribute is not
        supplied then the action on encountering unconnected outputs depends on the
        global directive "unass\_out" (see \autorefph{sec:directives}). The values
        must be {\bf true} to tie the outputs high, or {\bf false} to tie low.
        
        {\bf noshadow} when true inhibits the duplication of flip-flops
        whose output is connected to an output buffer. By default any
        flip-flop whose output is connected to an output buffer is placed
        in the IO pad containing the buffer. If the output of the flip-flop
        is used elsewhere then the flip-flop is duplicated in a CLB with
        all inputs in parallel with the pad flip-flop and the output of the
        duplicate used for internal logic. If {\bf noshadow} is true then
        duplication will not be done, the flip-flop will of necessity
        reside in a CLB and its output will be routed to the output buffer.
        This will result in much worse clock-to-output time. Setting this
        attribute ensures that the pad output genuinely reflects the
        flip-flop state and not that of a duplicate but at the cost of
        delayed output.
       
        The {\bf loc} attribute is mandatory. The array dimension
        must exactly match the number of bits in the variable type. If the output
        is differential the array dimension must be twice the number of
        bits, the locations being used in pairs for the +ve and the -ve inputs
        (in that order). The location array matches the associated signal
        in order LSB to MSB for the array members and the signal bits.
        in order LSB to MSB for the array members and the signal bits.
            
        Pin strings are matched left to right to an output variable
        starting at the the least-significant bit. For a compound type pins are
        matched left to right from inner to outer nested types. For a struct the
        strings are matched left to right in field order.

        FPGA package pins are usually already defined in a platform header file,
        for example -
        \begin{lstlisting}[gobble=12, language=threepl]
            static(v, "[2][2]uint:2");
            iopins(a_pins, "[12]str",
                        {
                            "A0", "B0", "A1", "B1", "A2", "C2",
                            "A3", "B3", "A4", "B4", "C3", "C4"
                        }
                );
        \end{lstlisting}
        There might be additional attributes included for these pins.

        \begin{lstlisting}[gobble=12, language=threepl]
           output(,, v, "[2][2]uint:2", loc=a_pins[2..9]);
        \end{lstlisting}

        The resulting output connections would be -\\
        v[0][0] (LSB, MSB) is connected to pins "A1", "B1" \\
        v[0][1] (LSB, MSB) is connected to pins "A2", "C2" \\
        v[1][0] (LSB, MSB) is connected to pins "A3", "B3" \\
        v[1][1] (LSB, MSB) is connected to pins "A4", "B4" \\

        The following attribute is read-only and is created when {\bf
        output()} is called and can be retrieved from the output mode
        variable if required by using inbuilt function {\bf attribute()}.
        
        \btabularx
        \hline
        \multicolumn{4}{|l|}{\bf output mode} \\
        \hline
        ids          & 3PL & []str & element IDs       \\
        \hline
        \btabularxend
        
        {\bf ids} is a string array that contains the netlist identifiers of the
        output buffer elements. This attribute is read-only and is initialised
        when the output variable is created. This identifier can be used should
        there be a necessity to apply a constraint (other than location which is
        done via the port name) to the output buffer.

        The following attributes are required for constraint generation and are
        handled at completion of immediate execution by library procedure {\bf
        generateconstraints()} \autorefph{sec:lpgenerateconstraints}. The
        attribute values are single values which will be applied to all output
        bits.
        
        \btabularx
        \hline
        \multicolumn{3}{|l|}{\bf output mode attributes} \\
        \hline
        slew         & str or []str & slew rate         \\
        drive        & int or []int & drive current     \\
        tnm          & str or []str & timing group id   \\                  
        pullup       & log or []log & pullup            \\
        pulldown     & log or []log & pulldown          \\
        iostandard   & str or []str & voltage standard  \\
        \hline
        \btabularxend
       
        The {\bf slew} attribute is the string "fast" or "slow" and selects
        one of two slew rates for Xilinx output buffers (see Xilinx documentation).
        
        The {\bf drive} attribute is an integer which specifies the output drive
        current for Xilinx output buffers (see Xilinx documentation).
        
        {\bf tnm} is the name of a timing group in which the output pads
        are to be included for timing constraint purposes.
        
        If {\bf pullup} or {\bf pulldown} are true the output will be given
        a weak pullup or pulldown. It is not allowed for both these
        attributes to be true.
        
        {\bf iostandard} is a string specifying one of the allowed
        I/O standards (see Xilinx documentation).
        
        The following attributes are tolerated but ignored. They may be included
        as attributes for the input side of a bi-directional pin.
        
        \btabularx
        \hline
        \multicolumn{4}{|l|}{\bf input mode} \\
        \hline
        ibuf\_delay\_value & constraint & int & extra delay to IBUF                      \\
        ifd\_delay\_value  & constraint & int & extra delay to IFD                       \\
        \hline
        \btabularxend

        Any attributes that are not listed above will result in a fatal error.


        There are two forms of assignment to output variables -
        
        Unconditional assignment is used where 3-state output is not required. This
        has the form
        \begin{lstlisting}[gobble=8, language=threepl]
           output(ov, type,, attr);
           ov := value;
        \end{lstlisting}
        If the type is compound then separate assignments can be made to
        different components, i.e. array members or struct fields. Output
        variables can have multiple assignments but the last assignment
        prior to termination of immediate execution is the one which
        will take effect. Any prior assignments have no effect.
        
        For unconditional assignment it is usually more succinct to include
        the source as the third argument to output() rather than performing a
        separate explicit assignment, e.g. the case above could instead be
        written -
        \begin{lstlisting}[gobble=8, language=threepl]
           output(,, value, attr);
        \end{lstlisting}
        Here the first argument is omitted since the identifier is not needed as
        the variable is not explicitly assigned. The second argument is omitted
        because the source is a target mode from which the type can be inferred.
        The third argument can also be a list of values of type "log", e.g.
        \begin{lstlisting}[gobble=8, language=threepl]
           static(l, "log");
           output(,, {false, l, true}, loc=pins[3..5]);
        \end{lstlisting}
        where the list values can be immediate or allowed target modes,
        constants or variables, as long as they are type "log".
        
        Conditional assignment is used to exercise 3-state control over
        the output. This has the form
        \begin{lstlisting}[gobble=8, language=threepl]
           output(ov, type,, attr);
           when (enable)
               ov |- value;
        \end{lstlisting}
        Conditional assignment can only be done once\footnote{This restriction
        could be relaxed if warranted} and it
        cannot follow an unconditional assignment. Similarly an
        unconditional assignment cannot follow a conditional assignment.
        A conditional assignment must assign the whole variable, not
        just some components of a compound type. Thus an output mode
        variable appearing in an unconditional assignment cannot have
        subscripts or struct fields. If fields or array members are to be
        assigned from different sources a value variable of the same type
        as the output mode variable can be used to collect up the disparate
        sources into a single variable for assignment.

        
        Note: for each pin a netlist port "PORT\_pin" is created where 'pin' is
        the pin identifying string. This port name is also the netlist identifier
        of the netlist from the output of the OBUF to the input of the pad.
        This convention can be used when generating associated constraints.
        
        Note also that where an output mode variable does not have an identifier
        (the first argument is omitted) an identifier of the form OUPUTnnn is
        created for the variable where nnn is just an integer in sequence from 0.

        \subsection{output() as a parameter}
        
            {\bf output()} may be used as a parameter in a module, procedure or
            function definition. It must be passed an argument which is mode OUTPUT.
            
            The call to {\bf output()} to create the parameter need not have an
            identifier (1st argument), but in that case it cannot be referred
            to in the module, procedure or function body.

            The type argument (2nd argument) is optional. If provided it must match
            the type of the argument exactly. If not provided the parameter type
            will be determined from the argument variable type. If there is no
            argument the type can be inferred from the source argument type,
            described below.

            The source argument (3rd argument) if supplied will overwrite any source
            in the argument variable. The type must match that of the variable
            as described above, or it may determine the variable type if it cannot
            be inferred otherwise.

            Attributes may be supplied as the 4th and following arguments. These
            attributes will be added to or supplant those of the argument output mode
            variable.
            
            All operations on the created parameter variable will be applied instead
            to its matching argument variable. If there is no argument to the output
            mode parameter then it will be  created as a local variable.

            Although it is conventional to utilise this for an output parameter,
            using {\bf output()} for an input parameter/argument will nevertheless
            function correctly.

            {\bf Note that the changes to a source value or to attributes in the
            argument variable will take place even if the parameter output mode
            variable is not used!}



    \section{part - specify the FPGA}\label{sec:ippart}

        {\bf Call:} \\
        part(p);

        {\bf Output parameters:} \\
        none

        {\bf Input parameters:} \\            
        {\bf p} is the part specification (string).
        
        {\bf Parameter=argument allowed:} \\
        No

        {\bf In-line target execution:} \\
        no

        {\bf description:} \\
        This procedure notifies the compiler of the FPGA part specification.
        The procedure creates directive "part", which contains the
        part specification string, and also creates a directive "family"
        which has the values "XC2S", "XCV", "XC3S", "XC2V", "XC2V" etc as
        appropriate.

        The part specification must be as expected by the vendor software,
        e.g. "xcv300ebg432-7".




    \section{printf - print a formatted line to stdout}\label{sec:ipprintf}

        {\bf Call:} \\
        printf(fmt, v1, v2,....);
        
        {\bf Output parameters:} \\
        none
        
        {\bf Input parameters:} \\
        \begin{tabular}{| l | l | l | l |}
        \hline
        param  & mode     & type  & \\
        \hline \hline
        fmt & I & str & format string   \\ \hline
        v1  & I & any & input variable  \\ \hline
        v2  & I & any & input variable  \\ \hline
        .   & I & any & input variable  \\ \hline
        .   & I & any & input variable  \\ \hline
        \end{tabular}
        
        {\bf Parameter=argument allowed:} \\
        No
        
        {\bf In-line target execution:} \\
        No
        
        {\bf Description:} \\
        This procedure is equivalent to printf() in C. The format string is that used
        by the Java Formatter class.
        
        The format argument may be followed by a single argument which is an array of
        pointers, for example {\bf inargs} in the calling module, procedure or function.
        In that case the pointers are resolved and provide the input variable list to be
        printed.
             
        See the note on character and byte reads and writes in the description
        of inbuilt procedure {\bf file()} \autorefph{sec:ipfile}.
   
        See also - \\
        {\bf fprintf()} \autorefph{sec:ipfprintf} \\
        {\bf sprintf()} \autorefph{sec:ipsprintf} \\
        {\bf writeln()} \autorefph{sec:ipwriteln} \\
        {\bf writebyte()} \autorefph{sec:ipwritebyte}.




    \section{presetdirective - enter a directive during parsing}\label{sec:ippresetdirective}
        
        {\bf Call:} \\
        presetdirective(key, value);
        
        {\bf Output parameters:} \\
        none
        
        {\bf Input parameters:} \\
        \begin{tabular}{| l | l | l | l |}
        \hline
        param  & mode     & type  & \\
        \hline \hline
        key   & I & str & directive key     \\ \hline
        value & I & any & directive value (optional)   \\ \hline
        \end{tabular}
        
        {\bf Parameter=argument allowed:} \\
        No
        
        {\bf In-line target execution:} \\
        No
        
        {\bf Description:} \\
        {\bf Note: this procedure is executed when encountered during parsing.
        It is not an inbuilt procedure that can be be called during immediate
        execution.}
        
        This procedure creates a directive.
        
        It is intended to facilitate conditional file inclusion, see
        \autorefph{sec:fileinclusion}.
        
        {\bf The procedure call can only be preceded on the line by white space.
        The call syntax may include white space for formatting. The remainder
        of the line following the semicolon is ignored.} These limitations arise
        because the statement is not recognised by the parser but is
        intercepted by the source file reader, the line not being passed to
        the token generator and parser.
        
        The key is a string. The value may one of the 3PL immediate primitive
        types {\bf int}, {\bf log}, {\bf float} or {\bf str}. If the second
        argument is not provided the entry will be type {\bf log} and value {\bf
        true}.

        The key and value are expressions which cannot contain variables or
        calls to user-defined functions. They can contain constants, calls to
        inbuilt functions, evaluations of global directives via inbuilt function
        {\bf directive()} (See \autorefph{sec:ifdirective}) and expression
        operators. Binary integer constants of the form 0b.... or 0B.... are not
        allowed in this context, but hexadecimal and octal forms are accepted.
        
        Because this inbuilt procedure is recognised prior to the token
        generator and parser an erroneous call may shadow a prior error. For
        example
        \begin{lstlisting}[gobble=8, language=threepl]
           module "nofile"
           presetdirective("name", 3.x);
        \end{lstlisting}
        would result in a fatal error for the second statement even when the
        first statement fails to open file "nofile". The first error is passed
        to the parser as a token but the second error causes an exception
        in the file reader before the parser has processed the error token.
        
        see also inbuilt procedures \\
        {\bf directive()} \autorefph{sec:ipdirective} \\
        {\bf removedirective} \autorefph{sec:ipremovedirective}. \\
        and inbuilt functions \\
        {\bf directive()} \autorefph{sec:ifdirective}.\\
        {\bf direxists()} \autorefph{sec:ifdirexists}.\\
        {\bf directives()} \autorefph{sec:ifdirectives}. \\


    \section{println - print a line}\label{sec:ipprintln}
        
        {\bf Call:} \\
        println(v);
        
        {\bf Output parameters:} \\
        none
        
        {\bf Input parameters:} \\
        \begin{tabular}{| l | l | l | l |}
        \hline
        param & mode & type & \\
        \hline \hline
        v & I & int, log, str & value to be printed \\ \hline
        \end{tabular}
        
        {\bf Parameter=argument allowed:} \\
        No
        
        {\bf In-line target execution:} \\
        No
        
        {\bf Description:} \\
        Print a line to standard out. The argument may be types int, uint, log,
        str or an enumerated type (the value identifier is printed).
        
        See also {\bf msg()} \autorefph{sec:ipmsg}.
        
        See the note on character and byte reads and writes in the description
        of inbuilt procedure {\bf file()} \autorefph{sec:ipfile}.


    \section{priority - create a priority variable}\label{sec:ippriority}
        \index{variable!priority}
        
        {\bf Call:} \\
        priority(ident, type, init, ...); \\
        or  \\
        priority(ident, size, init, ...); \\
        
        
        {\bf Output parameters:} \\
        none
        
        {\bf Input parameters:} \\
        \begin{tabular}{| l | l | l | l |}
        \hline
        param  & mode & type & \\
        \hline \hline
        ident & - &  -        & variable identifier          \\ \hline
        size  & I & type uint & optional array size          \\
        type  & I & type str  & optional type or type string \\ \hline
        init  & I & log       & optional initialisation      \\ \hline
        ...   & I &  -        & optional attributes          \\ \hline
        \end{tabular}
        
        {\bf Parameter=argument allowed:} \\
        No
        
        {\bf In-line target execution:} \\
        No
        
        {\bf Description:} \\
        Create a priority variable. Its type will be an array of "log".
        
        {\bf size} is an optional 2nd argument which is the dimension of the
        "[]log" array. Note that if a priority mode variable is created by using
        the inbuilt procedure {\bf var()} \autorefph{sec:ipvar} and the array
        size is to be specified then a proper type
        argument must be used (see below).

        {\bf type} is a type variable or a string expression giving the variable
        type explicitly. The type must be an array of log or a fatal error will
        occur.
        
        The array may be larger than required in which case unused entries,
        whether trailing or interspersed, will be pruned.

        The size/type argument may be omitted and for a variable which is not a
        parameter it will be given the default type "[20]log".
        In the case of a parameter the variable will take the type of the
        argument.
        
        The optional third argument, {\bf init}, allows the inputs to the
        priority array to be assigned initially rather than by the use of
        explicit assignment statements later.
        
        Additional trailing arguments may be passed to set attributes. These
        must be of the form attr=value.
        
        A priority variable can be assigned or re-assigned at any time,
        the assigned variables only taking effect at the conclusion of
        immediate execution.
        
        A priority variable has a single clock domain for reading and writing.
        When a priority variable is created using the inbuilt procedure {\bf
        priority()} it has no associated clock domain. The clock domain is
        established by
        \blist
        \item   the first assignment to the variable
        \item   the first use of the value of the variable
        \blend
        
        {\bf Note that evaluation of the variable in any context, for example
        in inbuilt or library modules, procedures or functions, will establish
        a clock domain if the variable does not yet have a clock domain.}
        This will always be the current clock domain (unless the code has
        been explicitly written to establish the clock domain of its argument
        to something other than the current clock domain, in which case
        this will presumably be well documented).
        
        When a priority variable is assigned, the clock domain of the assigned
        value must be the same as the current clock domain or null.
        
        When a priority variable is assigned, If it has no clock domain the
        current clock domain will be established as the clock domain of the
        variable. If it has a clock domain it must be the same as the current
        clock domain.
        
        When the value of a priority variable is used that value will have the
        clock domain of the variable.
        
        Applicable attributes are shown in the following table. The
        attributes can be passed as arguments to this procedure or may be
        applied to the variable using the inbuilt procedure
        {\bf attributes()} \autorefph{sec:ipattributes}.
        
        When {\bf priority()} is used to create a module, procedure or function
        parameter it cannot be passed attributes. The created parameter variable
        will adopt the attributes of the associated argument.

                \btabularx
        \hline
        \multicolumn{4}{|l|}{\bf priority mode attributes} \\
        \hline
        domain     & 3PL & clock   & force the clock domain \\     
        readonly   & 3PL & boolean & set to read only       \\     
        multiplein & 3PL & boolean & OR multiple inputs       \\     
        \hline
        \btabularxend

        {\bf domain} forces the clock domain (both input and output) to
        the specified clock domain. If assigned {\bf null} the current
        clock domain is selected. If the variable already has a clock domain
        as a result of having been assigned or evaluated then that clock domain
        is compared with the intended domain. If the assigned clock domain
        agrees with that already in place then there is no further action,
        otherwise a fatal error occurs.
               
        {\bf readonly} true specifies that from now on this variable cannot
        be assigned to. This attribute cannot be assigned false.
        
        Normally an assignment to a member of a priority variable array
        overwrites any previous assignment, the final assignment prior to
        completion of immediate execution prevailing. If attribute {\bf
        multiplein} is set true then all assignments to a member of a priority
        variable array are collected and on completion of immediate execution
        these are ORed to form the input to that array member.






    \section{priunlock - unlock priority execution}\label{sec:ippriunlock}\index{variable!priority}
        
        {\bf Call:} \\
        priunlock(ident);
        
        {\bf Output parameters:} \\
        none
        
        {\bf Input parameters:} \\
        \begin{tabular}{| l | l | l | l |}
        \hline
        param  & mode & type & \\
        \hline \hline
        ident  & Q &  -       & priority variable identifier \\ \hline
        \end{tabular}
        
        {\bf Parameter=argument allowed:} \\
        No
        
        {\bf In-line target execution:} \\
        yes
        
        {\bf Description:} \\
        This procedure creates in-line code which unlocks a priority variable
        previously locked by a {\bf waitprilock} statement. It executes in
        one clock cycle.
        
        For a description of usage, along with control structures
        {\bf waitpri} and {\bf waitprilock}, see section \autorefph{sec:waitpri}.



    \section{pulse - generate a pulse - DEPRECATED!}\label{sec:ippulse}
        
        
        {\bf Call:} \\
        pulse(v \verb'<-' o);
        
        {\bf Output parameters:} \\
        \begin{tabular}{| l | l | l | l |}
        \hline
        param & mode & type & \\
        \hline \hline
        v & V & log & value carrying the pulse \\ \hline
        \end{tabular}
        
        {\bf Input parameters:} \\
        \begin{tabular}{| l | l | l | l |}
        \hline
        param & mode & type & \\
        \hline \hline
        o & I & log & if true, OR with previous value \\ \hline
        \end{tabular}
        
        {\bf Parameter=argument allowed:} \\
        No
        
        {\bf In-line target execution:} \\
        yes
        
        {\bf Description:} \\
        Pulse a logical value variable for one clock cycle when this statement
        executes. The value is true for the clock cycle leading up to the
        clock edge on which execution occurs.
        
        This procedure has been replaced by the inbuilt
        procedure {\bf assert()}(\autorefph{sec:ipassert}). The latter
        procedure is identical in function and the name was changed as it was
        felt that 'pulse' was too hardware-oriented. In addition this procedure
        may be executed on several successive clock cycles in which case it
        really could be viewed as a level rather than a pulse. The important
        thing is that it is asserted true in time for the next clock edge.
        
        If the input argument is provided and is true the output value variable
        will be assigned the OR of its previous value and the execution pulse.
        If it has no previous assigned value it is assigned the execution pulse.
        Using this optional input argument a signal pulsed from more than one
        place can be generated. {\bf Note: when execution pulses are ORed in
        this way, use of the value of the variable must follow the last call to
        pulse() since each call re-assigns the value variable!}
        
        Note that if the ORing option is not selected and multiple calls to
        {\bf pulse()} are made with the same output argument, each call will
        re-assign the output value variable, which may be what is wanted or
        may not!
        


    \section{put - put a map entry}\label{sec:ipput}
        
        {\bf Call:} \\
        put(map, key, value);
        
        {\bf Output parameters:} \\
        none
        
        {\bf Input parameters:} \\
        \begin{tabular}{| l | l | l | l |}
        \hline
        param  & mode & type & \\
        \hline \hline
        map    & I     & map             & a map        \\ \hline
        key    & I     & str, int, uint  & a string key \\ \hline
        value  & I V C & any             & the value    \\ \hline
        \end{tabular}
        
        {\bf Parameter=argument allowed:} \\
        Yes
        
        {\bf In-line target execution:} \\
        No
        
        {\bf Description:} \\
        The item {\bf value} is placed in the map {\bf map} against key
        {\bf key}. If that key is already in the map the value will
        be replaced. If the last argument is omitted null will be put
        against the key. The key should be type "str" but if it is
        type "int" or type "uint" it will be converted to type "str".


    \section{putall - put entries from another map}\label{sec:ipputall}
        
        {\bf Call:} \\
        putall(m, m2);
        
        {\bf Output parameters:} \\
        none
        
        {\bf Input parameters:} \\
        \begin{tabular}{| l | l | l | l |}
        \hline
        param  & mode & type & \\
        \hline \hline
        m      & I     & map  & a map   \\ \hline
        m2     & I     & map  & another map   \\ \hline
        \end{tabular}
        
        {\bf Parameter=argument allowed:} \\
        No
        
        {\bf In-line target execution:} \\
        No
        
        {\bf Description:} \\
        All items from map {\bf m2} are added to map {\bf m}. Where a
        key already exists its value will be replaced.



 

    \section{queue - create one or more queue variables}\label{sec:ipqueue}
        \index{variable!queue}
        
        {\bf Call:} \\
        queue(ident, type, init, ...);
        
        {\bf Output parameters:} \\
        none
        
        {\bf Input parameters:} \\
        \begin{tabular}{| l | l | l | l |}
        \hline
        param  & mode & type & \\
        \hline \hline
        ident or \verb'{'id1, id2... \verb'}' & - &  -       & variable identifier(s)       \\ \hline
        type                   & I & type str & optional type or type string \\ \hline
        init                   & I & any      & optional initialisation      \\ \hline
        ...                    & I &  -       & optional attributes          \\ \hline
        \end{tabular}
        
        {\bf Parameter=argument allowed:} \\
        No
        
        {\bf In-line target execution:} \\
        No
        
        {\bf Description:} \\
        Create one or more queue variables. If more than one variable is to be
        created the list of variables must be in braces (\{\dots\}).
        {\bf type} is a type variable or a string expression giving the
        variable type. This may be omitted in the case of parameter
        creation in which case the variable will take the type of the
        argument.
        
        The optional third argument allows the queue to be created with an
        initial datum, i.e. the queue will be initially available. This
        argument must be immediate mode. If the type is compound the
        initialisation argument will usually be a compound value, but it may 
        be initialised by a single immediate value in which case all members of
        the compound type will be initialised to that value.
        
        Additional trailing arguments may be passed to set attributes. These
        must be of the form attr=value.
        
        A queue variable has a clock domain for reading and another clock domain
        for writing. When a static variable is created using the inbuilt
        procedure {\bf queue()} it has no associated clock domains.
        
        The clock domain for writing is established by -
        \blist
        \item   the first write (assignment) to the queue variable (or
                component of the variable if the type is compound)
        \item   the use of the write availability operator \verb'>' on the queue
        \item   a call to {\bf queuewritestatus()} on the queue
        \item   a call to {\bf reset()} on the queue
        \blend
        If the write clock domain is not yet established in any of these
        situations then it will be set to the current clock domain,
        otherwise it must be the same as the current clock domain.

        The clock domain for reading is established by -
        \blist
        \item   the first read or examine of the variable (or component of
                the variable if the type is compound)
        \item   the use of the read availability operator \verb'<' on the queue
        \item   a call to {\bf queuereadstatus()} on the queue
        \blend
        If the read clock domain is not yet established in any of these
        situations then it will be set to the current clock domain,
        otherwise it must be the same as the current clock domain.
        
        The read and write clock domains will usually be the same, in which case
        the queue variable is called {\bf synchronous}.
        
        When a queue variable is written (assigned), If it has no write clock
        domain the current clock domain will be established as the write clock
        domain of the variable. If it has a write clock domain it must be the
        same as the current clock domain.
        
        When a queue variable is written (assignedthe clock domain of
        the assigned value must be the same as the current clock domain or null.
        
        If the {\bf reset()} inbuilt procedure is used on a queue variable -
        \blist
        \item   if the queue variable has a write clock domain established, it
                must be the same as the current clock domain
        \item   if the queue variable does not have a write clock domain established,
                it will be given the current clock domain
        \blend
        
        When a queue variable read or examined the resulting value will have the
        read clock domain of the queue variable.
        
        Architectures supported -

        \begin{tabular}{| l | l |} \hline
        Spartan-II, Spartan IIE         & YES   \\ \hline
        Spartan-3                       & YES   \\ \hline
        Virtex, Virtex-E                & YES   \\ \hline
        Virtex-2, Virtex-2 Pro          & YES   \\ \hline
        Virtex-4, Virtex-5, Virtex-6    & YES   \\ \hline
        XC9500, XC9500XV, XC9500XL      & NO    \\ \hline
        Coolrunner XPLA3, Coolrunner-II & NO    \\ \hline
        \end{tabular}
        
        Applicable attributes are shown in the following table. The
        attributes can be passed as arguments to this procedure or may be
        applied to the variable using the inbuilt procedure
        {\bf attributes()} \autorefph{sec:ipattributes}.
        
        When {\bf queue()} is used to create a module, procedure or function
        parameter it cannot be passed attributes. The created parameter variable
        will adopt the attributes of the associated argument.

                \btabularx
        \hline
        \multicolumn{4}{|l|}{\bf queue mode attributes} \\
        \hline
        depth          & 3PL     & int     & queue buffer depth                          \\         
        mindepth       & 3PL     & int     & minimum queue buffer depth                  \\         
        buffersize     & 3PL     & int     & queue buffer depth                          \\         
        minbuffersize  & 3PL     & int     & minimum queue buffer depth                  \\         
        readdomain     & 3PL     & clock   & force the read (output) clock domain        \\    
        writedomain    & 3PL     & clock   & force the write (input) clock domain        \\    
        domain         & 3PL     & clock   & force both the read and write clock domains \\    
        alu            & 3PL     & boolean & use an ALU if appropriate                   \\       
        dsp            & 3PL     & boolean & use a DSP block                             \\       
        fifo           & 3PL     & boolean & use a hardware FIFO                         \\       
        queuereg       & 3PL     & boolean & extra output register to improve timing     \\       
        continuous     & 3PL     & boolean & no startup dependence                       \\   
        ignoremodules  & 3PL     & boolean & all reads treated as if in same module      \\   
        readonly       & 3PL     & boolean & set to read only                            \\
        queuefieldzero & netlist & boolean & clear unassigned fields                     \\
        \hline
        \btabularxend
        
        {\bf These attributes are case-insensitive.}
        
        {\bf depth} sets the depth of the queue buffer. The default depth is 2
        for a synchronous queue and 16 for an asynchronous queue.
        The values allowed depend on the FPGA family.

        \begin{tabular}{| l | l | l |} \hline
                                & using cmemory  & using rmemory \\ \hline \hline
        Spartan-II, Spartan IIE & 16             &  256 to  4096 \\ \hline
        Spartan-3               & 16             &  512 to 16384 \\ \hline
        Virtex                  & 16 to  32      &  256 to  4096 \\ \hline
        Virtex-E                & 16 to  32      &  256 to  4096 \\ \hline
        Virtex-II               & 16 to  64      &  512 to 16384 \\ \hline
        Virtex-II               & 16 to  64      &  512 to 16384 \\ \hline
        Virtex 4                & 16             &  512 to 16384 \\ \hline
        Virtex 5                & 16 to 128      & 1024 to 32768 \\ \hline                                        
        \end{tabular}
        
        A depth of 0 may also be set. This is a special case, creating an
        unbuffered queue. The use of unbuffered queues is very restricted -
        see section \autorefph{sec:unbufferedqueues}.

        {\bf buffersize} alternative older attribute name for {\bf depth}.

        {\bf mindepth} sets the depth of the queue buffer but rounds the
        supplied depth up to the next allowed value.

        {\bf minbuffersize} alternative older attribute name for {\bf mindepth}.

        {\bf readdomain} forces the read (output) clock domain to the specified
        clock domain. {\bf writedomain} forces the write (input) clock domain to
        the specified clock domain.
        {\bf domain} forces both the read (output) clock domain and the write
        (input) clock domain to the specified clock domain. 
        If {\bf null} is assigned, the current clock domain
        is selected. If the variable already has a clock domain as a result of
        having been assigned or evaluated then that clock domain is compared with
        the intended domain. If the assigned clock domain agrees with that already
        in place then there is no further action, otherwise a fatal error occurs.
        
        {\bf alu} specifies that assignments of expressions containing
        arithmetic operators '+' and '-' as the final or only operation can be
        combined into a single common adder/subtracter using CLBs. If the {\bf alu}
        attribute is false it can override the global directive {\bf ALU} to
        suppress the use of an ALU for an individual static variable.
        
        The {\bf alu} attribute can save logic and incurs no time penalty for the
        operations drawn in to the adder/subtracter, but all other assignments
        suffer an increase in propagation time since assigned values must now
        traverse the shared adder/subtracter. This tradeoff must be considered
        in each case.
        
        The {\bf alu} attribute cannot be added or changed once the variable has been
        assigned.
        
        The behaviour of the {\bf dsp} attribute depends on the FPGA -
        \blist
        \item   For Virtex-5 and Virtex-6 it has a similar effect to the
                {\bf alu} option in that it draws arithmetic and bitwise logic
                operations in assignments to the variable into a single ALU,
                but using a DSP block rather than CLBs.
        \item   For other FPGAs the attribute has no effect.
        \blend
        
        The {\bf dsp} attribute will not use a Virtex-5 or Virtex-6 DSP block as an
        ALU if the word width exceed 48 bits (multiplier operands are restricted
        to 25 by 18 bits).
        
        Note that for Virtex-5 or Virtex-6 if the {\bf alu} attribute is true but not
        the {\bf dsp} attribute then arithmetic operators '+' and '-' as the final or
        only operation will be combined into a single common adder/subtracter
        using CLBs.
        
        The {\bf fifo} attribute when true specifies that the queue buffer will
        use inbuilt block RAM FIFO logic if possible. It is ignored if the FPGA
        does not have this logic or the queue variable is asynchronous and uses
        the read or write status outputs. This attribute cannot be used if the
        queue is type "null", has the default depth of 2, or has a default depth
        which such that CLB RAM is used (depth \verb'<=' 16). In these cases an
        informative message is produced and the attribute is ignored. FIFO
        hardware can be enabled globally by setting directive {\bf global.FIFO}
        in which case this attribute can be assigned false to counteract the
        directive, i.e. the attribute when set overrides the directive. The FIFO
        directive is false by default.
        
        Use of the FIFO hardware is a compromise between speed and logic
        usage. The FIFO hardware saves CLB gates and flip-flops, which would
        otherwise be used, but has greater path delays which may cause
        failure to meet timing constraints.
        
        The {\bf queuereg} attribute when true causes an extra pipelined output
        register to be added. This reduces the clock-to-output time. The queue
        buffer logic incorporates the extra buffer with no change to
        functionality.

        {\bf Continuous} true specifies that if possible assignment will occur
        as soon as the FPGA is operational after configuration. This can
        increase operating speed by eliminating clock-enable logic. There are
        two situations where this can occur -
        \blist
        \item   the queue variable has a non-conditional assignment - the
                assignment operates continuously following configuration loading
                regardless of any 3PL startup mechanisms.
        \item   a {\bf when} statement has only unconditional assignments to
                variables in its nested code block(s) - the assignments are
                enabled directly from the {\bf when} test expression (or its
                inverse) and are not controlled from execution threads
                initiated by the 3PL startup mechanisms.
        \blend
        see {\bf start()}(\autorefph{sec:ipstart}).
        Note that if directive "continuous" is defined as type log with
        value true at the end of immediate execution, all queue variables
        will be treated as though they had the continuous attribute true.

        {\bf ignoremodules} true specifies that all queue reads will be
        treated as though from within the same module. Thus regardless of
        which module a queue read is contained within it will pop the queue
        when executed and the availability of the next datum will be seen
        within all modules in which the queue is read. This is distinct from
        the default behaviour of queue reads from multiple modules where a
        queue read from each relevant module must be executed before the next
        datum (if any) is seen as available to all the modules. This attribute
        is ignored for unbuffered queues.
        
        The purpose of this option is to allow code within multiple modules
        to read queues while looking after exclusive access where necessary
        by means of explicit user code. This can also be done also by ensuring
        that all the queue reads are contained within one module but this
        restriction can be inconvenient and may result in more obscure code.
        
        {\bf readonly} true specifies that from now on this variable cannot
        be assigned to. This attribute cannot be assigned false.
        
        {\bf queuefieldzero} false inhibits the clearing of unassigned fields or
        array members when a queue variable is assigned. The default is {\bf true}.
        If this attribute is made false for a queue variable which has a
        non-primitive type and an incomplete assignment is made, i.e. not all
        struct fields or array members are assigned, the unassigned elements
        for the datum will be undefined instead of explicitly zero. The value
        in making this attribute false is that it saves logic resources and
        reduces the delays on some data paths. Where a queue variable has a primitive
        type, or the whole variable is assigned in one statement or in multiple statements
        a par or sync block, this attribute has no effect.

        
        




    \section{ref - create one or more variables}\label{sec:ipref}
        \index{variable}
        
        {\bf Call:} \\
        ref(ident, mode, type, init);
        
        {\bf Output parameters:} \\
        none
        
        {\bf Input parameters:} \\
        \begin{tabular}{| l | l | l | l |}
        \hline
        param  & mode & type & \\
        \hline \hline
        ident or \verb'{'id1, id2...\verb'}' & - &  -       & variable identifier(s)          \\ \hline
        mode                   & - &  -       & optional mode keyword           \\ \hline
        type                   & I & type str & optional type or type string \\ \hline
        init                   & I &  -       & optional initial value          \\ \hline
        \end{tabular}
        
        {\bf Parameter=argument allowed:} \\
        No
        
        {\bf In-line target execution:} \\
        No
        
        {\bf Description:} \\
        Create a variable. This procedure behaves the same as {\bf var()}
        \autorefph{sec:ipvar} except where used to create an input parameter to
        a module, procedure or function. In that case its argument is passed by
        reference and not by value. This situation arises where an immediate or
        value mode variable is passed as an argument and a pointer to the
        variable via the parameter is required within the module, procedure or
        function. The argument passed to a {\bf ref()} parameter must be a
        variable and not an expression, but it may have subscripts or fields.
        
        See {\bf var()} \autorefph{sec:ipvar} for details.



    \section{remove - remove a list entry}\label{sec:ipremove}
        
        {\bf Call:} \\
        remove(l, index);
        
        {\bf Output parameters:} \\
        none
        
        {\bf Input parameters:} \\
        \begin{tabular}{| l | l | l | l |}
        \hline
        param  & mode & type & \\
        \hline \hline
        l      & I     & list & a list   \\ \hline
        index  & I     & int  & an index   \\ \hline
        \end{tabular}
        
        {\bf Parameter=argument allowed:} \\
        No
        
        {\bf In-line target execution:} \\
        No
        
        {\bf Description:} \\
        The item in list {\bf l} at position {\bf index} is removed.
        The index may range from 0 to {\bf entries(l) - 1}
        inclusive. If {\bf index} is outside this range a fatal error
        occurs.

    \section{removedirective - remove a directive}\label{sec:ipremovedirective}
        
        {\bf Call:} \\
        removedirective(key);
        
        {\bf Output parameters:} \\
        none
        
        {\bf Input parameters:} \\
        \begin{tabular}{| l | l | l | l |}
        \hline
        param  & mode     & type  & \\
        \hline \hline
        key   & I & str & directive key     \\ \hline
        \end{tabular}
        
        {\bf Parameter=argument allowed:} \\
        No
        
        {\bf In-line target execution:} \\
        No
        
        {\bf Description:} \\
        A directive is removed from the directive map. If the directive
        is not found a fatal error occurs.
        
        see also inbuilt procedures \\
        {\bf directive()} \autorefph{sec:ipdirective} \\
        {\bf presetdirective()} \autorefph{sec:ippresetdirective} \\
        and inbuilt functions \\
        {\bf directive()} \autorefph{sec:ifdirective}.\\
        {\bf direxists()} \autorefph{sec:ifdirexists}.\\
        {\bf directives()} \autorefph{sec:ifdirectives}.


    \section{reset - reset a static or queue variable}\label{sec:ipreset}
        
        {\bf Call:} \\
        reset(v1, v2,\dots);
        
        {\bf Output parameters:} \\
        none
        
        {\bf Input parameters:} \\
        \begin{tabular}{| l | l | l | l |}
        \hline
        param & mode & type & \\
        \hline \hline
        v & S Q & any & variable reference      \\ \hline
        - & - - &  -  & -                       \\ \hline
        - & - - &  -  & etc                     \\ \hline
        \end{tabular}
        
        {\bf Parameter=argument allowed:} \\
        No
        
        {\bf In-line target execution:} \\
        Yes
        
        {\bf Description:} \\
        This procedure resets one or more static or queue variables.
        
        For a static variable reset it to its initial value. The argument may
        have immediate or target subscripts or field identifiers, i.e. can
        be a variable or any components of a variable (any array members or
        struct fields). A reset overrides any simultaneous assignment. A reset
        uses different inputs to the static variable than assignments and can
        use less logic and be faster than an assignment.

        For a queue, clear the queue buffer making the queue unavailable. The
        argument must not have subscripts or field identifiers. It would be
        unusual to reset a single queue variable, the procedure normally being
        used to reset all queue variables in a pipeline to return to initial
        conditions. Because a queue variable can buffer more than one datum it
        is not possible to ascertain how many data are lost when a queue
        variable is reset, but resetting an entire pipeline results in
        a deterministic state.

        A static or queue variable can be reset at more than one place in the
        program, the resets being ORed. A reset overrides any simultaneous
        assignment.
        
        If {\bf reset()} is used on a static variable -
        \blist
        \item   if the static variable has a clock domain established, it
                must be the same as the current clock domain
        \item   if the static variable does not have a clock domain established,
                it will be given the current clock domain
        \blend
        
        If {\bf reset()} is used on a queue variable -
        \blist
        \item   if the queue variable has a write clock domain established, it
                must be the same as the current clock domain
        \item   if the static variable does not have a write clock domain
                established, the write clock domain will be set to the current
                clock domain
        \blend
        A queue reset will remove the write availability during
        the reset blocking any write. If the queue is synchronous the
        reset takes one clock cycle. If the queue is asynchronous the reset
        takes one write clock cycle plus a following read clock cycle plus
        another following write clock cycle.
        
        Note that if an immediate value, or a variable which is value mode
        and is constant, is assigned to a static variable or one of its
        components, the constant is checked against the initial value for
        that variable or component. If they are the same the variable or
        component is reset to its initial value rather than being assigned.
        Thus {\bf reset()} need not be used explicitly since assigning a value
        equal to the initial value will implicitly invoke a reset rather than
        an assignment.



    \section{rmemory - create a registered output memory}\label{sec:iprmemory}
        \index{variable!memory!registered output}
        
        {\bf Call:} \\
        rmemory(id, ports, atype, dtype, init, init0, init1 ...);
        
        {\bf Output parameters:} \\
        none
        
        {\bf Input parameters:} \\
        \begin{tabular}{| l | l | l | l |}
        \hline
        param  & mode       & type  & \\
        \hline \hline
        id     & - &  -           & memory identifier         \\ \hline
        ports  & I & int          & number of ports           \\ \hline
        atype  & I & uint         & address type              \\ \hline
        dtype  & I & any          & data type                 \\ \hline
        init   & I & []any        & optional initialisation (variable or compound
                                    constant value)                    \\ \hline
        init0  & I & any          & optional output latch A initialisation
                                    (variable or compound value)  \\ \hline
        init1  & I & any          & optional output latch B initialisation
                                    (variable or compound value)  \\ \hline
        ...    & I &  -           & optional attributes       \\ \hline
        \end{tabular}
        
        {\bf Parameter=argument allowed:} \\
        Yes
        
        {\bf In-line target execution:} \\
        No
        
        {\bf Description:} \\
        Procedure {\bf rmemory()} creates a registered output RAM.        
        {\bf rmemory()} may create a module, procedure or function
        parameter in which case only the first argument is needed (the memory
        identifier).

        In the case of Xilinx FPGAs, registered output memory is block RAM.
        This can have one or two access ports.

        Either port can be read or written (or both).
        
        {\em id} is an identifier for the memory. The variable identifier can
        also be constructed from a string by using the indirect operator on a
        string variable or parenthesised expression, see subsection
        \autorefph{subsec:strvarcreat}.

        {\em ports} is the number of ports. This argument may be by position or
        by ports=x where x is a constant or immediate variable. For Xilinx FPGAs
        theis must be 1 or 2.

        {\em atype} is the address type. This can be any target type but would
        usually be a "uint:". The address type has the constraint that the total
        number of bits must not exceed the maximum block RAM address width for
        the FPGA type. For Virtex, Virtex E and Spartan 2 this is 12 bits. For
        Virtex2, Virtex2 PRO and Spartan 3 this is 14 bits. This would almost
        always be an unsigned integer. This argument may be by position or by
        atype=x where x is a string specifying a type or else a type variable.

        {\em dtype} is the data type. This can be any target type. This argument
        may be by position or by dtype=x where x is a string specifying a type
        or else a type variable.
        
        {\em init} is an optional initialisation array. The initialisation data
        is a one-dimensional variable or array which is no larger than the depth
        of the memory. It may be smaller in which case the memory is zero
        padded. The types and structure of the initialisation variable or
        structured constant after the first dimension must match the memory data
        type. Immediate constants are padded to fit target elements but must not
        exceed the target width. If the initialisation argument is omitted the
        memory will be initialised to all zero. This argument may be by position
        or by init=x where x is an immediate array variable. For compound types
        the immediate type must match the target data type but omitting the data
        widths. The target type can be converted to a matching immediate type
        by the builtin function {\em targtoimtype()}.
        
        {\em init0} and {\em init1} are optional arguments to initialise the data
        output latches. The argument type must be the same as the memory data
        type. For complex types the argument may be a compound constant. These
        arguments may be by position or by init0=x or init1=x where x is an
        immediate value matching the memory data type.
        
        Where the {\em init}, {\em init0} or {\em init1} arguments are supplied
        in the call to rmemory() the arguments may be compound constants, e.g.
        \begin{lstlisting}[gobble=8, language=threepl]
            struct (st,
                f1, "uint:30",
                f2, "uint:30"
            );
            rmemory(m, 1, atype="uint:3", dtype=st, init0={44, 55});
        \end{lstlisting}
        though this is not usually practical for {\em init} which is usually
        a large array. Otherwise the arguments will be immediate variables
        of appropriate type. Initialising variables must be a matching type
        to the target data type of the memory. The initial value variable
        can be created from the target data type by calling builtin
        function {\bf targtoimtype()} \autorefph{sec:iftargtoimtype} -

        \begin{lstlisting}[gobble=8, language=threepl]
            struct (st,
                f1, "uint:30",
                f2, "uint:30"
            );
            var(init, "[8]" + targtoimtype(st));
            init[0].f1 = ....;
            init[0].f2 = ....;
            .
            .
            .
            var(init0, targtoimtype(st));
            init0.f1 = ..;
            init0.f2 = ..;
            rmemory(m, 1, atype="uint:3", dtype=st, init=init, init0=init0);
        \end{lstlisting}
            
        {\em init}, {\em init0} or {\em init1} can be assigned to a memory
        variable, along with attributes, later in execution by calling
        builtin procedure {\bf attributes()}\autorefph{sec:ipattributes}, e.g.
        \begin{lstlisting}[gobble=8, language=threepl]
            rmemory(m, 1, atype="uint:3", dtype=st);
            attributes(m, init=init, init0=init0);
        \end{lstlisting}
        but in this case the arguments must be immediate variables, not compound
        constants. The call to attributes() can be any time after calling
        rmemory() and (obviously) before termination of execution.
        The arguments to attributes() must be of the form attr=value.
        
        There may be trailing arguments following the optional {\em init}, {\em
        init0} and {\em init1} which supply memory attrbutes, described in a
        table below. {\em init}, {\em init0} or {\em init1} are included amongst
        the attributes for convenience. These must all be of the form
        attribute=value. These may be included in the call to rmemory() or else
        passed later using builtin procedure attributes().

                \btabularx
        \hline
        \multicolumn{4}{|l|}{{\bf rmemory mode attributes}} \\
        \hline
        init       & 3PL       &                      & initial value                    \\
        continuous & 3PL       & boolean              & continuous enable                \\   
        writemode  & netlist   & string               & write behaviour                  \\   
        writemodea & netlist   & string               & write behaviour for port A       \\   
        writemodeb & netlist   & string               & write behaviour for port B       \\   
        domain     & 3PL       & clock                & force clock domain for all ports \\     
        domain0    & 3PL       & clock                & force clock domain port0         \\     
        domain1    & 3PL       & clock                & force clock domain port1         \\     
        \hline
        \btabularxend
        
        {\bf These attributes are case-insensitive.}
        
        {\bf init} sets an initial value for the memory. Any previous initial
        value will be overwritten. The initial value must be a compound type and
        the nested type must be an immediate type equivalent to the memory data
        type, i.e. same type but without a width field. The initial value is not
        used until completion of immediate execution when variables are
        processed into target code and hence it can be changed at any time
        during immediate execution.
        
        {\bf continuous} true specifies that the block RAM enables are tied high
        instead of being gated by the read enable ORed with the write enable.
        This can gain speed and reduce routing resource usage but at the
        cost of increased noise and power consumption (and hence also FPGA
        temperature).
        
        {\bf write\_mode} allows the default write mode to be changed for
        block RAMs on Xilinx XC3S (Spartan 3), XC2V (Virtex 2) and
        XC2VP (Virtex 2 PRO) and XC4V (Virtex 4). The default is
        "write\_first". Other write modes are "read\_first" and "no\_change".
        {\bf write\_mode\_a} and {\bf write\_mode\_b} select the write mode on
        separate ports of dual-port block RAM.

        {\bf domain} forces the clock domain (both input and output) for all
        ports to the specified clock domain. If assigned {\bf null} the current
        clock domain is selected. If the variable already has a clock domain as
        a result of having been assigned or evaluated then that clock domain is
        compared with the intended domain. If the assigned clock domain agrees
        with that already in place then there is no further action, otherwise a
        fatal error occurs.

        {\bf domain0} and {\bf domain1} force the clock domain (both input and
        output) for port 0 or 1 of registered-output memory to the specified
        clock domain. If assigned {\bf null} the current clock domain is
        selected. If the variable already has a clock domain as a result of
        having been assigned or evaluated then that clock domain is compared
        with the intended domain. If the assigned clock domain agrees with that
        already in place then there is no further action, otherwise a fatal
        error occurs. These attributes cannot be set for combinatorial-output
        memory.

        
        Architectures supported -

        \begin{tabular}{| l | l |} \hline
        Spartan-II, Spartan-IIE         & YES   \\ \hline
        Spartan 3, Spartan 6            & YES   \\ \hline
        Virtex, Virtex-E                & YES   \\ \hline
        Virtex-II, Virtex-II Pro        & YES   \\ \hline
        Virtex 4, Virtex 5, Virtex 6    & YES   \\ \hline
        series 7, Zynq                  & YES   \\ \hline
        XC9500, XC9500XV, XC9500XL      & NO    \\ \hline
        Coolrunner XPLA3, Coolrunner-II & NO    \\ \hline
        \end{tabular}

        The range of memory depths that have been implemented is shown in the
        following table. The depth required is determined from the number of bits
        in the address type. Memories with depths smaller than the minimum listed
        are still implemented, but the physical memory will be the minimum size
        shown and the smaller depth will be achieved by tying higher address bits
        low.

        \begin{tabular}{| l | l | l |} \hline
                                 & 1 port       &   2 port    \\ \hline \hline
        Spartan-II, Spartan IIE,
        Virtex, Virtex-E         &  256 -  4096 &  256 -  4096 \\ \hline
        Spartan-3, Virtex-II,
        Virtex-II Pro,
        Virtex 4,                &  512 - 16384 &  512 - 16384 \\ \hline
        Spartan 6, Virtex 5,
        Virtex 6, series 7, Zynq & 1024 - 32768 & 1024 - 32768 \\ \hline
        \end{tabular}

        Registered output memory will employ the smallest address width which
        will encompass the address type. The data may be any
        width, multiple blocks being used if the data width
        exceeds the block data width associated with the address width. Both
        ports have the same address and data widths (the physical devices allow
        these to differ).
        
        Applicable attributes are shown in the following table. The
        attributes can be passed as arguments to this procedure or may be
        applied to the variable using the inbuilt procedure
        {\bf attributes()} \autorefph{sec:ipattributes}.
        
        When {\bf rmemory()} is used to create a module, procedure or function
        parameter it cannot be passed attributes. The created parameter variable
        will adopt the attributes of the associated argument.

    \section{rmemoryread - read a registered output memory}\label{sec:iprmemoryread}
        
        {\bf Call:} \\
        rmemoryread(id, port, addr);
        
        {\bf Output parameters:} \\
        none
        
        {\bf Input parameters:} \\
        \begin{tabular}{| l | l | l | l |}
        \hline
        param  & mode     & type  & \\
        \hline \hline
        id   & M             &  -  & memory being read  \\ \hline
        port & I             & int & port               \\ \hline
        addr & I V SV S Q PR & any & memory address     \\ \hline
        \end{tabular}
        
        {\bf Parameter=argument allowed:} \\
        Yes
        
        {\bf In-line target execution:} \\
        Yes
        
        {\bf Description:} \\
        Read from a port of a registered output memory.

        {\em id} is the memory mode variable identifier.
        
        {\em port} is the memory port.
        
        {\em addr} is the address.
        
        The address may be any type, including a compound type, but the type
        must match the type given when the registered output memory was created
        (see {\bf rmemory()} \autorefph{sec:iprmemory}) but primitive component
        widths may differ in which case padding or truncation will occur.

        If the address input argument contains queue reads execution of the write
        will wait on queue read availability.

        The value resulting from the read will be available from the following
        clock cycle at the output of the memory. This value can be accessed via
        the inbuilt function {\bf rmemoryout} \autorefph{sec:ifrmemoryout}.

        Multiple reads and writes on a port are allowed and the address and
        data inputs will be multiplexed, however there is no access arbitration
        and it is up to the application code to ensure exclusive access.
        Coincident accesses will result in data corruption.

        A call to {\bf rmemorywrite()} (\autorefph{sec:iprmemorywrite}) or {\bf
        rmemoryread()} (\autorefph{sec:iprmemoryread})  sets the clock domain
        for the memory port to the current clock domain. All further calls to
        either of these inbuilt procedures for the same port of the same
        registered output memory must be in the same clock domain (or a fatal
        error results).


    \section{rmemorywrite - write to a registered output memory}\label{sec:iprmemorywrite}
        
        {\bf Call:} \\
        rmemorywrite(id, port, addr, data);
        
        {\bf Output parameters:} \\
        none
        
        {\bf Input parameters:} \\
        \begin{tabular}{| l | l | l | l |}
        \hline
        param  & mode     & type  & \\
        \hline \hline
        id   & M             &  -  & memory being written   \\ \hline
        port & I             & int & port                   \\ \hline
        addr & I V SV S Q PR & any & memory address         \\ \hline
        data & I V SV S Q PR & any & data to be written     \\ \hline
        \end{tabular}
        
        {\bf Parameter=argument allowed:} \\
        Yes
        
        {\bf In-line target execution:} \\
        Yes
        
        {\bf Description:} \\
        Write to a port of a registered output memory. The registered output
        of the port is also changed to the written value.

        {\em id} is the memory mode variable identifier.

        {\em port} is the memory port.

        {\em addr} is the address.

        {\em data} is the data to be written.
        
        The data and address may be any type, including compound types, but
        the types must match the types given when the combinatorial
        output memory was created (see {\bf rmemory()} \autorefph{sec:iprmemory}) but
        primitive component widths may differ in which case padding or
        truncation will occur.
        
        If the address input or data input arguments contain queue reads
        execution of the write will wait on queue read availability.

        Multiple reads and writes on a port are allowed and the address and
        data inputs will be multiplexed, however there is no access arbitration
        and it is up to the application code to ensure exclusive access.
        Coincident accesses will result in data corruption.
        
        A call to {\bf rmemorywrite()} (\autorefph{sec:iprmemorywrite}) or {\bf rmemoryread()}
        (\autorefph{sec:iprmemoryread}) sets the clock domain for the memory port to the
        current clock domain. All further calls to either of these inbuilt
        procedures for the same port of the same registered output memory must
        be in the same clock domain (or a fatal error results).



    \section{rpt - print to the report file}\label{sec:iprpt}
        
        {\bf Call:} \\
        rpt(s);
        
        {\bf Output parameters:} \\
        none
        
        {\bf Input parameters:} \\
        \begin{tabular}{| l | l | l | l |}
        \hline
        param & mode & type & \\
        \hline \hline
        s & I & str & message string    \\ \hline
        \end{tabular}
        
        {\bf Parameter=argument allowed:} \\
        No
        
        {\bf In-line target execution:} \\
        No
        
        {\bf Description:} \\
        Print a message to the optional report file if it is active.
        
        {\bf s} is the message string.
                
        See also - \\
        {\bf println()} \autorefph{sec:ipprintln} \\
        {\bf texit()} \autorefph{sec:iptexit} \\
        {\bf msg()} \autorefph{sec:ipmsg}.


%    \section{select - a combinatorial selector with arbitration}\label{sec:ipselect}
%       
%        {\bf Call:} \\
%        select(o, so \verb'<-' s1, d1, s2, d2,...)
%        
%        {\bf output parameters:} \\
%        \begin{tabular}{| l | l | l | l |}
%        \hline
%        param  & mode & type & \\
%        \hline \hline
%        o  & V & any   & data output            \\ \hline
%        so & V & [log] & optional select output \\ \hline
%        \end{tabular}
%        
%        {\bf input parameters:} \\
%        \begin{tabular}{| l | l | l | l |}
%        \hline
%        param  & mode & type & \\
%        \hline \hline
%        s1 & V & log & 1st selector \\ \hline
%        d1 & V & any & 1st data     \\ \hline
%        s2 & V & log & 2nd selector \\ \hline
%        d2 & V & any & 2nd data     \\ \hline
%        .  & . &  .  & etc          \\ \hline
%        \end{tabular}
%        
%        {\bf In-line target execution:} \\
%        no
%        
%        {\bf Description:} \\
%        An inbuilt selector with arbitration using gate logic.
%        This has any number of input argument pairs.
%
%        The 1st input argument of each pair is a logical expression which selects
%        the following data input expression if true.
%
%        The 2nd input argument of each pair is a data input signal. If the
%        data input expressions are primitive they must be the same type but
%        their numeric widths may differ. The function value will be the width
%        of the widest data input. If the data inputs are not primitive then they
%        must all be exactly the same type, including primitive component widths.
%
%        The 1st output argument is the selector value. This must be mode VALUE
%        and the same type as the data inputs.
%
%        The optional 2nd output is an array of mode VALUE and type log one
%        member of which is true when the output is valid. The array dimension
%        must be equal to the number of input pairs and the member of this output
%        array that is true indicates which input was selected. The array is
%        always available.
%
%        The data output has a defined value when one or more select inputs is
%        available and true and the associated data input is available, the value
%        reflecting the selected data input. If more than one input is selected the
%        one with the lowest index prevails. If no inputs are selected, the value
%        is zero or false depending on data type(s).
%
%        If none of the inputs contain queues, the output value is
%        always available. If any inputs contain queues the output value
%        is available when an input selector and its associated data input are
%        available and the selector is true. Note that if some inputs contain
%        queues but one input selector and its associated data do not, the
%        output will be available when that selector is true. Note that the
%        use of this function as a gate with a clear value when no selectors are
%        true is only possible if the inputs contain no queues since any queues
%        in any inputs will render the output value unavailable when no
%        input is selected.


%    \section{selecta - a combinatorial selector from arrays}\label{sec:ipselecta}
%
%       
%        {\bf Call:} \\
%        selecta(o, so \verb'<-' s1, d1, s2, d2,...)
%        
%        {\bf output parameters:} \\
%        \begin{tabular}{| l | l | l | l |}
%        \hline
%        param  & mode & type & \\
%        \hline \hline
%        o  & V & any   & data output            \\ \hline
%        so & V & [log] & optional select output \\ \hline
%        \end{tabular}
%        
%        
%        {\bf input parameters:} \\
%        \begin{tabular}{| l | l | l | l |}
%        \hline
%        param  & mode & type & \\
%        \hline \hline
%        sel  & V & []log & selector array \\ \hline
%        data & V & []any & data array     \\ \hline
%        \end{tabular}
%        
%        {\bf In-line target execution:} \\
%        no
%        
%        {\bf Description:} \\
%        An inbuilt selector with arbitration using gate logic.
%        This has two input arguments and 1 or 2 output arguments.
%
%        The 1st input argument is a 1 dimensional array of type "log"
%        whose members select the corresponding data input. At most one
%        of these must be true at any time. The array may be an array of
%        pointers in which case the values pointed to will be used.
%
%        The 2nd input argument is a 1 dimensional array of inputs, one of which
%        is to be selected. The array may be an array of pointers in which case
%        the values pointed to will be used. If the data type is primitive, the
%        data inputs must all be the same primitive type but the data widths may
%        vary. The output type will be the same primitive type with the maximum
%        input width. If the data type is not primitive then all data inputs must
%        be of identical type and the output will be the same type.
%
%        The 1st output argument is the selector value. This must be mode VALUE
%        and the same type as the data inputs.
%
%        The optional 2nd output is an array of mode VALUE and type log one member of
%        which is true when the output is valid. The array dimension is equal to
%        the number of input pairs and the member of this output array that is true
%        indicates which input was selected. The array is always available.
%
%        The data output has a defined value when one or more select inputs is
%        available and true and the associated data input is available, the value
%        reflecting the selected data input. If more than one input is selected the
%        one with the lowest index prevails. If no inputs are selected, the value
%        is zero or false depending on data type(s).
%
%        If none of the inputs contain queues, the output value is
%        always available. If any inputs contain queues the output value
%        is available when an input selector and its associated data input are
%        available and the selector is true. Note that if some inputs contain
%        queues but one input selector and its associated data do not, the
%        output will be available when that selector is true. Note that the
%        use of this function as a gate with a clear value when no selectors are
%        true is only possible if the inputs contain no queues since any queues
%        in any inputs will render the output value unavailable when no
%        input is selected.
% 

    \section{selectvalue - create one or more selectvalue variables}\label{sec:ipselectvalue}
        \index{variable!selectvalue}
        
        {\bf Call:} \\
        selectvalue(ident, type, default, ...);
        
        {\bf Output parameters:} \\
        none
        
        {\bf Input parameters:} \\
        \begin{tabular}{| l | l | l | l |}
        \hline
        param  & mode & type & \\
        \hline \hline
        ident or \verb'{'id1, id2... \verb'}' & - & - & variable identifier(s) \\ \hline
        type    & I & str type & type or type string                           \\ \hline
        default & I &          & default value                                 \\ \hline
        ...     & I &  -       & optional attributes                           \\ \hline
        \end{tabular}
        
        {\bf Parameter=argument allowed:} \\
        No
        
        {\bf In-line target execution:} \\
        No
        
        {\bf Description:} \\
        Create one or more selectvalue variables. If more than one variable is to
        be created the list of variables must be in braces (\{\dots\}).
        
        {\bf type} is a type variable or a string expression giving the variable
        type. This may be omitted in the case of parameter creation in which case
        the variable will take the type of the argument.
        
        {\bf default} is optional and specifies the value the selectvalue mode
        variable should have when no input is selected. If omitted, 0 is used
        \footnote{If the "threestate" attribute is true the default value is
        implemented by the use of pullups or pulldowns. The specific FPGA
        platform may require a pullup explicitly on each signal, i.e. the
        default value may be constrained to be all 1s.}.
        
        Additional trailing arguments may be passed to set attributes. These
        must be of the form attr=value.
        
        A selectvalue variable has no clock domain when created. The first
        assignment to the variable establishes its clock domain as the current
        clock domain. A selectvalue variable which is evaluated prior to any
        assignment is given the clock domain current when the evaluation occurs.
        Later assignments to the variable must be in the same clock domain.
        
        Applicable attributes are shown in the following table. The
        attributes can be passed as arguments to this procedure or may be
        applied to the variable using the inbuilt procedure
        {\bf attributes()} \autorefph{sec:ipattributes}.
        
        When {\bf selectvalue()} is used to create a module, procedure or function
        parameter it cannot be passed attributes. The created parameter variable
        will adopt the attributes of the associated argument.

                \btabularx
        \hline
        \multicolumn{4}{|l|}{\bf selectvalue mode attributes} \\
        \hline
        domain     & 3PL & clock   & force the clock domain                    \\     
        direct     & 3PL & boolean & enable direct from test                   \\     
        readonly   & 3PL & boolean & set to read only                          \\     
        alu        & 3P  & boolean & use an ALU if appropriate                 \\        
        dsp        & 3PL & boolean & use a DSP block                           \\
        threestate & 3PL & boolean & use 3-state buffers rather than CLB logic \\
        \hline
        \btabularxend

        {\bf domain} forces the clock domain (both input and output) to
        the specified clock domain. If assigned {\bf null} the current
        clock domain is selected. If the variable already has a clock domain
        as a result of having been assigned or evaluated then that clock domain
        is compared with the intended domain. If the assigned clock domain
        agrees with that already in place then there is no further action,
        otherwise a fatal error occurs.

        {\bf direct} true specifies that if possible the output enable will be
        driven directly from a conditional signal, not through an execution
        thread and WHEN block. This can significantly simplify the logic and
        reduce the probability of failing to meet timing constraints locally.
        Note that when this attribute is in effect the selectvalue is operative
        following configuration loading regardless of any 3PL startup mechanisms
        - see {\bf start()}(\autorefph{sec:ipstart}).
        
        {\bf readonly} true specifies that from now on this variable cannot
        be assigned to. This attribute cannot be assigned false.
        
        {\bf alu} specifies that assignments of expressions containing
        arithmetic operators '+' and '-' as the final or only operation can be
        combined into a single common adder/subtracter using CLBs. If the {\bf alu}
        attribute is false it can override the global directive {\bf ALU} to
        suppress the use of an ALU for an individual static variable.
        
        The {\bf alu} attribute can save logic and incurs no time penalty for the
        operations drawn in to the adder/subtracter, but all other assignments
        suffer an increase in propagation time since assigned values must now
        traverse the shared adder/subtracter. This tradeoff must be considered
        in each case.
        
        The {\bf alu} attribute is only used at completion of immediate execution.
        
        The behaviour of the {\bf dsp} attribute depends on the FPGA -
        \blist
        \item   For Virtex-5 and Virtex-6 it has a similar effect to the
                {\bf alu} attribute in that it draws arithmetic and bitwise logic
                operations in assignments to the variable into a single ALU,
                but using a DSP block rather than CLBs.
        \item   For other FPGAs the attribute has no effect.
        \blend
        
        The {\bf dsp} attribute will not use a Virtex-5 or Virtex-6 DSP block as an
        ALU if the word width exceed 48 bits (multiplier operands are restricted
        to 25 by 18 bits).
        
        Note that for Virtex-5 or Virtex-6 if the {\bf alu} attribute is true but not
        the {\bf dsp} attribute then arithmetic operators '+' and '-' as the final or
        only operation will be combined into a single common adder/subtracter
        using CLBs.
        
        {\bf threestate} true specifies that three state buffers are to be used
        rather than CLB logic. If the device family does not support three state
        drivers then CLB logic will be used regardless of this attribute.
        Devices which have internal three-state buffers are Spartan-II, Spartan
        IIE, Virtex, Virtex-E, virtex-II, Virtex-II Pro, XC9500, XC9500XV and
        XC9500XL.




    \section{scanf - scan a line from stdin into variables}\label{sec:ipscanf}

        {\bf Call:} \\
        scanf(v1, v2,.... \verb'<-' fmt);
        
        {\bf Output parameters:} \\
        \begin{tabular}{| l | l | l | l |}
        \hline
        param  & mode     & type  & \\
        \hline \hline
        v1 & I & any & output variable  \\ \hline
        v2 & I & any & output variable  \\ \hline
        .  & I & any & output variable  \\ \hline
        .  & I & any & output variable  \\ \hline
        \end{tabular}
        
        {\bf Input parameters:} \\
        \begin{tabular}{| l | l | l | l |}
        \hline
        param  & mode     & type  & \\
        \hline \hline
        fmt  & I & str & format string      \\ \hline
        \end{tabular}
        
        {\bf Parameter=argument allowed:} \\
        No
        
        {\bf In-line target execution:} \\
        No
        
        {\bf Description:} \\
        This procedure scans a line from standard input extracting values from fields
        and assigning these to the output variables.
        
        The scanned line must consist of fields (tokens) separated by delimiters of commas
or spaces (any number of either, i.e. the delimiter pattern is [, ]+ ).

The format string is not the same as that used in scanf, sscanf or fscanf in C.
It is a string of single characters, one per conversion. Allowed
conversions are -

\begin{tabular}{| l | l |} 
\hline
'd' & integer                                   \\ \hline
'u' & unsigned integer                          \\ \hline
'x' & hexadecimal integer                       \\ \hline
'o' & octal integer                             \\ \hline
'f' & floating point value                      \\ \hline
's' & string                                    \\ \hline
'.' & convert according to the variable type    \\ \hline
'*' & skip the token                            \\ \hline
' ' & ignored                                   \\ \hline
tab & ignored                                   \\ \hline
\end{tabular}

The format string must not contain any other characters.
Spaces or tabs can be included anywhere if desired. Because the scanner cannot match explicit
fields from the format string and there are no conversion paramaters such as field widths or
precisions there is no need for a \verb"'%'" escape character. The number of actual
conversions (i.e. not including the skip) must be equal to the number of output variables.

The format argument may be null or "" in which case tokens will be converted and
assigned to the variables based on the variable types.

Tokens that are to be converted as hexadecimal values must not have a leading 0x.
Similarly a leading 0 will not be interpreted as indicating an octal value. The
conversion character alone decides which radix will be used.

          
        See the note on character and byte reads and writes in the description
        of inbuilt procedure {\bf file()} \autorefph{sec:ipfile}.
     
        See also - \\
        {\bf fscanf()} \autorefph{sec:ipfscanf} \\
        {\bf sscanf()} \autorefph{sec:ipsscanf} \\
        {\bf readln()} \autorefph{sec:ifreadln} \\
        {\bf readbyte()} \autorefph{sec:ifreadbyte}




    \section{sscanf - scan a string into variables}\label{sec:ipsscanf}

        {\bf Call:} \\
        sscanf(v1, v2,.... \verb'<-' line, fmt);
        
        {\bf Output parameters:} \\
        \begin{tabular}{| l | l | l | l |}
        \hline
        param  & mode     & type  & \\
        \hline \hline
        v1 & I & any & output variable  \\ \hline
        v2 & I & any & output variable  \\ \hline
        .  & I & any & output variable  \\ \hline
        .  & I & any & output variable  \\ \hline
        \end{tabular}
        
        {\bf Input parameters:} \\
        \begin{tabular}{| l | l | l | l |}
        \hline
        param  & mode     & type  & \\
        \hline \hline
        line & I & str & line to be scanned \\ \hline
        fmt  & I & str & format string      \\ \hline
        \end{tabular}
        
        {\bf Parameter=argument allowed:} \\
        No
        
        {\bf In-line target execution:} \\
        No
        
        {\bf Description:} \\
        This procedure scans an input string extracting values from fields
        and assigning these to the output variables.
        
        The scanned line must consist of fields (tokens) separated by delimiters of commas
or spaces (any number of either, i.e. the delimiter pattern is [, ]+ ).

The format string is not the same as that used in scanf, sscanf or fscanf in C.
It is a string of single characters, one per conversion. Allowed
conversions are -

\begin{tabular}{| l | l |} 
\hline
'd' & integer                                   \\ \hline
'u' & unsigned integer                          \\ \hline
'x' & hexadecimal integer                       \\ \hline
'o' & octal integer                             \\ \hline
'f' & floating point value                      \\ \hline
's' & string                                    \\ \hline
'.' & convert according to the variable type    \\ \hline
'*' & skip the token                            \\ \hline
' ' & ignored                                   \\ \hline
tab & ignored                                   \\ \hline
\end{tabular}

The format string must not contain any other characters.
Spaces or tabs can be included anywhere if desired. Because the scanner cannot match explicit
fields from the format string and there are no conversion paramaters such as field widths or
precisions there is no need for a \verb"'%'" escape character. The number of actual
conversions (i.e. not including the skip) must be equal to the number of output variables.

The format argument may be null or "" in which case tokens will be converted and
assigned to the variables based on the variable types.

Tokens that are to be converted as hexadecimal values must not have a leading 0x.
Similarly a leading 0 will not be interpreted as indicating an octal value. The
conversion character alone decides which radix will be used.

       
        See also - \\
        {\bf scanf()} \autorefph{sec:ipscanf} \\
        {\bf fscanf()} \autorefph{sec:ipfscanf} \\






    \section{setdirective - enter a directive - DEPRECATED!}\label{sec:ipsetdirective}
        
        {\bf Call:} \\
        setdirective(key, value);
        
        {\bf Output parameters:} \\
        none
        
        {\bf Input parameters:} \\
        \begin{tabular}{| l | l | l | l |}
        \hline
        param  & mode     & type  & \\
        \hline \hline
        key   & I & str & directive key                 \\ \hline
        value & I & any & directive value (optional)    \\ \hline
        \end{tabular}
        
        {\bf Parameter=argument allowed:} \\
        No
        
        {\bf In-line target execution:} \\
        No
        
        {\bf Description:} \\
        DEPRECATED - this procedure has been renamed {\bf directive()} but is
        otherwise unchanged. Procedure setdirective() has been retained
        temporarily for backward compatability.
        
        There is a single global directive map. The key is a string.  The value
        may one of the 3PL immediate primitive types {\bf int}, {\bf log}, {\bf
        float} or {\bf str}. If the second argument is not provided the entry
        will be type {\bf log} and value {\bf true}.
                
        see also inbuilt procedures \\
        {\bf directive()} \autorefph{sec:ipdirective}.\\
        {\bf presetdirective()} \autorefph{sec:ippresetdirective} \\
        {\bf removedirective} \autorefph{sec:ipremovedirective}. \\
        and inbuilt functions \\
        {\bf directive()} \autorefph{sec:ifdirective}.\\
        {\bf direxists()} \autorefph{sec:ifdirexists}.\\
        {\bf directives()} \autorefph{sec:ifdirectives}.


    \section{shell - execute a command in a shell}\label{sec:ipshell}
        
        {\bf Call:} \\
        shell(status, out, err \verb'<-' command, environ, dir, in);
        
        {\bf Output parameters:} \\
        \begin{tabular}{| l | l | l | l |}
        \hline
        param  & mode     & type  & \\
        \hline \hline
        status & I & int, uint & status returned from shell     \\ \hline
        out    & I & str       & standard output from shell     \\ \hline
        err    & I & str       & standard error from shell     \\ \hline
        \end{tabular}
        
        {\bf Input parameters:} \\
        \begin{tabular}{| l | l | l | l |}
        \hline
        param  & mode     & type  & \\
        \hline \hline
        command & I & str & command to be executed          \\ \hline
        environ & I & map & optional environment variables  \\ \hline
        dir     & I & str & optional working directory      \\ \hline
        in      & I & str & optional string to standard input to command \\ \hline
        \end{tabular}
        
        {\bf Parameter=argument allowed:} \\
        Yes
        
        {\bf In-line target execution:} \\
        No
        
        {\bf Description:} \\
        This procedure executes the command string passed as an argument using a
        shell. The default shell is {\bf /usr/bin/bash}, or if not found {\bf
        /bin/bash}, but this may be overridden by setting directive {\bf shell}
        to a string giving the path of the required shell.
        
        The first input parameter {\bf command} is a string containing the
        command path with any command line options and arguments.
        
        The optional second input parameter, {\bf environ}, is a map of environment
        variables. For each entry the key is the environment variable and the
        associated map value is the variable value. Value types are restricted
        to "int", "uint", "str" and "float", other types resulting in a fatal error.
        If no argument is supplied for this parameter the environment on the
        call to 3PL is passed to the shell.

	Note the shell may read user configuration files (such as {\bf
	.bash\_profile}) which may alter or replace {\bf environ} before {\bf command} is
	searched for and executed.
        To avoid this possibility consider using {\bf exec()} \autorefph{sec:ipexec}.
        
        The optional third input parameter, {\bf dir}, is a string specifying the working
        directory. If no argument is supplied for this parameter the current
        directory on the call to 3PL is used.
        
        The optional fourth input parameter, {\bf in}, is a string to be copied to
        standard input of the shell process. Note that if the shell or command
	closes its standard input stream, a fatal error occurs.
        
        If the first output parameter, {\bf status}, is given an argument, the
        status returned from the shell is copied to this variable. The argument
        must be immediate mode and type "int" or "uint".
        
        If the second output parameter, {\bf out}, is given an argument, the
        standard output from the shell is copied to this variable. The argument
        must be immediate mode and type "str". The resulting string is usually
        one or more lines, each terminated by a newline character.
        
        If the third output parameter, {\bf err}, is given an argument, the
        standard error from the shell is copied to this variable. The argument
        must be immediate mode and type "str". The resulting string is usually
        one or more lines, each terminated by a newline character.
        
        If the shell cannot be found, or cannot be executed, a fatal error occurs.
        
        This procedure waits until the shell exits. If the shell, or the command
        executed by the shell, is signalled then this procedure continues as for
        a normal exit.
        
        This procedure assumes that the operating system is a version of
        Unix, Linux or MacOS.

        See also {\bf exec()} \autorefph{sec:ipexec}.



%    \section{simread - read from a file during simulation}\label{sec:ipsimread}
%        
%        {\bf Call:} \\
%        simread(var \verb'<-' filename, step);
%        
%        {\bf Output parameters:} \\
%        \begin{tabular}{| l | l | l | l |}
%        \hline
%        param  & mode     & type  & \\
%        \hline \hline
%        var1 & V & any & value variable \\ \hline
%        var2 & V & any & value variable \\ \hline
%        .    & V & any & .              \\ \hline
%        \end{tabular}
%        
%        {\bf Input parameters:} \\
%        \begin{tabular}{| l | l | l | l |}
%        \hline
%        param  & mode     & type  & \\
%        \hline \hline
%        filename & I & str & external file path string  \\ \hline
%        step     & V & log & read signal                \\ \hline
%        \end{tabular}
%        
%        {\bf In-line target execution:} \\
%        No
%        
%        {\bf Description:} \\
%        This procedure is only used during simulation. When the 2nd input {\bf
%        step} argument is true, one line is read from file {\bf filename}, is
%        decoded according to the type of the output arguments and the output
%        values {\bf var1}, {\bf var2} etc. are assigned the decoded values.
%        The 2nd input argument {\bf step} must not contain queues.


%    \section{simwrite - write to a file during simulation}\label{sec:ipsimwrite}
%        
%        {\bf Call:} \\
%        simwrite(filename, var, step);
%        
%        {\bf Output parameters:} \\
%        none
%        
%        {\bf Input parameters:} \\
%        \begin{tabular}{| l | l | l | l |}
%        \hline
%        param  & mode     & type  & \\
%        \hline \hline
%        filename & I & str & external file path string  \\ \hline
%        step     & V & log & read signal                \\ \hline
%        var1     & V & any & value variable             \\ \hline
%        var2     & V & any & value variable             \\ \hline
%        .        & V & any & .                          \\ \hline
%        \end{tabular}
%        
%        {\bf In-line target execution:} \\
%        No
%        
%        {\bf Description:} \\
%        This procedure is only used during simulation. When the 2nd input
%        argument {\bf step} is true, the value of the 3rd and following
%        arguments {\bf var1},  {\bf var2} etc. are written as a line to file
%        {\bf filename}. None of the target arguments may contain queues.


    \section{special - vendor-specific special functions}\label{sec:ipspecial}
        
        {\bf Call:} \\
        special(\dots);
        
        {\bf Output parameters:} \\
        none
        
        {\bf Input parameters:} \\
        see description
        
        {\bf Parameter=argument allowed:} \\
        No
        
        {\bf In-line target execution:} \\
        No
        
        {\bf Description:} \\
        This procedure is used to implement special functions particular to the
        FPGA vendor. The first argument is an integer selecting the required
        special function. At the moment there are only two special functions
        as described below. Others could be added if required. Output arguments
        are not at the moment accepted as the existing special functions do
        not use any.

        \subsection{special 0 - add a line to the Xilinx NCF file}

            {\bf special(0, format\_string, port, target\_arguments)}
            
            The optional target arguments may be value, static, clock or
            selectvalue mode. Each target argument must be a single bit.
            
            The format string is written to the ISE NCF file. The format string
            may contain \%n or \%N conversions at which points netlist
            identifiers are to be inserted. The netlist identifiers are obtained
            from the input signals which are matched to the conversions in
            order. There will usually be only one of these though the format
            allows any number. \%n is replaced by the identifier of the next
            input signal argument. \%N is similar to the above but changes the
            identifier to upper case and removes any leading "GLOB\_" string.
            This is used for generating clock timing group names from the clock
            signal identifier. A \% followed by any other character is
            unchanged. The target mode signal arguments are matched to the \%n
            and \%N escape sequences left to right. The number of target mode
            arguments must must be suffient for the number of \%n and \%N escape
            sequences, any excess nets will being ignored.
            
            The port argument is optional and may be null (or missing if there
            are no target arguments to follow). If supplied it is a port name
            to which the constraint applies and no target arguments are required.
            If the port is ultimately unconnected to logic then the constraint
            will be omitted.
            
            This can also be called via library procedure
            {\bf addncfline(str(s))} \autorefph{sec:lpaddncfline} in
            {\bf common.3pl}, which only allows a
            single optional target argument, however this meets most requirements.


        \subsection{special 1 - add a line to the Xilinx XDC file}

            {\bf special(1, format\_string, port, target\_arguments)}
            
            The optional target arguments may be value, static, clock or
            selectvalue mode. Each target argument must be a single bit.
            
            The format string is copied to the XDC file. It must have a trailing
            ';' but does not need a newline character. The format string
            and input signal arguments are processed as above for the ISE
            NCF file constraints.
            
            The port argument is optional and may be null (or missing if there
            are no target arguments to follow). If supplied it is a port name
            to which the constraint applies and no target arguments are required.
            If the port is ultimately unconnected to logic then the constraint
            will be omitted.

            This can aslo be called via library procedure
            {\bf addxdcline(str(fmt), value(sig))} \autorefph{sec:lpaddxdcline} in
            {\bf common.3pl}, which only allows a
            single optional target argument, however this meets most requirements.
        







    \section{sprintf - copy a formatted line to a string variable}\label{sec:ipsprintf}

        {\bf Call:} \\
        sprintf(s \verb'<-' fmt, v1, v2,....);
        
        {\bf Output parameters:} \\
        \begin{tabular}{| l | l | l | l |}
        \hline
        param  & mode     & type  & \\
        \hline \hline
        s & I & str & output string   \\ \hline
        \end{tabular}
        
        {\bf Input parameters:} \\
        \begin{tabular}{| l | l | l | l |}
        \hline
        param  & mode     & type  & \\
        \hline \hline
        fmt  & I & str  & format string   \\ \hline
        v1   & I & any  & input variable  \\ \hline
        v2   & I & any  & input variable  \\ \hline
        .    & I & any  & input variable  \\ \hline
        .    & I & any  & input variable  \\ \hline
        \end{tabular}
        
        {\bf Parameter=argument allowed:} \\
        No
        
        {\bf In-line target execution:} \\
        No
        
        {\bf Description:} \\
        This procedure is equivalent to sprintf() in C. The format string is that used
        by the Java Formatter class.        
        
        The format argument may be followed by a single argument which is an array of
        pointers, for example {\bf inargs} in the calling module, procedure or function.
        In that case the pointers are resolved and provide the input variable list to be
        formatted.
        
        See also - \\
        {\bf printf()} \autorefph{sec:ipprintf} \\
        {\bf fprintf()} \autorefph{sec:ipfprintf} \\
        {\bf writeln()} \autorefph{sec:ipwriteln} \\
        {\bf writebyte()} \autorefph{sec:ipwritebyte} \\
        and library function {\bf sprintf()} \autorefph{sec:lfsprintf} \\





    \section{start - set the start condition for execution threads}\label{sec:ipstart}
        
        {\bf Call:} \\
        start(v, clk);
        
        {\bf Output parameters:} \\
        none
        
        {\bf Input parameters:} \\
        \begin{tabular}{| l | l | l | l |}
        \hline
        param & mode & type & \\
        \hline \hline
        v     & V,I    & log, uint  & start condition \\ \hline
        clk   & C      & optional clock \\ \hline
        \end{tabular}
        
        {\bf Parameter=argument allowed:} \\
        No
        
        {\bf In-line target execution:} \\
        No
        
        {\bf Description:} \\
        This procedure sets up the start of all execution threads for FPGAs.
        It cannot be used for XC2C or XC95 CPLDs.
      
        Every execution thread is started by a start pulse in its clock domain.
        
        These start pulses can be generated in four ways -
        \begin{enumerate}
        \item   If no call is made to start(), each clock domain will have an
                associated start pulse generated as soon as GWE is asserted
                after the configuration has been loaded.
        \item   if the argument to start() is value mode or input mode and
                is type "log" then a start pulse will be generated for every
                clock domain when that input is asserted.
        \item   if the argument to start() is immediate mode and type "uint" or
                "int" then the start will be delayed by the number of clock
                cycles specified by that integer 1st argument. The clock domain
                to which the count is to apply is specified by the optional
                2nd argument or else is the current clock. Following that delay
                start pulses will be generated in all clock domains.
        \item   if the argument to start() is immediate mode and type "float"
                then the start will be delayed by the time in second specified
                by the floating point 1st argument. The time will be in periods of
                the relevant clock domain, either the optional 2nd argument or else
                the current clock domain. Note that choosing the slowest clock will
                use the smallest counter.
        \end{enumerate}
        
        In all cases nothing occurs until the assertion of GWE.
        
        For case 2 the global execution start signal must be derived from
        external signals. The start level will usually either be a single
        external input signal or else an expression decoding an external bus or
        other communications interface. Any expression must not depend on any
        state devices which are clocked (e.g. static variables, queue variables,
        memory variables etc.) since these will not operate until the primary
        start pulse has occurred.     
        
        When start() is called for all except case 1 a clock domain is required.
        This will be the current clock domain by default or it can be specified in
        the optional 2nd argument. If the clock domain is null a fatal error will
        occur.
        
        In the case of Spartan 6, Virtex 4 and all later FPGAs there is a 
        START hardware element available which provides an output signal that
        is asserted high on completion of configuration startup. This can be used
        where the configuration startup waits for DCMs and PLLs to lock The start
        of execution of 3PL threads can be delayed until this signal is asserted,
        for example for Spartan 6 -
        
        \begin{lstlisting}[gobble=8, language=threepl]
           value(sl, "log");
           element("STARTUP_SPARTAN6");
           elementoutput(sl <- , "EOS");
           elementend();
           start(sl);
        \end{lstlisting}
        





    \section{statemachine - create a state machine - DEPRECATED!}\label{sec:ipstatemachine}
        
        {\bf Call:} \\
        statemachine(sf1, st1, e1, sf2, st2, e2, \dots, reset);
        
        {\bf Output parameters:} \\
        none
        
        {\bf Input parameters:} \\
        \begin{tabular}{| l | l | l | l |}
        \hline
        param  & mode     & type  & \\
        \hline \hline
        sf1   & S               & log & from state variable         \\ \hline
        st1   & S               & log & to state variable           \\ \hline
        e1    & I, V, SV, S, PR & log & state transition expression \\ \hline
        sf2   & S               & log & from state variable         \\ \hline
        st2   & S               & log & to state variable           \\ \hline
        e2    & I, V, SV, S, PR & log & state transition expression \\ \hline
        .     & .               &     & .                           \\ \hline
        reset & V, SV, S, PR    &     & optional reset              \\ \hline
        \end{tabular}
        
        {\bf Parameter=argument allowed:} \\
        No
        
        {\bf In-line target execution:} \\
        No
        
        {\bf Description:} \\
        No in-line code is created (the state machine is effectively a
        module) but the state variables are in the scope in which this
        procedure is called. Although not an error, it should not be called from
        within a target control structure such as {\em when}, {\em while} or
        {\em do while} as it does not generate in-line code and hence would
        not appear within the body of the control structure or be mediated by
        it in any way.

        States are represented by static logical variables, one and only one
        of which is true at any time (this is called 'one-hot encoding').

        The procedure has a multiple of 3 input arguments, with an optional
        extra trailing input argument, and no output arguments. The input
        arguments are triples, each of which describes a state transition.
        The 1st and 2nd arguments of a triple are the previous and next state
        of the transition and the 3rd argument of the triple is the logical
        expression causing the transition. The optional trailing input
        argument is a reset logical expression which, when it evaluates as
        true, will reset the state machine to its initial state. This reset
        overrides all other state transitions and resets all state variables to
        initial states (the first state true and all other states false).

        The states are unsubscripted variable identifiers. For each state
        appearing in a transition triple a target static variable of type log
        will be created in the current scope. If such a variable already
        exists it will be used instead (this allows variables to be created
        in advance for use in previous expressions prior to the call which
        creates the state machine). State variable identifiers must not have
        been previously used in any accessible scope except as described
        above. In addition if a state variable is pre-declared it must not be
        assigned a value or be initialised. State variables created by this
        procedure can never be assigned by user code but otherwise can be
        used as normal static logical variables. State variables for one
        state machine cannot be used as state variables in another state
        machine (the same identifiers can be used where the scope is not
        shared).

        The transition expressions may be immediate or target expressions.
        They may contain queue reads. Transition expressions must evaluate to
        type log.

        The optional trailing reset expression must evaluate to type log
        and must not contain queues.

        The first state encountered in the transition list is the initial
        state and hence will be initialised to true. This applies even when
        the initial state variable is pre-declared in which case this
        procedure will initialise that variable.

        At the moment there is little checking of the validity of the state
        transitions. The only check is that all states can be entered.
        Invalid state transitions may result in multiple or no state
        variables being true. Validation checking may be improved later.
        
        {\bf This procedure has been deprecated. Use the language construct
        {\em petrinet}.}


    \section{static - create one or more static variables}\label{sec:ipstatic}
        \index{variable!static}
        
        {\bf Call:} \\
        static(ident, type, ...);
        
        {\bf Output parameters:} \\
        none
        
        {\bf Input parameters:} \\
        \begin{tabular}{| l | l | l | l |}
        \hline
        param  & mode & type & \\
        \hline \hline
        ident or \verb'{'id1, id2...\verb'}' & - &  - & variable identifier(s) \\ \hline
        type                   & I & type str & optional type or type string   \\ \hline
        init                   & I &          & optional initial value         \\ \hline
        ...                    & I &  -       & optional attributes            \\ \hline
        \end{tabular}
        
        {\bf Parameter=argument allowed:} \\
        No
        
        {\bf In-line target execution:} \\
        No
        
        {\bf Description:} \\
        Create one or more value variables. If more than one variable is to
        be created the list of variables must be in braces (\{\dots\}).
        {\bf type} is a type variable or a string expression giving the
        variable type. This may be omitted in the case of parameter
        creation in which case the variable will take the type of the
        argument.
        
        The third argument is an optional initial value which must be immediate
        mode. If the type is compound the initialisation argument will usually
        be a compound value, but it may be initialised by a single immediate
        value in which case all members of the compound type will be initialised
        to that value. If more than one variable is created all will be
        initialised to the same value (the compound type members would have to
        be the same primitive type). An initial value may also be set later by
        calling inbuilt procedure {\bf attributes()}
        \autorefph{sec:ipattributes} with key "init". If no initial value has
        been given the initial value will be 0, false etc. as appropriate. If no
        initial value has been provided then attribute "init" will not appear in
        the attributes map returned by inbuilt function {\bf attributes()}
        \autorefph{sec:ifattributes}.
        
        Additional trailing arguments may be passed to set attributes. These
        must be of the form attr=value.
        
        A static variable has a single clock domain for reading and writing.
        When a static variable is created using the inbuilt procedure {\bf
        static()} it has no associated clock domain. The clock domain is
        established by
        \blist
        \item   the first assignment to the variable (or component of
                the variable if the type is compound)
        \item   the first use of the value
                of the variable or one of its components
        \blend
        
        {\bf Note that evaluation of the variable in any context, for example
        in inbuilt or library modules, procedures or functions, will establish
        a clock domain if the variable does not yet have a clock domain.}
        This will always be the current clock domain (unless the code has
        been explicitly written to establish the clock domain of its argument
        to something other than the current clock domain, in which case
        this will presumably be well documented).
        
        When a static variable is assigned, the clock domain of the assigned
        value must be the same as the current clock domain or null.
        
        When a static variable is assigned, If it has no clock domain the
        current clock domain will be established as the clock domain of the
        variable. If it has a clock domain it must be the same as the current
        clock domain.
        
        If the {\bf reset()} inbuilt procedure is used on a static variable -
        \blist
        \item   if the static variable has a clock domain established, it
                must be the same as the current clock domain
        \item   if the static variable does not have a clock domain established,
                it will be given the current clock domain
        \blend
        
        When the value of a static variable is used that value will have the
        clock domain of the variable.
        
        Applicable attributes are shown in the following table. The
        attributes can be passed as arguments to this procedure or may be
        applied to the variable using the inbuilt procedure
        {\bf attributes()} \autorefph{sec:ipattributes} and may be read back
        using inbuilt function {\bf attributes()} \autorefph{sec:ifattributes}.
        
        When {\bf static()} is used to create a module, procedure or function
        parameter it cannot be passed attributes. The created parameter variable
        will adopt the attributes of the associated argument.

                \btabularx
        \hline
        \multicolumn{4}{|l|}{\bf static mode attributes} \\
        \hline
        init       & 3PL  &           & initial value             \\
        domain     & 3PL  & clock     & force the clock domain    \\
        continuous & 3PL  & boolean   & no startup dependence     \\
        readonly   & 3PL  & boolean   & set to read only          \\
        alu        & 3PL  & boolean   & use an ALU if appropriate \\
        dsp        & 3PL  & boolean   & use a DSP block           \\
        loc        & UC:C & string(s) & locations                 \\
        iob        & UC:C & boolean   & locate in I/O blocks      \\
        ids        & UC:C & string(s) & element IDs               \\
        \hline
        \btabularxend
        
        3PL - 3PL uses the attribute itself for control logic or optimisation.
        
        UC:C - 3PL does not use the attribute - it must be used by 3PL source
        code to generate constraint file entries.
        
        {\bf These attributes are case-insensitive.}
        
        {\bf init} sets an initial value for the variable. Any previous initial
        value will be overwritten. If the static variable type is compound the
        initialisation value will usually be a compound value, but it may be
        initialised by a single immediate value in which case all members of the
        compound type will be initialised to that value (the compound type
        members would have to be the same primitive type). The initial value
        is not used until completion of immediate execution when variables
        are processed into target code and hence it can be changed at any time
        during immediate execution.

        {\bf domain} forces the clock domain (both input and output) to
        the specified clock domain. If assigned {\bf null} the current
        clock domain is selected. If the variable already has a clock domain
        as a result of having been assigned or evaluated then that clock domain
        is compared with the intended domain. If the assigned clock domain
        agrees with that already in place then there is no further action,
        otherwise a fatal error occurs.

        {\bf Continuous} true specifies that if possible assignment will occur
        as soon as the FPGA is operational after configuration. This can
        increase operating speed by eliminating clock-enable logic. There are
        two situations where this can occur -
        \blist
        \item   the static variable has a non-conditional assignment - the
                assignment operates continuously following configuration loading
                regardless of any 3PL startup mechanisms.
        \item   a {\bf when} statement has only unconditional assignments to
                variables in its nested code block(s) - the assignments are
                enabled directly from the {\bf when} test expression (or its
                inverse) and are not controlled from execution threads
                initiated by the 3PL startup mechanisms.
        \blend
        see {\bf start()}(\autorefph{sec:ipstart}).
        Note that if directive "continuous" is defined as type log with value
        true, all static variables which do
        not have this attribute set will be treated as though they had the
        continuous attribute true. The directive prevails if the attribute has
        not been set, otherwise the attribute is used.
        
        The "continuous" attribute should be set prior to assignments to the
        variable to have effect. Similarly the "continuous" directive should
        be set as early as possible in immediate execution or else prior
        to immediate execution using presetdirective().

        {\bf readonly} true specifies that from now on this variable cannot
        be assigned to. This attribute cannot be assigned false.
        
        {\bf alu} specifies that assignments of expressions containing
        arithmetic operators '+' and '-' as the final or only operation can be
        combined into a single common adder/subtracter using CLBs. If the {\bf alu}
        attribute is false it can override the global directive {\bf ALU} to
        suppress the use of an ALU for an individual static variable.
        
        The {\bf alu} attribute can save logic and incurs no time penalty for the
        operations drawn in to the adder/subtracter, but all other assignments
        suffer an increase in propagation time since assigned values must now
        traverse the shared adder/subtracter. This tradeoff must be considered
        in each case.
        
        The {\bf alu} attribute can only be applied to a static variable of
        primitive type; it will be ignored following a warning message
        otherwise.

        The {\bf alu} attribute is only used at completion of immediate execution.
        
        The behaviour of the {\bf dsp} attribute depends on the FPGA -
        \blist
        \item   For Virtex-5 and Virtex-6 it has two effects -
                \blist
                \item   it causes the static variable to be implemented using
                        the output register of a DSP block instead of
                        flip-flops.
                \item   For Virtex-5 and Virtex-6 it has a similar effect to
                        the {\bf alu} attribute in that it draws arithmetic and
                        bitwise logic operations in assignments to the
                        variable into a single ALU, but using a DSP block
                        rather than CLBs.
                \blend
        \item   For Virtex-2, Virtex-2Pro and Virtex-4\footnote{Virtex-4 has a
                DSP block but for reasons of local (non)usage the full
                capabilities of Virtex4 are not supported by 3PL. The DSP
                block is treated as a simple multiplier block.} it causes the
                static variable to be implemented using the output register
                of a multiplier block instead of flip-flops.
        \item   For other FPGAs the attribute has no effect.
        \blend
        
        The {\bf dsp} attribute will not implement the static variable using
        a multiplier or DSP register in the following cases -
        \blist
        \item   the variable is initialised to a non-zero value
        \item   the variable has the {\bf iob} attribute
        \blend
        The attribute will be ignored and a warning message issued in these
        cases.
        
        Note that the use of the {\bf dsp} attribute reduces the number of
        flip-flops required by using multiplier or DSP block output registers
        instead, but this may make timing constraints harder to meet because -
        \blist
        \item   the clock edge to output time of these may be worse that that
                for CLB flip-flops
        \item   the bits can no longer be distributed among CLBs and routing
                must instead converge on one or more multiplier or DSP blocks
        \blend
   
        In addition the {\bf loc} attribute if supplied will be ignored (with a
        warning message). For Virtex-2 and Virtex-2 PRO each multiplier element
        has a 36 bit register. For Virtex-5 and Virtex-6 the output register is
        48 bits. For variables whose type is not primitive each component will
        be implemented in a separate multiplier or DSP register. For components
        larger than the hardware register size sufficient DSP registers will be
        used to accomodate the component (but as noted above registers larger
        than 48 bits cannot use the Virtex-5 or Virtex-6 DSP block as an ALU).
        In the case of Virtex-5 or Virtex-6 , setting this attribute true also
        sets the {\bf alu} attribute true and the ALU logic will be implemented in the
        same DSP block. Note that if the limitations listed above are not met
        the {\bf dsp} attribute will be cleared (an associated warning message will be
        printed), the variable will be implemented using CLb flip-flops and if
        an ALU is generated it will use CLB logic.
        
        Note that for Virtex-5 or Virtex-6 if the {\bf alu} attribute is true but not
        the {\bf dsp} attribute then the static variable will be implemented using
        flip-flops and arithmetic operators '+' and '-' as the final or only
        operation can be combined into a single common adder/subtracter using
        CLBs.
        
        {\bf loc} provides absolute locations for the CLB flip-flops comprising
        the static variable (register). Locations are matched to flip-flops
        starting at the least significant bit. For a single bit static this
        attribute may be a string, otherwise it must be a string array. Each
        string must have the location format expected by the vendor for the FPGA
        type. 3PL does not at present provide support for relative placement
        of flip-flops in static variables though this could easily be added
        if required.

        {\bf iob} specifies that the flip-flops implementing the static variable
        are to be placed in IOB blocks rather than CLB blocks (this can also be
        requested by using the -pr option in the Xilinx map utility). If the
        value of the variable is used in addition to being connected to
        output buffers then the flip-flops will be duplicated in parallel, one
        set in the I/O blocks and the other in CLBs.
        
        {\bf ids} is a string array that contains the identifiers of the netlist
        elements (in this case D-type flip-flops) that implement the static mode
        variable (register). This attribute is read-only and is initialised when
        the variable is created. It can be used to apply constraints, such as
        absolute or relative location, to these elements.The array is ordered from
        least to most significant bit. For example to place a static variable in
        specific FPGA slice locations (for ISE) -
        
        \begin{lstlisting}[gobble=8, language=threepl]
        static(s, "uint:8");
        var(slocs, "[8]str", {"", "", "", .......}); // array of slice locations
        var(sids, "[8]str", attribute(s, ids));      // array of netlist element identifiers
        for (int(i) ; i<8 ; i++)
            addncfline("INST \""+ sids[i] + "\" LOC=" + slocs[i]);
        \end{lstlisting}
    
        For Vivado -
        
        \begin{lstlisting}[gobble=8, language=threepl]
        static(s, "uint:8");
        var(slocs, "[8]str", {"", "", "", .......}); // array of slice locations
        var(sids, "[8]str", attribute(s, ids));      // array of netlist element identifiers
        for (int(i) ; i<8 ; i++)
            addxdcline("set_property LOC " + slocs[i] + "[ get_cells " + sids[i] + "]");
        \end{lstlisting}
        
        
        



    \section{stop - stop a target execution stream}\label{sec:ipstop}
        
        {\bf Call:} \\
        stop();
        
        {\bf Output parameters:} \\
        none
        
        {\bf Input parameters:} \\
        none
        
        {\bf Parameter=argument allowed:} \\
        No
        
        {\bf In-line target execution:} \\
        yes
        
        {\bf Description:} \\
        Stop execution of the containing execution stream. Ultimately this
        will stop the module infinite loop in which it is contained.
        Parallel execution streams will continue to execute until their mutual
        completion point is reached. A stop() can be skipped or cancelled by
        an {\bf unless} construct on the containing block which will force the
        code to skip to the end of the block.


    \section{str - create one or more immediate string variables}\label{sec:ipstr}
        \index{variable!immediate!str}
        
        {\bf Call:} \\
        str(ident, init);
        
        {\bf Output parameters:} \\
        none
        
        {\bf Input parameters:} \\
        \begin{tabular}{| l | l | l | l |}
        \hline
        param  & mode & type & \\
        \hline \hline
        ident or \verb'{'id1, id2... \verb'}' & - &  -  & variable identifier(s)   \\ \hline
        init                   & I & log & optional initial value   \\ \hline
        \end{tabular}
        
        {\bf Parameter=argument allowed:} \\
        No
        
        {\bf In-line target execution:} \\
        No
        
        {\bf Description:} \\
        Create one or more immediate variables of type STR.
        
        {\bf ident} is the variable identifier. If more than one
        variable is to be created the list of variables must be in braces
        {\em id1, id2,\{\dots\}}. The variable identifier can also be constructed
        from a string by using the indirect operator on a string variable or
        parenthesised expression, see subsection \autorefph{subsec:strvarcreat}.
        
        A string expression may be provided as the second
        argument to initialise the variables, otherwise they will be
        initially "". If the type of the initialisation argument is
        INT, LOG or FLOAT it will be converted to a string.


    \section{struct - create a struct}\label{sec:ipstruct}
        \index{type!struct}
        
        {\bf Call:} \\
        struct(ident, f1, t1, f2, t2, \dots);
        
        {\bf Output parameters:} \\
        none
        
        {\bf Input parameters:} \\
        \begin{tabular}{| l | l | l | l |}
        \hline
        param  & mode & type & \\
        \hline \hline
        ident             & - &  -       & struct identifier         \\ \hline
        f1 or \verb'{'fa, fb...\verb'}' & - &  -       & field identifier(s)       \\ \hline
        t1                & I & str type & field type or type string \\ \hline
        f2 or \verb'{'fc, fd...\verb'}' & - &  -       & field identifier(s)       \\ \hline
        t2                & I & str type & field type or type string \\ \hline
         .                & . &  .       & etc.                      \\ \hline
        \end{tabular}
        
        {\bf Parameter=argument allowed:} \\
        No
        
        {\bf In-line target execution:} \\
        No
        
        {\bf Description:} \\
        Create a struct.
        
        {\bf ident} is the identifier of the struct.
        
        {\bf f1} is the identifier of the first field. More than one
        field of the same type can be specified by including a list
        of field identifiers in braces ({\em \{fa, fb,\dots\}}).

        {\bf t1} is the type or string type description of the first field.
        
        Optionally other argument pairs, follow.

        If a field identifier is the keyword {\bf extends} then the second
        component of the argument pair must be a struct identifier. The
        fields of that struct will be included in this struct at that point.

        The type of a field can of course be another struct in which case
        that struct is nested (rather than being extended).
        
        See \autorefph{sec:ssecstructs} for a description of structs and also
        an alternative string specification for creating them.



    \section{texit - print stack trace and terminate execution}\label{sec:iptexit}
        
        {\bf Call:} \\
        texit(message);
        
        {\bf Output parameters:} \\
        none
        
        {\bf Input parameters:} \\
        \begin{tabular}{| l | l | l | l |}
        \hline
        param  & mode     & type  & \\
        \hline \hline
        message & I & str & optional message   \\ \hline
        \end{tabular}
        
        {\bf Parameter=argument allowed:} \\
        No
        
        {\bf In-line target execution:} \\
        No
        
        {\bf Description:} \\
	If the argument is provided, print the message on standard out,
	prepended by the source code location.
        Print a stack trace to standard out, then terminate immediate execution.
        This lists all module, procedure and function calls and the source code
        file names and line numbers at which they occurred.
        
        See also - \\
        {\bf exit()} \autorefph{sec:ipexit} \\
        {\bf trace()} \autorefph{sec:iptrace}



    \section{trace - print a stack trace}\label{sec:iptrace}
        
        {\bf Call:} \\
        trace();
        
        {\bf Output parameters:} \\
        none
        
        {\bf Input parameters:} \\
        none
        
        {\bf Parameter=argument allowed:} \\
        No
        
        {\bf In-line target execution:} \\
        No
        
        {\bf Description:} \\
        Print a stack trace to standard out. This lists all module, procedure
        and function calls and the source code file names and line numbers at
        which they occurred.
        
        See also - \\
        {\bf exit()} \autorefph{sec:ipexit} \\
        {\bf texit()} \autorefph{sec:iptexit}


    \section{type - create a type variable}\label{sec:iptype}
        \index{type!type}
        
        {\bf Call:} \\
        type(v, descr);
        
        {\bf Output parameters:} \\
        none
        
        {\bf Input parameters:} \\
        \begin{tabular}{| l | l | l | l |}
        \hline
        param  & mode       & type  & \\
        \hline \hline
        v     & - &  -  & created immediate variable    \\ \hline
        descr & I & str & type description string       \\ \hline
        \end{tabular}
        
        {\bf Parameter=argument allowed:} \\
        No
        
        {\bf In-line target execution:} \\
        No
        
        {\bf Description:} \\
        This procedure creates a variable of type "type". The first
        argument is the identifier of the created variable. The second
        argument is the type specification string or another type variable
        or type returned from a function.
        
        A type variable can be used as the type argument to the
        various variable creation procedures instead of a string
        argument. It can also be passed as an argument to
        modules, procedures or functions where the matching
        parameter is also type "type".
        
        Note that a type variable can also be simply assigned a type
        description string which will result in the type being generated
        from the string.


    \section{useclock - set the current clock}\label{sec:ipuseclock}
        
        {\bf Call:} \\
        useclock(c);
        
        {\bf Output parameters:} \\
        none
        
        {\bf Input parameters:} \\
        \begin{tabular}{| l | l | l | l |}
        \hline
        param  & mode       & type  & \\
        \hline \hline
        c   & C  &  -    & clock mode variable \\ \hline
        \end{tabular}
        
        {\bf Parameter=argument allowed:} \\
        No
        
        {\bf In-line target execution:} \\
        No
        
        {\bf Description:} \\
        The procedure call sets the current clock domain, given by the clock
        mode variable argument, and pushes the previous current clock domain
        onto a stack. If there is no argument the previous clock domain is
        restored from the stack, or if that stack is now empty there is now no
        current clock domain.
        
        To explicitly set a null current clock domain the argument can
        be the keyword {\bf null}. The previous clock domain can later be
        restored by calling {\bf useclock()} with no argument. Another
        way to get a null current clock domain is to call {\bf useclock()}
        with no argument enough times to empty the clock stack.
        
        The clock cannot be changed within an execution stream. This procedure
        can only be called at the top level in a module (i.e. not nested within
        target control structures). Calls elsewhere will result in a fatal 3PL
        immediate execution error.


    \section{value - create one or more value variables}\label{sec:ipvalue}
        \index{variable!value}
        
        {\bf Call:} \\
        value(ident, type, init, ...);
        
        {\bf Output parameters:} \\
        none
        
        {\bf Input parameters:} \\
        \begin{tabular}{| l | l | l | l |}
        \hline
        param  & mode & type & \\
        \hline \hline
        ident or \verb'{'id1, id2... \verb'}' & - &  -       & variable identifier(s)       \\ \hline
        type                   & I & str type & optional type or type string \\ \hline
        init                   & I & str      & optional initial value       \\ \hline
        ...                    & I &  -       & optional attributes          \\ \hline
        \end{tabular}
        
        {\bf Parameter=argument allowed:} \\
        No
        
        {\bf In-line target execution:} \\
        No
        
        {\bf Description:} \\
        Create one or more value variables. If more than one variable is to be
        created the list of variables must be in braces (\{\dots\}).
        
        {\bf type}
        is a type variable or a string expression giving the variable type.
        This may be omitted in the case of parameter creation in which case
        the variable will take the type of the argument. It may also be
        omitted if an initialisation argument is given in which case it uses
        the type of that argument. In either case for an immediate constant
        the type is constructed from the size of the constant and its sign.

        A value mode variable may be created as type "empty", in which case it
        can be assigned any type. Following assignment the variable will
        retain the type of the RHS value assigned to it. Any further
        assignments to that variable will be subject to type checking as
        normal. A value variable can be cleared and its type changed to
        "empty" by calling inbuilt procedure {\bf emptyvalue}
        \autorefph{sec:ipemptyvalue}.

        The third argument is an optional initial value which may be a target
        expression or an immediate expression, i.e. a value may be assigned an
        initial constant -
        
        \begin{lstlisting}[gobble=12, language=threepl]
            static(a, "uint:10");
            value(v1, "uint:8", 5);         // constant initial value
            value(v2, "uint:12", a + 37);   // target expression initial value
        \end{lstlisting}

        If the value is an array it may be initialised by a single immediate
        value in which case all members of the array will be initialised to
        that value, otherwise it can be initilised by a compound constant or
        an immediate mode variable with a matching type.
        
        \begin{lstlisting}[gobble=12, language=threepl]
            var(n, "[3]int", {9, 6, 3});
            value(v1, "[3]uint:8", {5, 7, 9});  // constant initial array
            value(v2, "[3]uint:8", n);          // variable initial array
            value(v3, "[10]uint:12", 37);       // single value for array
        \end{lstlisting}
        
        For a compound type the initialisation does not have to initialise
        every member of the type, e.g.
        
        \begin{lstlisting}[gobble=12, language=threepl]
            var(n, "[3]int", {9, 6, 3});
            value(v1, "[6]uint:8", {5, 7, 9});  // constant smaller initial array
            value(v2, "[6]uint:8", {5,, 9});    // constant initial array with gap
            value(v2, "[6]uint:8", n);          // smaller variable initial array
        \end{lstlisting}
        
        Where an initial value is provided the type argument can occassionally
        be omitted, in which case the value variable will take the type of the
        initial value, e.g.
        
        \begin{lstlisting}[gobble=12, language=threepl]
            static({s1, s2}, "uint:8");
            value(v,, s1 + s2);
        \end{lstlisting}
        
        Here the initial expression is the sum of two unsigned 8-bit integers
        which will give a 9-bit result, hence {\bf v} will have type "uint:9".
        For this to be accepted the initial expression must have an
        unambiguous target type. If the initialisation argument is a compound
        value then the type argument must be given.

        In the call to {\bf value()} additional trailing arguments may be
        passed to set attributes. These must be of the form attr=value.

        A value variable has no initial default clock domain. A value variable
        which is evaluated prior to its assignment is given an initial value
        which has the clock domain that was current when the evaluation
        occurred. Later assignment to the variable must be in the same clock
        domain. A clock domain may be established by setting the domain
        attribute (see below).
        
        Applicable attributes are shown in the following table. The
        attributes can be passed as arguments to this procedure or may be
        applied to the variable using the inbuilt procedure
        {\bf attributes()} \autorefph{sec:ipattributes}.
        
        When {\bf value()} is used to create a module, procedure or function
        parameter it cannot be passed attributes. The created parameter variable
        will adopt the attributes of the associated argument.

                \btabularx
        \hline
        \multicolumn{4}{|l|}{\bf value mode attributes} \\
        \hline
        domain   & 3PL     & clock   & resolve the clock domains of assigned values \\     
        readonly & 3PL     & boolean & set to read only \\     
%        maxdelay & netlist & log     & maximum delay    \\     
        \hline
        \btabularxend
        
        {\bf These attributes are case-insensitive.}

        {\bf domain} resolves the clock domains of any values assigned to the value
        variable to the specified clock domain. If assigned {\bf null} the current
        clock domain is selected. If the values are already resolved their clock
        domains are compared with the intended domain. If the assigned clock domain
        agrees with that within the values then there is no further action,
        otherwise a fatal error occurs. If a value has not been assigned, the clock
        domain will be set and a future assignment must have the same clock domain.
        The attribute will affect all variables contained in the value expression.
        Setting the clock domain has a particular use for value variables which are
        outputs of an {\bf import()} \autorefph{sec:ipimport} procedure. Note that if
        this attribute is read back a null is returned.
        
        {\bf readonly} true specifies that from now on this variable cannot
        be assigned to. This attribute cannot be assigned false.
        
%        {\bf maxdelay} is a place-and-route constraint. This places an
%        upper limit on the propagation time through the signal net or nets
%        comprising this value.



    \section{var - create one or more variables}\label{sec:ipvar}
        \index{variable}
        
        {\bf Call:} \\
        var(ident, mode, type, init);
        
        {\bf Output parameters:} \\
        none
        
        {\bf Input parameters:} \\
        \begin{tabular}{| l | l | l | l |}
        \hline
        param  & mode & type & \\
        \hline \hline
        ident or \verb'{'id1, id2...\verb'}' & - &  -       & variable identifier(s)          \\ \hline
        mode                   & - &  -       & optional mode keyword           \\ \hline
        type                   & I & type str & optional type or type string \\ \hline
        init                   & I &  -       & optional initial value          \\ \hline
        \end{tabular}
        
        {\bf Parameter=argument allowed:} \\
        No
        
        {\bf In-line target execution:} \\
        No
        
        {\bf Description:} \\
        Create a variable.
        
        {\bf ident} is the variable identifier. If more than one
        variable is to be created the list of variables must be in braces
        {\em id1, id2,\{\dots\}}. The variable identifier can also be constructed
        from a string by using the indirect operator on a string variable or
        parenthesised expression, see subsection \autorefph{subsec:strvarcreat}.
        
        {\bf mode} is an optional mode keyword. This keyword must be
        one of
        \blist
        \item    immediate
        \item    value
        \item    selectvalue
        \item    static
        \item    queue
        \item    priority
        \blend
        The mode may be omitted (completely - no "," is needed) in which case the
        mode will be {\bf immediate} by default unless {\bf var} is being used to
        create a module, procedure or function parameter in which case it will use
        the mode of the matching argument (this can only be done for a single
        variable, not a variable list).
        
        {\bf var()} cannot create a variable whose mode is rmemory or cmemory
        but it can create a module, procedure or function parameter which
        is passed an rmemory or cmemory mode variable as argument. In that case
        the resulting parameter variable will reference the argument variable.
        
        See further decription at \autorefph{sec:callparams}.

        Variables of other modes cannot be created by this procedure. The inbuilt
        procedures for creating these are -

        \begin{tabular}{| l | l |}
        \hline
        mode    & procedure \\
        \hline
        \hline
        input   & input()   \autorefph{sec:ipinput}   \\ \hline
        output  & output()  \autorefph{sec:ipoutput}  \\ \hline
        cmemory & cmemory() \autorefph{sec:ipcmemory} \\ \hline
        rmemory & rmemory() \autorefph{sec:iprmemory} \\ \hline
        clock   & clock()   \autorefph{sec:ipclock}   \\ \hline
        \end{tabular}
        
        {\bf type} is a type variable or a string expression giving the
        variable type. This may be omitted in the case of parameter
        creation in which case the variable will take the type of the
        argument (this can only be done for a single variable, not
        a variable list). For priority mode the type must be an array of log.
        
        Where the first argument is a list of identifiers all of the
        created variables will have the same type.
        
        {\bf init} is an optional initialisation expression, or compound
        value, matching the type. For a compound type the initialisation may
        be a single immediate value in which case all members of the compound
        type will be initialised to that value. If omitted the variables, or
        their components, except for a value mode variable, will be initialised
        by default to 0, false "" etc. The initialisation argument cannot be
        used if the type argument is omitted, but the type argument may be
        "empty" in which case the type will be derived from the initialisation.
        Note that the initialisation argument may be IMMEDIATE, VALUE or CLOCK
        mode and the created variable will have the same mode if the mode
        argument is omitted (otherwise the modes must agree).
        
        Where the type is an array with missing dimensions the dimensions
        can be derived from an initialisation array, be it a variable or
        a compound value. For multiple variables all will be initialised
        to the same array values but note that the arrays are distinct -
        see \autorefph{sec:compoundarg}. If an argument is supplied the
        parameter will take its dimensions from the argument instead of the
        initialisation array.
        
        Where the variable created is a parameter to a module,
        procedure or function, missing type dimensions can be derived
        from the dimensions of the argument. If no argument is supplied
        the dimensions can still be derived from a default initialiser.
        If the call to {\bf var()} creates more than one parameter, i.e.
        the first argument is a bracketed identifier list, missing dimensions
        cannot be supplied from associated arguments, but they can still
        be derived from an initialiser. Remember that where an identifier
        list is supplied all the variables created will have the same type,
        though the dimensions for that type may have been derived from
        an initialiser. This means that the arguments matching those
        parameter variables must have the same dimensions. If parameters
        are to match arguments which are arrays of different dimensions then
        separate calls to {\bf var()} must be used for each parameter.
        
        Value mode variables can be initialised to an immediate or target
        mode expression.
        
        For convenience there are specific variable creation inbuilt procedures
        for some modes or types -
        
        \begin{tabular}{| l | l | l | l | l |}
        \hline
        mode  & type & procedure  & description & section \\
        \hline \hline
        immediate   & int    & int()         & immediate integer variable &
                                              \autorefph{sec:ipint}         \\ \hline
        immediate   & log    & log()         & immediate logical variable &
                                              \autorefph{sec:iplog}         \\ \hline
        immediate   & str    & str()         & immediate string variable &
                                              \autorefph{sec:ipstr}         \\ \hline
        immediate   & float  & float()       & immediate floating point variable &
                                              \autorefph{sec:ipfloat}       \\ \hline
        immediate   & type   & type()        & immediate type variable &
                                              \autorefph{sec:iptype}        \\ \hline
        immediate   & enum   & enum()        & immediate enumerated type variable &
                                              \autorefph{sec:ipenum}        \\ 
                    &        & enumv()       &                             &
                                              \autorefph{sec:ipenumv}       \\ \hline
        immediate   & struct & struct()      & immediate struct variable &
                                              \autorefph{sec:ipstruct}      \\ \hline
        immediate   & class  & class()       & immediate class variable &
                                              \autorefph{sec:ipclass}      \\ \hline
        value       & any    & value()       & value variable &
                                              \autorefph{sec:ipvalue}       \\ \hline
        selectvalue & any    & selectvalue() & selectvalue variable &
                                              \autorefph{sec:ipselectvalue} \\ \hline
        static      & any    & static()      & static variable &
                                              \autorefph{sec:ipstatic}      \\ \hline
        queue        & any    & queue()        & queue variable &
                                              \autorefph{sec:ipqueue}        \\ \hline
        priority    & any    & priority()    & priority variable &
                                              \autorefph{sec:ippriority}    \\ \hline
        cmemory     & any    & cmemory()     & combinatorial output memory variable &
                                              \autorefph{sec:ipcmemory}     \\ \hline
        rmemory     & any    & rmemory()     & registered output memory variable &
                                              \autorefph{sec:iprmemory}     \\ \hline
        clock       &  -     & clock()       & clock variable &
                                              \autorefph{sec:ipclock}        \\ \hline
        \end{tabular}
        
        Target mode variables have no associated clock domain when they are
        created. For a description of the establishment of clock domains
        see the specific target variable creation inbuilt procedures
        {\bf value()} \autorefph{sec:ipvalue},
        {\bf selectvalue()} \autorefph{sec:ipselectvalue},
        {\bf static()} \autorefph{sec:ipstatic},
        {\bf queue()} \autorefph{sec:ipqueue} and
        {\bf priority()} \autorefph{sec:ippriority} or the descriptions of
        selecvalue variable assignment \autorefph{sec:target_selvalass}
        value variable assignment \autorefph{sec:im_valass}
        static variable assignment \autorefph{sec:stat_ass}       
        queue variable assignment \autorefph{sec:queue_ass} or 
        priority variable assignment \autorefph{sec:priass}.
        
        See also {\bf var()} \autorefph{sec:ipref} where a module, procedure
        or function input parameter needs to ensure that its argument is passed by
        reference and not by value. This situation arises where an immediate or
        value mode variable is passed as an argument and a pointer to the variable
        via the parameter is required within the module, procedure or function.





    \section{writebyte - write a byte to a file}\label{sec:ipwritebyte}
        
        {\bf Call:} \\
        writebyte(file, e);
        
        {\bf Output parameters:} \\
        none
        
        {\bf Input parameters:} \\
        \begin{tabular}{| l | l | l | l |}
        \hline
        param  & mode & type  & \\
        \hline \hline
        file & I & file & file              \\ \hline
        e    & I & any  & value to print    \\ \hline
        \end{tabular}
        
        {\bf Parameter=argument allowed:} \\
        No
        
        {\bf In-line target execution:} \\
        No
        
        {\bf Description:} \\
        This procedure writes a byte to an output file. The file
        must have already been opened as an output file. The immediate
        value may be types {\bf bits}, {\bf uint}, {\bf int} or {\bf log}.
        A log argument will write 0 or 1. The other types must have a
        value between 0 and 255 or a fatal error will occur.
        
        See the note on character and byte reads and writes in the description
        of inbuilt procedure {\bf file()} \autorefph{sec:ipfile}.

        See also - \\
        writeln()\autorefph{sec:ipwriteln} \\
        printf()\autorefph{sec:ipprintf} \\
        fprintf()\autorefph{sec:ipfprintf} \\
        sprintf()\autorefph{sec:ipsprintf}.







    \section{writeln - write a line to a file}\label{sec:ipwriteln}
        
        {\bf Call:} \\
        writeln(file, e);
        
        {\bf Output parameters:} \\
        none
        
        {\bf Input parameters:} \\
        \begin{tabular}{| l | l | l | l |}
        \hline
        param  & mode & type  & \\
        \hline \hline
        file & I & file & file              \\ \hline
        e    & I & any  & value to print    \\ \hline
        \end{tabular}
        
        {\bf Parameter=argument allowed:} \\
        No
        
        {\bf In-line target execution:} \\
        No
        
        {\bf Description:} \\
        This procedure writes a line to an output file. The file
        must have already been opened as an output file. The immediate
        value may be types {\bf bits}, {\bf uint}, {\bf int}, {\bf log},
        {\bf str} or {\bf float}.
        
        See the note on character and byte reads and writes in the description
        of inbuilt procedure {\bf file()} \autorefph{sec:ipfile}.

        See also - \\
        printf()\autorefph{sec:ipprintf} \\
        fprintf()\autorefph{sec:ipfprintf} \\
        sprintf()\autorefph{sec:ipsprintf} \\
        writebyte()\autorefph{sec:ipwritebyte}.




